\section{Introduction}\label{s0}

\subsection{Short summary}
The study of nodal domains is a classical topic in spectral geometry. The founding work of R.\,Courant \cite{CH} from 1923 showed that an eigenfunction of the Laplacian corresponding to the $k$-th eigenvalue has at most $k$ nodal domains. In~the past century many contributions have led to a better understanding of nodal domains. An important step for this was a theorem of \AA.\,Pleijel \cite{Pl} from 1956, which shows that asymptotically the number of nodal sets of an eigenfunction becomes significantly smaller than the bound given by Courant, provided the dimension is strictly larger than~$1$. This initially involved the Dirichlet condition but was more recently extended to other boundary conditions by I.\,Polterovich \cite{Po} and C.\,Léna \cite{Le}; see also \cite{DFV}.

On the other hand, it~is natural to consider the same question for other operators and this leads naturally to the consideration of Dirichlet sub-Laplacians, initially called Hörmander operators \cite{Ho}, which share with the Laplacian the property of hypoellipticity. An important step in
their analysis was taken by L.\,P.\,Rothschild and E.\,M.\,Stein \cite{RS}, who proved their maximal hypoellipticity. They proceeded by comparison with operators on nilpotent Lie groups, a technique that is also known as the nilpotentization procedure. Using this approximation, G.\,Métivier \cite{Me} proved a beautiful Weyl formula for the asymptotic behavior of the counting function of eigenvalues, provided this approximation can be done ``uniformly''. Motivated by a recent paper \cite{EL} by S.\,Eswarathasan and C.\,Letrouit on Courant's nodal theorem and many discussions on related problems with C.\,Letrouit, our aim in this paper is to try to extend Pleijel's theorem to the sub-Riemannian context.

Our paper is divided into two, rather independent, parts. In~the first part, we~show how the validity of Pleijel's theorem in the sub-Riemannian case can be reduced to the specific analysis of sub-Laplacian on nilpotent groups. In~the second part, we~analyze the validity of Pleijel's theorem in open sets of specific groups related to the Heisenberg group. This leads us to the question of finding new explicit and close to optimal bounds on the constants for Sobolev inequalities, Faber--Krahn inequalities and isoperimetric inequalities.

\subsection{The Hörmander operator}
We consider in an open set $\Omega \subset \mathbb R^n$ the Dirichlet realization of a sub-Laplacian (also called Hörmander's operator)
\[
-\Delta_\bX^\Omega :=\sum_{j=1}^p X_j^\star X_j,
\]
where $X_1,\ldots,X_p$ are real $C^\infty$ vector fields on $\overline\Omega$ satisfying the so-called Hörmander condition \cite{Ho}, which reads:

\begin{assumption}\label{CHOr}
For some $r\geq 1$ the $X_j$, $j=1,\ldots,p$, and their brackets up to order $r$ generate at each point $ x\in \overline{\Omega}$ the tangent space $T_x \Omega$.
\end{assumption}

Here and in what follows we use the convention that the vector fields $X_1,\ldots,X_p$ are brackets of order $1$, that $[X_i,X_j]$ for $1\leq i,j\leq p$ are brackets of order $2$ and so~on.

The terminology ``sub-Riemannian Laplacian'' or in short ``sub-Laplacian'' is posterior to the work of L. Hörmander and corresponds to the development of sub-Riemannian geometry in the middle of the eighties; see for example \cite{AgBaBo,BeRi,Je,Str2}.

More generally, we~are given a connected $C^\infty$ Riemannian manifold $M$ of dimension~$n$ with a given measure $\mu$ (with a $C^\infty$-density with respect to the Lebesgue measure in a local system of coordinates) and a system $\bX= (X_1,\cdots,X_p)$ of real $C^\infty$ vector fields on $M$ satisfying Assumption~\ref{CHOr}. We~consider the operator
\[
-\Delta_\bX^{M,\mu} :=\sum_{j=1}^p X_j^\star X_j,
\]
where $X_j^\star$ is the formal adjoint obtained by using the $L^2$ scalar product with respect to the given measure $\mu$. In~local coordinates $X_j^\star =- X_j +c_j$ for a function $c_j$. In~the case when $M$ has a boundary, we~always impose a Dirichlet condition.

The sub-Laplacian $-\Delta_\bX^{M,\mu}$ is known to be hypoelliptic \cite{Ho} (meaning that if $\omega\subset M$ is open and $u$ is a distribution in $M$ with $-\Delta_\bX^{M,\mu} u$ smooth in~$\omega$, then $u$ is smooth in~$\omega$). By~Rothschild--Stein \cite{RS} it is maximally hypoelliptic, that is, it~satisfies
\[
\| X_k X_\ell \,u \|_{L^2(M,\mu)} \leq C \bigl(\|\Delta_\bX^{M,\mu} u \|_{L^2(M,\mu)} + \|u\|_{L^2(M,\mu)} \bigr), \quad \forall k,\ell, \forall u \in C_\rc^\infty(M).
\]
The latter result is proved through a technique of nilpotentization, which will also be important for us. Moreover, if~the boundary is $C^\infty$ and noncharacteristic for $\bX$ (\ie at each point of the boundary there exists a vector field $ X_j$ that is transverse to the boundary at the given point), then we have $C^\infty$-regularity up to the boundary. We~emphasize that we will \emph{not} need this latter condition for our results; see Section~\ref{srest}.

The operator $-\Delta_\bX^{M,\mu}$ has compact resolvent for instance when $M$ is closed, and we can ask all the questions about its discrete spectrum that have been solved throughout the years for the Dirichlet realization of the Euclidean Laplacian on a bounded open set. These include:
\begin{itemize}
\item Simplicity of the ground state or, more generally, its multiplicity.
\item Local structure of the nodal sets, density of the nodal sets,...
\item Courant's theorem: comparison between the minimal labeling $ k$ of an eigenvalue~$\lambda_k$ and the number $\nu_k$ of the nodal domains of the eigenfunction in the eigenspace corresponding to $\lambda_k$.
\item Pleijel's theorem.
\end{itemize}
We focus in this paper on the last two items. They will be described in more detail in the next two subsections.

To clarify our terminology: nodal domains of a real (eigen)function $u$ are the connected components of $\{ x\in M:\ u(x) \neq 0\}$; nodal sets are their boundaries.

\subsection{Courant's Theorem}
As is well-known, Courant's theorem in the case of the Dirichlet Laplacian on a bounded open subset of $\mathbb R^n$ states that an eigenfunction associated with the $k$-th eigenvalue has at most $k$ nodal domains:
\[
\nu_k \leq k.
\]
If one looks at the standard proof of Courant's theorem, this inequality mainly appears as a consequence of a restriction statement (the restriction of an eigenfunction to its nodal domain is the ground state of the Dirichlet realization of the Laplacian in this domain), the minimax characterization of the eigenvalues, and the Unique Continuation theorem (UCT). Hence the question is to determine under which conditions these three results extend to sub-Riemannian Laplacians.

Concerning the restriction statement, having rather limited information about the nodal sets, we~successfully adapt to the sub-Riemannian case a proof proposed in \cite{MP}, which permits to avoid regularity assumptions on $\partial \Omega$. The variational characterization then holds. The UCT was proved by K.\,Watanabe \cite{Wa} in the $C^\infty$ category in dimension 2, but H.\,Bahouri \cite{Ba, Ba2} gave a discouraging counterexample to UCT with two vector fields in $ \mathbb R^3$.
Here
we are fortunate to know that J.-M.\,Bony proved at the end of the sixties \cite{Bo} that UCT holds when the vector fields are analytic.\footnote{Note nevertheless that the sub-Laplacians with analytic vector fields are not in general hypoelliptic analytic. While the sub-Laplacian on the Heisenberg group $\mathbb H_n$ is known to be hypoelliptic analytic (its fundamental solution is explicitly known and analytic outside the origin
in the exponential coordinates; see \cite{Fo0}), the sub-Laplacian on $\mathbb H_n\times \mathbb R^k$ with $k\geq 1$ is known to be non-hypoelliptic analytic as a direct consequence of a result by Baouendi--Goulaouic \cite{BaGo}. We~refer to \cite{Me2} for a characterization of the nilpotent groups of rank $2$ whose associated sub-Laplacian is hypoelliptic analytic and to \cite{Hel} for other counterexamples when the rank of nilpotency is strictly larger than $2$. Here we have limited our references to the case of sub-Laplacians on nilpotent groups. Starting in the seventies there have been a lot of contributions on the subject, in particular around an ``evoluting'' Trèves conjecture.}
Hence Courant's theorem holds in the analytic category, as shown by S.\,Eswarathasan and C.\,Letrouit \cite{EL}. At the end of the next section, we~will extend statements given in \cite{EL} to the case when the boundary is not necessarily non-characteristic.

\subsection{Pleijel's Theorem}
In the same spirit, one can hope for an asymptotic control of $ \nu_k/k$ for large $k$ that improves over Courant's bound when the dimension is strictly larger than one. In~the case of the Dirichlet Laplacian on a bounded open subset of~$\mathbb R^n$ with $n\geq 2$, Pleijel's theorem \cite{Pl} says that there is an constant, independent of the open set and denoted by $\gamma(\mathbb R^n)$, such that
\begin{equation}
\label{eq:pleijelintro}
\limsup_{k\to +\infty} \frac{\nu_k}{k} \leq \gamma(\R^n)
\end{equation}
and, importantly,
\begin{equation}
\label{eq:pleijelintro2}
\gamma(\R^n)<1
\quad\text{for}\ n\geq 2.
\end{equation}
Because of this latter inequality, Pleijel's bound \eqref{eq:pleijelintro} provides an asymptotic improvement over Courant's theorem.

Later in this paper we will use the expression ``Pleijel's theorem holds'' (for a given operator) to mean the assertion that $\limsup_{k\to\infty} \nu_k/k<1$.

The proof of Pleijel's theorem is a nice combination of two ingredients. The first one is Weyl's formula, which describes the asymptotic behavior of the eigenvalue counting function $N(\lambda,-\Delta^\Omega)$ of the Dirichlet Laplacian $-\Delta^\Omega$:
\[
N(\lambda, -\Delta^\Omega) \sim \mathcal W(\mathbb R^n) \,|\Omega|\,\lambda^{\sfrac n2}
\quad\text{as}\ \lambda\to\infty.
\]
Here $\mathcal W(\R^n)$ is a certain explicit constant, depending only on $n$. The second ingredient is the Faber--Krahn inequality, which gives the following lower bound on the lowest eigenvalue $\lambda_1(-\Delta^\Omega)$ of the operator $-\Delta^\Omega$:
\[
\lambda_1(-\Delta^{\Omega}) \geq C^\FK(\R^n) \,|\Omega|^{-\sfrac 2n}.
\]
The constant $C^\FK(\R^n)$ is equal to the first Dirichlet eigenvalue on a ball of unit volume; for an expression of this constant in terms of Bessel functions see \eqref{eq:isofkexplicit} below. Pleijel's proof combines these two ingredients and leads to the inequality \eqref{eq:pleijelintro} with
\begin{equation}\label{PlConst}
\gamma (\R^n)= (C^\FK(\R^n))^{-\sfrac n2} \,\mathcal W(\mathbb R^n)^{-1}.
\end{equation}
Using the explicit expressions for the Weyl constant $\mathcal W(\mathbb R^n)$ and the Faber--Krahn constant $C^\FK(\R^n)$ one can establish \eqref{eq:pleijelintro2}; in this regard we refer to \cite[Part\,II, Lem.\,9]{BerMe}.

\subsection{A Pleijel bound in the sub-Riemannian case}
Our goal in the present paper is to generalize Pleijel's theorem to the case of sub-Laplacians, and so we are naturally led to the sub-Riemannian analogues of the two ingredients of its proof, namely Weyl's formula and the Faber--Krahn inequality.

We can be optimistic on the side of Weyl's formula. Since the pioneering work of G.\,Métivier \cite{Me} we are rich in results on the asymptotic distribution of eigenvalues, at least if we add to Assumption \ref{CHOr} a certain equiregularity condition (or Métivier's condition), which permits to approximate the vector fields $X_j$ at each point $x$ by the generators of a nilpotent Lie algebra $\mathcal G_x$. As~we already mentioned, we~work in the setting of a connected $C^\infty$ manifold $M$ (with or without boundary) of dimension $n$ with a given measure $\mu$ (with a $C^\infty$-density with respect to the Lebesgue measure in a local system of coordinates) and a system of $p$ $C^\infty$ vector fields $X_1,\ldots,X_p$ satisfying Assumption \ref{CHOr} with some $r$. In~addition, we~assume the vector fields satisfy the following:

\begin{assumption}\label{equireg}
For each $j\leq r$, the dimension of the space spanned by the brackets of $X_1,\ldots,X_p$ of length $\leq j$ at each point is constant.
\end{assumption}

In the language of sub-Riemannian geometry, the span of the vector fields satisfying Assumption \eqref{equireg} is called an \emph{equiregular distribution}. The simplest nontrivial example occurs with $r=2$, $p=2$ and $n=3$ with the Heisenberg group $\mathbb H$ and more generally in contact geometry; see for example \cite{AgBaBo,ArRi}.

We denote by $\mathcal{D}_j(x)$ the span at $x\in M$ of all vector fields obtained as brackets of length $\leq j$ of the $X_k$'s. We~recall that we use the convention $\mathcal D_1(x)=\allowbreak\Span\{X_1(x),\ldots,X_p(x)\}$. We~set $n_j :=\text{dim}(\mathcal{D}_j(x))$, which by Assumption \ref{equireg} above does not depend on the point $x\in M$. We~can then introduce, setting $n_0:=0$,\vspace*{-5pt}\enlargethispage{.5\baselineskip}
\begin{equation}\label{eq:defQ}
Q:= \sum_{j=1}^r j \,(n_j-n_{j-1}).
\end{equation}
This plays the role of an ``effective dimension''.

Under Assumptions \ref{CHOr} and \ref{equireg}, G.\,Métivier shows (using, in particular, the techniques of \cite{RS,Rot}) that there is a constant $c(M,\bX)$ such that\vspace*{-3pt}
\begin{equation}\label{MetForm1}
N(\lambda,-\Delta_\bX^{M,\mu}):= \# \{j:\ \lambda_j (-\Delta_\bX^{M,\mu}) \leq \lambda\} \sim c (M,\bX) \,\lambda^{\sfrac Q2}
\quad\text{as}\ \lambda\to\infty.
\end{equation}
We will come back later to the structure of $c (M,\bX)$ and to its computation in particular cases. Note that in the case $r=2$ related results are obtained in \cite{MS1,MS2,Me,Mo}, and Métivier's theorem (together with many other results) has been recently revisited in the light of sub-Riemannian geometry in \cite{CHT0,CHT,CHT1}.

This concludes our discussion of the first ingredient in Pleijel's proof, namely an analogue of Weyl's formula.

Concerning the second ingredient, namely an analogue of the Faber--Krahn inequality, our knowledge is rather poor. This is the question about minimizing the first Dirichlet eigenvalue among open sets of given measure. In~the case of the Heisenberg group, one can think of a result by P.\,Pansu \cite{Pa} concerning the isoperimetric inequality. C.\,Léna's approach \cite{Le} for treating the Neumann problem for the Laplacian could be helpful (see \cite{HS}) if the set in $M$ where the system $\bX$ is not elliptic is ``small'' in some sense, but our equiregularity assumption excludes this case.

In the first part of this paper we will follow another way by revisiting the nilpotentization procedure. This permits us to deduce Faber--Krahn inequalities for sub-Laplacians from Faber--Krahn inequalities for sub-Laplacians on nilpotent groups.

More precisely, under the above two assumptions we will prove that the Faber--Krahn inequality holds on subsets of small measure with a constant that is arbitrarily close to an integral over the constants of the Faber--Krahn inequalities on the nilpotent approximations $\mathcal G_x$; see Proposition \ref{faberkrahnbm}. This result is in the spirit of a result of Bérard--Meyer \cite[Part\,II, Lem.\,16]{BerMe}, who have shown that on a Riemannian manifold the Faber--Krahn inequality holds on subsets of small measure with a constant that is arbitrarily close to the constant in the Faber--Krahn inequality on $\R^n$. We~emphasize, however, that in the case of a Riemannian manifold the model space is the same at each point, namely $\R^n$, while in the sub-Riemannian setting the approximating model spaces $\mathcal G_x$ may vary with $x$. Our techniques are quite different from those employed in~\cite{BerMe} and its generalizations, e.g., in \cite{DFV}.

Combining our result about Faber--Krahn inequalities with Métivier's Weyl-type formula, we~obtain a sufficient condition for the validity of a Pleijel-type bound; see Theorem \ref{t:manifold}, which is the main result of the first part of this paper. The upper bound on $\limsup_{k\to\infty} \nu_k/k$ is of the form\enlargethispage{\baselineskip}
\begin{equation}
\label{eq:mainintro}
\biggl(\int_M (c_x^\FK)^{-\sfrac Q2} \,d\mu(x) \biggr) \biggl(\int_M c_x^\Weyl\,d\mu(x) \biggr)^{-1},
\end{equation}
where $c_x^\FK$ is a certain local Faber--Krahn constant, defined in terms of the nilpotentization of $-\Delta_\bX^{M,\mu}$ at $x\in M$, and $c_x^\Weyl$ is a certain local Weyl constant, defined in terms of the same nilpotentization.\footnote{This strengthening of our original result \cite{FrHe1,FrHe2} is due to Y.\,Colin de Verdière, who kindly allowed us to include his argument.} The precise definitions will be given below.

The role of the Borel measure $D\mto \int_D c_x^\Weyl\,d\mu(x)$ on $M$ is emphasized in \cite{CHT1}, where it is called the \emph{Weyl measure}. Similarly, here we introduce what may be called the \emph{Faber--Krahn measure} $D\mto \int_D (c_x^\FK)^{-\sfrac Q2} \,d\mu(x)$.

It is interesting to compare \eqref{eq:mainintro} with the Pleijel formula \eqref{PlConst}, to which it reduces in the case of open subsets of $\R^n$. More generally, in the Riemannian case (where $\mathcal D_1(x) = T_x M$ and where $\mu$ is the Riemannian volume measure) the expression \eqref{eq:mainintro} reduces to \eqref{PlConst} and we recover the result of Bérard and Meyer \cite{BerMe}. However, our result is already new in this case when $\mu$ is different from the Riemannian volume measure. In~the general sub-Riemannian case, the integration with respect to the measure $\mu$ takes into account that the model spaces $\mathcal G_x$ may vary with the point $x\in M$. In~this respect it is also interesting to note that \eqref{eq:mainintro} depends on $M$ and the vector fields $X_1,\ldots,X_p$, but does \emph{not} depend on the measure $\mu$. Indeed, both integrals in \eqref{eq:mainintro} do not depend on $\mu$; see Remark \ref{rem:indepmeas}.

According to \eqref{eq:mainintro}, a sufficient condition for the validity of Pleijel's theorem is the following bound on the ``local Pleijel constants'':
\[
\bigl(c_x^\FK \bigr)^{-\sfrac Q2} \,\bigl(c_x^\Weyl \bigr)^{-1} < 1
\quad\text{for all}\ x\in M ;
\]
see Corollary \ref{maincor}. We~emphasize that the latter condition involves the corresponding Faber--Krahn constants for Dirichlet realizations of sub-Laplacians in open set of nilpotent groups.

This provides a motivation for the second part of this paper, which is devoted to the validity of a Pleijel-type bound for the nilpotent groups $\mathbb H_n \times \mathbb R^k$, where $\mathbb H_n $ is the Heisenberg group of homogeneous dimension $2n+2$ and where $k\in\N_0$. While we have not been able to establish a Pleijel-type theorem in the most important case $(n,k)=(1,0)$, we~have succeeded in proving it if one admits the celebrated conjecture of Pansu concerning the isoperimetric constant on $\mathbb H_1$; see Proposition \ref{pansupleijel}. Even a nonsharp, but sufficiently good bound on the isoperimetric constant on $\mathbb H_1$ would imply a Pleijel-type theorem for $\mathbb H_1$.

We also have (unconditional) positive results under the assumption that the homogeneous dimension $Q=2n+2+k$ of $\mathbb H_n\times\R^k$ is sufficiently large. Indeed, the validity of Pleijel's theorem remains open for only four pairs $(n,k)$; see Theorem \ref{maincomp}, which is the main result of the second part of this paper.

\subsection{Organization}
This paper is divided into two main parts.

The first part begins in Section~\ref{srest} with a somewhat technical result concerning the restriction of a Sobolev function to a nodal domain, which plays an important role in the arguments of Courant and Pleijel and their generalizations. The main result of the first part is Theorem \ref{t:manifold} in Section~\ref{sec:mainsubriem}, which gives a sufficient condition for the validity of Pleijel's theorem via nilpotentization. In~the preceding Section~\ref{s1} we discuss the setting of this theorem and give the proof in the following Section~\ref{sec:mainsubriemproof}. The theorem is illustrated by an example in Section~\ref{AppA}.

The second part of this paper deals with the case of the Heisenberg group $\mathbb H_n$ and, more generally, with $\mathbb H_n\times\R^k$ with $k\in\N_0$. The main results of that part are summarized in Section~\ref{sec:mainhnrk}; see, in particular, Theorem \ref{maincomp} and Proposition \ref{pansupleijel}. The proofs of these results rely on an explicit form of the Weyl asymptotics, treated in Section~\ref{sec:weylhnrk}, and bounds on the Faber--Krahn constants. For the latter, we~proceed via two different techniques that are spread out over Sections \ref{sec:fksob} (continued in Section~\ref{s10}) and \ref{sec:isoperimetric}.

In two appendices we discuss an assumption appearing in the main result of the first part (Appendix \ref{AppA2}) and review an approach to Weyl asymptotics (Appendix \ref{AppC}).

\subsubsection*{Acknowledgements}
Thanks to D.\,Mangoubi for inviting the second author to propose open questions in an Oberwolfach workshop in August 2023 and for previous discussions. Great thanks also to C.\,Letrouit for stimulating discussions about his joint paper with S. Eswarathasan and sub-Riemannian geometry when we started to attack these questions. The comments of Y.\,Colin de Verdière, to whom we are most grateful, on a preprint version \cite{FrHe1} of this paper helped to significantly strengthen our results. Thanks to J.\,Viola and F.\,Nicoleau for helping us with the use of Mathematica or Wolframalpha. Thanks to V.\,Colin for his help in topology. We~are also grateful to L.\,Hillairet and P.\,Pansu for their help in the understanding of sub-Riemannian results and to L.\,Capogna and N.\,Garofalo for exchanges about the isoperimetric problem on the Heisenberg group. Thanks to E.\,Milman for discussions on Faber--Krahn constants. Thanks to I.\,Silvestre Roselló for pointing out some typographical errors.

\part{Courant's and Pleijel's bound in the sub-Riemannian case}

\section{The restriction of a Sobolev function to a nodal domain}\label{srest}

\subsection{Presentation}
In this section we will show that the restriction of a continuous Sobolev function to a nodal domain satisfies the boundary values zero in the sense of Sobolev functions. This holds under remarkably weak assumptions on the vector fields used to define the Sobolev spaces. The precise formulation can be found in Theorem~\ref{restriction} below.

That this restriction property requires a proof seems to have been overlooked by R.\,Courant in his original proof of the Courant nodal theorem, but this was proved in the case when $M$ is a compact manifold in \cite{BerMe}. In~the case with boundary, the proof goes in the same way when the Dirichlet boundary problem is regular (see \cite{An} or more recently \cite{LMP}). Without this assumption, it~was proved in a paper by Müller-Pfeiffer \cite{MP} that seems to be little known in the spectral theory community. An alternative proof, due to D.\,Bucur, is presented in the book \cite{LMP}. Here we show that the proof in \cite{MP} can be generalized to the sub-Riemannian setting. Notice that Bucur's proof could also have been adapted.

\subsection{Setting}\label{ss2.2}
Let $M$ be a $C^1$-manifold (without boundary) of dimension $n$ and let~$\mu$ be a $C^1$ nonnegative Borel measure on $M$. (By $C^1$-measure we mean that in each chart the measure is absolutely continuous with a positive $C^1$-density.) Let $X_1,\ldots,X_p$ be $C^1$-vector fields on $M$.

Notice that in this section our results are established under weaker assumptions than in the rest of this paper. In~particular, we~emphasize that we do not make an assumption on the span of the vector fields. In~fact, even the trivial case $p=0$, where there are no vector fields at all, is formally included in our analysis.

We define first-order differential operators $X_1^\star,\ldots,X_p^\star$ by
\[
\int_M \psi X_j \phi\,d\mu = \int_M \phi X_j^\star \psi\,d\mu
\quad\text{for all}\ \phi,\psi\in C^1_\rc(M).
\]
We define $S^1(M)$ to be the set of all functions $u\in L^2(M)$ for which there are $f_1,\ldots f_p\in L^2(M)$ such that for $j=1,\ldots,p$ one has
\[
\int_M u X_j^\star\phi\,d\mu = \int_M \phi f_j\,d\mu
\quad\text{for all}\ \phi\in C^1_\rc(M).
\]
The $f_j$'s are necessarily unique and we denote them by $f_j=:X_j u$, thus extending the usual notation in case $u\in C^1_\rc(M)$. We~set
\[
q[u] := \sum_{j=1}^p \|X_j u\|_{L^2}^2
\quad\text{for all}\ u\in S^1(M).
\]
One easily sees that $S^1(M)$ is a vector space that is complete with respect to the norm $\sqrt{q[u] + \|u\|_{L^2}^2}$. We~denote by $S^1_0(M)$ the closure of $C^1_\rc(M)$ in $S^1(M)$. When~$\Omega$ is an open subset of $M$, then $\Omega$ itself is a manifold and therefore the spaces $S^1(\Omega)$ and $S^1_0(\Omega)$ are defined. We~also introduce the space $S^1_\loc(M)$ of functions such that $u|_\Omega \in S^1(\Omega)$ for any open $\Omega\subset M$ with $\overline\Omega$ compact.

We emphasize that the spaces $S^1(M)$ and $S^1_0(M)$ depend on $X_1,\ldots,X_p$ and, if~$M$ is not compact, on $\mu$, even if this is not reflected in the notation. Note that if $p=0$, then $S^1(M)=S^1_0(M)=L^2(M)$.
If $M=\mathbb R^n$ (or an open set in $\mathbb R^n$), $p=n$, and $X_i=\partial_{x_i}$, then we recover the classical Sobolev spaces.

We begin by recording truncation properties, which will play an important role in our arguments.

\begin{lemma}\label{truncation}
If $u,v\in S^1_\loc(M)$, then $\max\{u,v\}, \min\{u,v\}\in S^1_\loc(M)$ and
\[
X_j \max\{u,v\}(x) = \begin{cases}
X_j u(x) & \text{a.e.\ in}\ \{ x\in M:\ u(x)\geq v(x)\},\\
X_jv(x) & \text{a.e.\ in}\ \{ x\in M:\ u(x)\leq v(x)\},
\end{cases}
\]
and
\[
X_j \min\{u,v\}(x) = \begin{cases}
X_j v(x) & \text{a.e.\ in}\ \{ x \in M:\ u(x)\geq v(x) \},\\
X_ju(x) & \text{a.e.\ in}\ \{ x \in M:\ u(x)\leq v(x) \}.
\end{cases}
\]
In particular, if~$u,v\in S^1(M)$, then $\max\{u,v\}, \min\{u,v\}\in S^1(M)$.
\end{lemma}

\begin{proof}
This lemma is well-known (see, e.g., \cite[Lem.\,3.5]{GaNh} when the underlying manifold is $\R^k$). For the sake of completeness we give an outline of the main steps of the proof.

First, one shows that if $\eta\in C^1(\R)$ with $\eta'$ bounded, then $u\in S^1_\loc(M)$ implies $\eta(u)\in S^1_\loc(M)$ with $X_j\eta(u) = \eta'(u)X_j u$. (Here one can argue as in \cite[Th.\,1.18]{HeKiMa}.) Next, one applies this result to $\eta(t)=\sqrt{t^2+\epsilon^2}$ and deduces, after passing to the limit $\epsilon\to 0$, that $u\in S^1_\loc(M)$ implies $|u|\in S^1_\loc(M)$ with
\[
X_j|u|(x) =
\begin{cases}
X_j u(x) & \text{a.e.\ in}\ \{ x\in M:\ u(x)\geq 0\},\\
- X_ju(x) & \text{a.e.\ in}\ \{ x\in M:\ u(x)\leq 0 \}.
\end{cases}
\]
Note that this shows, in particular, that $X_j u(x)=0$ a.e. in $\{x\in M:\ u(x)=0\}$. (Here one can argue similarly as in \cite[Lem.\,1.19]{HeKiMa}.) Since
\begin{align*}
\max\{u,v\} & = \tfrac12(u(x)+v(x) + |u(x)-v(x)|), \\
\min\{u,v\} & = \tfrac12(u(x)+v(x) - |u(x)-v(x)|),
\end{align*}
this implies the assertion of the lemma.
\end{proof}

\subsection{The restriction theorem}
The following theorem is the main result of this section.

\begin{theorem}\label{restriction}
Let $\Omega\subset M$ be open and let $u\in S^1_0(\Omega)\cap C(\Omega)$. If $\omega$ is a connected component of $\{ x\in\Omega:\ u(x)\neq 0\}$, then $u|_\omega \in S^1_0(\omega)$.
\end{theorem}

The proof below of this theorem is essentially taken from \cite{MP}. It~is simpler than Bucur's proof presented in \cite{LMP} and, in particular, avoids the notion of capacity. Similar to that proof, it~relies on truncation properties of Sobolev functions and on the following simple lemma.

\begin{lemma}\label{s10easy}
Let $\Omega\subset M$ be open and let $u\in S^1(\Omega)$. Assume that $u$ vanishes outside of a compact set of $M$ and vanishes on $\partial\Omega$ in the sense that for any $y\in\partial\Omega$ and any $\epsilon>0$ there is a neighborhood $U$ of $y$ in $M$ such that $|u|<\epsilon$ a.e.\ in $U\cap\Omega$. Then $u\in S^1_0(\Omega)$.
\end{lemma}

\begin{proof}[Proof of Lemma \ref{s10easy}]
This is modeled after \cite[Lem.\,1.26]{HeKiMa}; see also \cite[Lem.\,1]{HeNoSw}. We~first note that we may assume that $\overline\Omega$ is compact. Otherwise, we~consider an open set $\widetilde \Omega$ with compact closure such that $u$ vanishes almost everywhere in $\Omega\setminus\widetilde \Omega$ and apply the result on $\Omega\cap\widetilde \Omega$. Further, by~considering $u_+$ and $u_-$ separately (using Lemma \ref{truncation}), we~may assume that $u\geq 0$. Consider $u_\epsilon:=(u-\epsilon)_+$. We~shall show that $u_\epsilon\in S^1_0(\Omega)$. Since $u_\epsilon\to u$ in $S^1(\Omega)$, this will imply the assertion.

We claim that the support of $u_\epsilon$ is a compact subset of $\Omega$. Indeed, any $y\in\partial\Omega$ has an open neighborhood $U_y$ in $M$ with $u<\epsilon$ in $U_y\cap\Omega$. By~compactness of $\partial\Omega$, a finite union of these $U_y$ cover $\partial\Omega$ and $u_\epsilon$ vanishes almost everywhere in the intersection of this finite union with $\Omega$.

By the Meyers--Serrin type result in \cite[Th.\,1.13]{GaNh} we see that $u_\epsilon$ can be approximated in $S^1(\Omega)$ by functions from $C^1(\Omega)\cap S^1(\Omega)$. The result in that reference is stated for open subsets of $\R^n$, but since it is a local result, the result remains valid in our situation by localizing via a partition of unity to coordinate neighborhoods. The proof in \cite{GaNh} proceeds by convolution with a compactly supported function. The fact that the support of $u_\epsilon$ is a compact subset of $\Omega$ implies that the approximating functions belong to $C^1_\rc(\Omega)$. This proves that, indeed, $u_\epsilon\in S^1_0(\Omega)$.
\end{proof}

\begin{proof}[Proof of Theorem \ref{restriction}]
For the sake of concreteness let us assume that $u> 0$ in $\omega$, the argument in the opposite case being similar.

\subsubsection*{Step 1} Let
\[
z:= \begin{cases} u & \text{in}\ \omega,\\ 0 & \text{in}\ \Omega\setminus\omega. \end{cases}
\]
We claim that $z\in S^1_\loc(\Omega)$.

Let $\psi\in C^1_\rc(\Omega)$. It~is easy to see that $z\in C(\Omega)$ and consequently also $\psi z\in C(\Omega)$. Moreover, $\psi z=0$ on $(\Omega\setminus\omega)\cup(\Omega\setminus\supp\psi)$. Thus $\psi z|_\omega = \psi u|_\omega$ is in $S^1(\omega)\cap C(\omega)$ and extends continuously to $\partial\omega$, where it vanishes. By~Lemma \ref{s10easy} this implies that $\psi z|_\omega \in S^1_0(\omega)$. This implies that $\psi z\in S^1_0(\Omega)$. (Indeed, since $\psi z|_\omega\in S^1_0(\omega)$, it~can be approximated in $S^1(\omega)$ by functions in $C^1_\rc(\omega)$. Extending these functions by zero to $\Omega$ gives functions in $C^1_\rc(\Omega)$ and, since $\psi z$ vanishes in $\Omega\setminus\omega$, these functions approximate $\psi z$ in $S^1(\Omega)$, so $\psi z\in S^1_0(\Omega)$.) Since $\psi\in C^1_\rc(\Omega)$ is arbitrary, we~deduce $z\in S^1_\loc(\Omega)$, as claimed.

\subsubsection*{Step 2} Let $(\phi_j)\subset C^1_\rc(\Omega)$ such that $\phi_j\to u$ in $S^1(\Omega)$ (such functions exist since $u\in S^1_0(\Omega)$) and set
\[
v_j := \min\{ z, (\phi_j)_+ \}|_\omega.
\]
We claim that $(v_j)$ is a bounded sequence in $S^1_0(\omega)$ and that it converges to $u|_\omega$ in $L^2(\omega)$.

To prove this, we~consider $w_j:=\min\{ z, (\phi_j)_+ \}$. By~Step 1 and truncation properties of Sobolev spaces (Lemma \ref{truncation}) we deduce that $w_j\in S^1_\loc(\Omega)$. Moreover, these truncation properties also imply that
\begin{align*}
\sum_{i=1}^p \int_\omega |X_i v_j|^2\,d\mu & = \sum_{i=1}^p \int_\omega |X_i w_j|^2\,d\mu \\
& = \sum_{i=1}^p \int_{\{u \leq \phi_j \}} |X_i u|^2\,d\mu + \sum_{i=1}^p \int_{\{ u> \phi_j >0 \}} |X_i \phi_j|^2\,d\mu \\
& \leq \sum_{i=1}^p \int_\Omega (|X_i u|^2 + |X_i \phi_j|^2)\,d\mu.
\end{align*}
Since $u\in S^1(\Omega)$ and since $(\phi_j)$ converges in $S^1(\Omega)$, the right side is uniformly bounded.

Concerning the $L^2$-norm, we~find
\[
\int_\omega (v_j - u)^2\,d\mu = \int_{\omega\cap\{ (\phi_j)_+< u\}} (u - (\phi_j)_+)^2\,d\mu
\leq \int_\Omega (u-\phi_j)^2\,d\mu \to 0.
\]

We have shown, in particular, that $v_j\in S^1(\omega)$. It~remains to prove that $v_j\in S^1_0(\omega)$. We~note that $w_j$ is continuous in $\Omega$. Its restriction $v_j$ to $\omega$ extends continuously to~$\partial\omega$ and vanishes there. (The vanishing on $\partial\omega\cap\Omega$ comes from the continuity of $u$. The vanishing on $\partial\omega\cap\partial\Omega$ comes from the compact support property of $\phi_j$.) Thus, again by Lemma \ref{s10easy}, we~have $w_j|_\omega\in S^1_0(\omega)$.

\subsubsection*{Step 3} We can now finish the proof. Since $(v_j)$ is a bounded sequence in $S^1_0(\omega)$, after passing to a subsequence we may assume that it converges weakly in $S^1_0(\omega)$ to some $v\in S^1_0(\omega)$. In~particular, it~converges weakly in $L^2(\omega)$ to $v$. Meanwhile, by~Step~ 2, it~converges strongly in $L^2(\omega)$ to $u|_\omega$, so the uniqueness of the weak limit in $L^2(\omega)$ implies that $v=u|_\omega$ a.e. In~particular, $u|_\omega\in S^1_0(\omega)$, as claimed.
\end{proof}

\subsection{Application to Courant's theorem}\label{ss2.7}
Using Theorem \ref{restriction} we can recover the results in \cite{EL} concerning the Dirichlet realization $-\Delta^{\Omega,\mu}_\Xb$ in an open set $\Omega$ of $M$, but without having to assume that $\partial \Omega$ is noncharacteristic for the system of vector fields~$X_j$ (see Assumption 1.5 in \cite{EL} or Assumption ($H\Gamma$) in \cite{De}). We~record this as follows.

\begin{theorem}\label{courant}
For any $k\in \mathbb N$, any eigenfunction of $-\Delta_\Xb^{\Omega,\mu}$ with eigenvalue $\lambda_k$ has at most
$k+ \mult(\lambda_k)-1$ nodal domains, where $\mult(\lambda_k)$ denotes the multiplicity of~$\lambda_k$.
If, moreover, one of the two following assumptions holds
\begin{itemize}
\item $n=2$,
\item $M$, $\mu$ and $(X_1,\dots, X_p)$ are real-analytic,
\end{itemize}
then we get an upper bound by $k$.
\end{theorem}

The first part of the statement relies on a remark of D.\,Mangoubi that permits to avoid the use of the Unique Continuation Theorem (see \cite[Exer.\,4.1.15]{LMP}). For the second part of the theorem, \cite{EL} refers in the first case to \cite{Wa}. For the second case corresponding to the standard statement of Courant's nodal theorem, the proof is based on a result of J.-M.\,Bony \cite{Bo}. For the proof of both parts of the theorem we apply our Theorem \ref{restriction}. The required continuity of eigenfunctions in $\Omega$ follows from~\cite{Ho}.

Let us further discuss the assumption of being noncharacteristic in \cite{EL}. A characteristic point of $\partial \Omega$ relative to $\mathcal D_1$ is a point $x$ for which all the elements of $\mathcal D_1(x)$ belong to the tangent space $T_x\partial \Omega$. The assumption that $\partial\Omega$ is noncharacteristic for~$\mathcal D_1$ guarantees that eigenfunctions belong to $C^\infty(\overline{\Omega})$; see \cite{De}, \cite{EL} and the book \cite{DGN} (particularly its Chapter 3). This regularity allowed \cite{EL} to prove Theorem \ref{courant} (under their noncharacteristic assumptions). This extra assumption seems to be crucial for the regularity of eigenfunctions at the boundary.\footnote{Notice that the question of the analyticity at the boundary (in case the sub-Laplacians are hypoelliptic analytic in $\Omega$) seems open.} Remarkably, however, it~is not necessary for the validity of Theorem \ref{restriction}.

The condition that $\partial\Omega$ is noncharacteristic for $\mathcal D_1$ may appear rather strict. Examples where this condition holds in the case of the Heisenberg group are given in \cite{DGN,MM}. An example in $\mathbb H$ is given (see \cite[Example 3.4]{MM} with $k=1$) by the domain
\[
(\sqrt{x^2+y^2} -2)^2 + 16 \,t^2 < 1.
\]
Note that for this example $\partial \Omega$ is homeomorphic to $\mathbb T^2$.

Meanwhile, topological considerations show that for $\Omega\subset\mathbb H$ the condition that $\partial\Omega$ is noncharacteristic for $\mathcal D_1$ is never satisfied if $\partial \Omega$ is homeomorphic to $\mathbb S^2$.\footnote{Many thanks to V.\,Colin for this remark.
If $\Omega \subset \mathbb R^3$ and if we have an equiregular distribution ($n_1=p=2$) and $r=2$, then
the transversality condition is never satisfied if $\partial \Omega$ is homeomorphic to $\mathbb S^2$.
The reason is that the noncharacteristic condition implies the existence of a continuous (with respect to $x$) unique straight line in $T_x \partial \Omega$ (this is the intersection of $\mathcal D_1(x)$ with $T_x \partial \Omega$) and this is impossible when $\partial \Omega$ is homeomorphic to the sphere by the Hairy Ball theorem (or by the Poincaré--Hopf theorem \cite{Miln} using the fact that the Euler characteristic is not $0$).}

Finally it was shown in \cite{De} that the measure in the boundary of the characteristic points is zero. Other related results are obtained in the case of the Heisenberg group by D.\,Jerison \cite{Jer}.

\section{The nilpotent approximation}\label{s1}

\subsection{Nilpotentization of vector fields and measures}
Throughout this section we follow the presentation of Rothschild \cite{Rot}, which is based on assumptions and definitions given earlier by Goodman \cite{Go}, Folland--Stein \cite{FS}, Folland \cite{Fo}, Métivier~\cite{Me}, Rothschild--Stein \cite{RS}. Since this period in the seventies, a huge literature has been devoted to the topic of sub-Riemannian geometry, for which we refer to the appendix in \cite{CHT} and references therein. We~attempt to combine the two formalisms in this section as well as in Appendices \ref{AppA2} and \ref{AppC}.

We consider the situation presented in the introduction and suppose that Assumptions \ref{CHOr} and \ref{equireg} are satisfied. We~recall that $\mathcal D_1(x)$ was defined after Assumption~\ref{equireg}, and also that $n$ denotes the dimension of $M$ and $n_1$ the (constant) dimension of $\mathcal D_1(x)$.

Clearly, we~have $p\geq n_1$, but this inequality may be strict (meaning that the vectors $X_1(x),\ldots,X_p(x)$ are not linearly independent at some and then, by~equiregularity, any $x\in M$). In~order to deal with this situation, we~will apply Lemma \ref{locindep}, which says that for every point in $M$ there is an open neighborhood $W\subset M$ of the given point and vector fields $\tilde X_1,\ldots \tilde X_{n_1}$ defined in $W$ such that
\begin{equation}
\label{eq:loclinindep}
\Span\{ \tilde X_1(x),\ldots \tilde X_{n_1}(x) \} = \mathcal D_1(x)
\quad\text{for all}\ x\in W
\end{equation}
and
\begin{equation}
\label{eq:loclinindeplapl}
-\Delta_\bX^{M,\mu} f = \sum_{j=1}^{n_1} \tilde X_j^\star \tilde X_j f
\quad\text{for all} \ f\in C^2(W).
\end{equation}
The vector fields $\tilde X_1,\ldots,\tilde X_{n_1}$ again satisfy Assumptions \ref{CHOr} and \ref{equireg} (with the same $r$ and $n_j$).

We will first discuss the nilpotent approximation on an open set $W\subset M$ where \eqref{eq:loclinindep} and \eqref{eq:loclinindeplapl} are satisfied. Later we will argue that this gives a nilpotent approximation on all of $M$. Of course, in the special case $p=n_1$ we can immediately take $W=M$, which simplifies the argument.

It is known that under our equiregularity assumption, for any $x\in W$ there is an open neighborhood $U\subset W$ of $x$ and vector fields $Y_1,\ldots,Y_n$ defined in $U$ such that for any $x'\in U$ we have
\[
\text{Span}(Y_1(x'),\ldots,Y_{n_j}(x'))=\mathcal{D}_j(x')
\quad\text{for all}\ j=1,\ldots, r
\]
and
\begin{equation}\label{eq:assadd}
Y_i(x') =\tilde X_i(x') \quad\text{for all}\ i=1,\ldots,n_1.
\end{equation}
A family of vector fields satisfying the first assumption is said to be adapted to the flag at $x'$.

Given an adapted flag $(Y_1,\ldots,Y_n)$ at $x\in W$ satisfying \eqref{eq:assadd} with $x'=x$, we~can define canonical privileged coordinates of the first kind\footnote{Other choices are possible but we only need that some privileged coordinates exist and will only consider this one, which is actually the one introduced by G.\,Métivier \cite{Me} in the proof of his Theorem~3.1.} at $x$ by the mapping $\theta_x$ given by
\begin{equation} \label{eq:pc}
\theta_x (y) := u = (u_{i}) \quad\mbox{ if } y = \exp (\sum_{i=1}^n u_{i}Y_i) \cdot x,
\end{equation} where $\exp$ denotes the exponential map defined in some small neighborhood of $x$.
Thus we identify a neighborhood of $x\in M$ via $\theta_x$ with a neighborhood of $0$ in $\mathbb R^n$. It~has been shown by G.\,Métivier (see below) that everything depends smoothly on~$x$. In~particular, $\theta_x$ is also $C^\infty$ with respect to $x$.

We denote by $Y_{i,x}$ the image of $Y_i$ by $\theta_x$,
which is simply $Y_i$ written in the local canonical coordinates around $x$. Thus $Y_{i,x}$ is a vector field defined in an open neighborhood of $0$ in $\R^n$.

On $\mathbb R^n$, with coordinates $u=(u_{i})$, we~introduce the family of dilations given by
\begin{equation}
\label{eq:dilation}
\delta_t (u_{i}) = (t^{w_i} u_{i}),
\end{equation}
where positive integers $w_1,\ldots,w_n$ are defined as follows: for any $i\in\{1,\ldots,n\}$ there is a unique $j\in\{1,\ldots,n\}$ such that $n_{j-1}+1 \leq i\leq n_j$, and we set $w_i=j$. We~note that the homogeneous dimension $Q$, defined in \eqref{eq:defQ}, satisfies
\begin{equation}\label{eq:defQa}
Q =\sum_{i=1}^n w_i.
\end{equation}
Via the family of dilations we have a natural definition of homogeneous functions of degree $s$ on $\mathbb R^n\setminus\{0\}$, a definition of ``homogeneous norm'' (of degree one) and corresponding notions of vanishing function to order $p$.
A differential operator of the form
$f(u) \sfrac{\partial}{\partial u_{i}}$ is of order $w_i -s$ if $f$ is homogeneous of degree $s$.\footnote{Here we follow the sign convention in \cite{Me}, which is the opposite of that of many works, including~\cite{Je}.} When $f$ is defined in a pointed neighborhood of $0$ and can be expanded into a sum of homogeneous terms of increasing order, we~will say that $f$ is of order $\leq s$ if the term of lowest order is homogeneous of degree $s$.

G.\,Métivier \cite[Th.\,3.1]{Me} proves the following theorem (in addition to the regularity of $\theta_x$ already mentioned above).

\begin{theorem}\label{lemmaMet}
For any $x$ and $j=1,\ldots,n_1$, $\tilde X_{j,x}$ is of order $\leq 1$. Furthermore,
\begin{itemize}
\item For $j=1,\ldots,n_1$ we have
\[
\tilde X_{j,x} = \widehat X_{j,x} + R_{j,x},
\]
where $\widehat X_{j,x}$ is homogeneous of order $1$ and $R_{j,x}$ is of order $\leq 0$.
\item
The $\widehat X_{j,x}$, $j=1,\ldots,n_1$, generate a nilpotent Lie algebra $\mathcal G_x$ of dimension $n$ and rank $r$.
\item
The mapping $x \mto \widehat X_{j,x}$ is smooth.
\end{itemize}
\end{theorem}

By the nilpotent approximation, we~can associate with each point $x\in W$ a nilpotent Lie group $G_x$ (identified with the Lie algebra $\mathcal G_x$ in the $u$-coordinates) and a corresponding sub-Laplacian
\[
\widehat \Delta_x= \sum_{j=1}^{n_1} \widehat X_{j,x}^{\,2}
\]
in $\mathcal U_2(\mathcal G_x)$ (the elements in the enveloping algebra $\mathcal U(\mathcal G_x)$ that are homogeneous of degree $2$). The Hörmander condition (Assumption \ref{CHOr}) for the vector fields $\tilde X_1,\ldots,\tilde X_{n_1}$ implies that for every $x\in W$ the vector fields $\widehat X_{1,x},\ldots,\widehat X_{n_1,x}$ satisfy the corresponding Hörmander condition on $\R^n$.

At this point it is important to notice that for any $x\in W$ the operator $-\widehat\Delta_x$ depends only on the vector fields $X_1,\ldots,X_p$ and \emph{not} on the choice of the auxiliary vector fields $\tilde X_1,\ldots,\tilde X_{n_1}$ satisfying \eqref{eq:loclinindep} and \eqref{eq:loclinindeplapl}. This claim is justified in \hbox{\cite[\S A.5.4]{CHT1}}. The choice of other auxiliary vector fields corresponds to a different choice of privileged coordinates. In~passing we note that this argument also shows that instead of the above canonical privileged coordinates of the first kind one can also choose other so-called privileged coordinates, leading to the same operator and measure; see also the example in Section~\ref{AppA}.

Next, we~introduce the notion of nilpotentized measure. That is, given the measure~$\mu$ on $M$, for any point $x\in W$ we define a measure $\widehat \mu_x$ on $\R^n$. We~refer to \cite[App.\,A.5.6]{CHT1} or to our Appendix \ref{AppC} for a definition in the formalism of sub-Riemannian geometry and explain here ``by hand'' how it can be constructed for our specific choice of privileged coordinates.
On $\mathbb R^n$ we have the Lebesgue measure
\[
du= \prod_{i=1}^n du_{i},
\]
and in these local coordinates the measure $\mu$ is of the form
\[
d\mu = a(x,u) \,du,
\]
where $(x,u)\mto a(x,u)$ is $C^\infty$ in both variables $x$ and $u$. In~a small neighborhood of~$0$, the nilpotentized measure at $x$ can be defined by
\begin{equation}\label{eq:nilpmeas}
d\widehat \mu_x:= a(x,0) \,du.
\end{equation}
Note that for $u$ small and locally in $x$ we have a good control on $a(x,u)/a(x,0)$ and its inverse.

It is important to note that for every $x\in M$, the nilpotentized measure $\widehat\mu_x$ is invariant with respect to the group operation on $G_x$. As~a consequence, the formal adjoint $(\widehat X_{j,x})^\star$ of $\widehat X_{j,x}$ with respect to the scalar product in $L^2(G_x,\widehat\mu_x)$ is equal to $-\widehat X_{j,x}$; see, e.g., \cite[Rem.\,A.5]{CHT1}. As~a consequence, we~have
\[
-\widehat \Delta_x = \sum_{j=1}^{n_1} (\widehat X_{j,x})^\star \widehat X_{j,x}
\]
and this operator is selfadjoint in $L^2(G_x,\widehat\mu_x)$. We~note that this sub-Laplacian arises from the construction in the introduction when we replace $M$ by $G_x=\R^n$, $\mu$ by $\widehat\mu_x$ and $X_1,\ldots,X_p$ by $\widehat X_{1,x},\ldots,\widehat X_{n_1,x}$.

The nilpotentized measure depends only on $\mu$ and the vector fields $X_1,\ldots, X_p$ and \emph{not} on the choice of the auxiliary vector fields $\tilde X_1,\ldots,\tilde X_{n_1}$ satisfying \eqref{eq:loclinindep} and \eqref{eq:loclinindeplapl}. This claim is implicit in \cite[\S A.5.6]{CHT1}, where the nilpotentized measure is defined through arbitrary privileged coordinates. The choice of other auxiliary vector fields corresponds to a different choice of privileged coordinates.

This concludes our presentation of the nilpotent approximation on $W$. In~order to obtain a nilpotent approximation on all of $M$ we apply Lemma \ref{locindep} to cover $M$ by open sets $W$ on which \eqref{eq:loclinindep} and \eqref{eq:loclinindeplapl} are satisfied. The fact that $-\widehat\Delta_x$ and $\widehat\mu_x$ depend only on $\mu$ and $X_1,\ldots,X_p$ implies that when $x$ belongs to two different sets $W$, then the corresponding nilpotentized sub-Laplacians and nilpotentized measures coincide. Therefore the nilpotent approximation is well defined on $M$.

\subsection{The local Weyl constant}
We now turn our attention to the asymptotic distribution of the eigenvalues of the operator $-\Delta_{\Xb}^{M,\mu}$, which is given by Métivier's Weyl formula \cite{Me} that we have already mentioned in \eqref{MetForm1}. We~will describe the constant that appears in this asymptotic formula.

For each fixed $x\in M$, we~consider the selfadjoint operator $-\widehat \Delta_x$ in $L^2(G_x,\widehat\mu_x)$. Its spectral projections $\1(-\widehat\Delta_x<\lambda)$, $\lambda>0$, are integral operators in $G_x$, that is,
\[
(\1(-\widehat\Delta_x<\lambda) f)(u) = \int_{G_x} \1(-\widehat\Delta_x<\lambda)(u,v) f(v)\,d\widehat\mu_x(v)
\]
for all $u\in G_x$ and $f\in L^2(G_x,\widehat\mu_x)$ with a certain integral kernel $\1(-\widehat\Delta_x<\lambda)(u,v)$, the \emph{spectral function}. Since the operator $-\widehat\Delta_x$ is invariant under the group operation in $G_x$, the integral kernel satisfies
\[
\1(-\widehat\Delta_x<\lambda)(u,u) = \1(-\widehat\Delta_x<\lambda)(0,0)
\quad\text{for all}\ u\in G_x.
\]
Moreover, since $-\widehat\Delta_x$ is homogeneous of degree $-2$ under dilations in $G_x$, we~deduce that
\begin{equation}
\label{eq:specfcndiag1}
\1(-\widehat\Delta_x<\lambda)(0,0) = c_x^\Weyl \ \lambda^{\sfrac Q2}
\quad\text{for all}\ \lambda>0
\end{equation}
with
\begin{equation}
\label{eq:specfcndiag2}
c_x^\Weyl := \1(-\widehat\Delta_x<1)(0,0).
\end{equation}
Explicit formulas for $c_x^\Weyl$ can be obtained in certain special cases, for instance, in the case where $G_x$ is a Heisenberg group; see Section~\ref{sec:weylhnrk}. In~general, it~is known that~$c^\Weyl_x$ is positive for every $x\in M$ and that $x\mto c^\Weyl_x$ is continuous.

Using this definition we can state a more precise version of \eqref{MetForm1}.

\begin{theorem}\label{weylmet}
The spectral counting function of the selfadjoint realization of $-\Delta^{M,\mu}_\bX$ in $L^2(M,\mu)$ satisfies, as $\lambda\to +\infty$,
\begin{equation}\label{eq:met1}
N(\lambda,-\Delta_{\Xb}^{M,\mu}):= \# \{j:\ \lambda_j (-\Delta_{\Xb}^{M,\mu}) \leq \lambda\} \sim \Bigl(\int_M c_x^\Weyl \,d\mu(x) \Bigr)\,\lambda^{\sfrac Q2}.
\end{equation}
\end{theorem}

As we already mentioned, this result is due to Métivier \cite{Me}. In~the special case where $r=2$, there were important contributions on the subject starting from the end of the seventies \cite{MS1,MS2,Mo}. For recent developments related to Theorem \ref{weylmet} we refer to \cite{CHT0,CHT,CHT1}; see also our Appendix \ref{AppC} for its relation to \cite{CHT1}.

\begin{remark}\label{rem:weylsubsets}
Theorem \ref{weylmet} remains valid when we consider the Dirichlet realization of the operator $-\Delta_{\Xb}^{M,\mu}$ in an open set $\Omega\subset M$. In~this case the integral on the right side of \eqref{eq:met1} is restricted to $\Omega$.
\end{remark}

\begin{remark}
The constant $c_x^\Weyl$ itself appears through a Weyl-type formula. Indeed, let $x\in M$ and let $\Omega\subset G_x$ be open with $\widehat\mu_x(\Omega)<\infty$. Then, if~$-\widehat\Delta_x|_\Omega$ is the Dirichlet realization of $-\widehat \Delta_x$ on $\Omega$, one has
\[
N(\lambda,-\widehat\Delta_x|_\Omega) \sim c_x^\Weyl \ \widehat\mu_x(\Omega) \ \lambda^{\sfrac Q2}
\quad\text{as}\ \lambda\to\infty.
\]
This follows from the previous remark since the nilpotentization of the sub-Laplacian on a nilpotent group is the sub-Laplacian itself.
\end{remark}

\begin{remark}\label{rem:weylindepmeas}
The integral $\int_M c_x^\Weyl \,d\mu(x)$ depends on $M$ and on the vector fields $X_1,\ldots, X_p$, but is independent of the measure $\mu$. This can probably be extracted from \cite{Me} and is made explicit in \cite[\S 2.1]{CHT1}.
\end{remark}

\subsection{The local Faber--Krahn constant}
As will be discussed in detail below, for any $x\in M$ we have a Faber--Krahn inequality on the nilpotent group $G_x$, that is, for every $x\in M$ there is a constant $c>0$ such that
\begin{equation} \label{eq:FKx1}
\langle - \widehat \Delta_x v,v\rangle_{L^2(G_x,\widehat\mu_x)} \!\geq c\,\widehat\mu_x(\Omega)^{-\sfrac 2Q} \| v\|_{L^2(G_x,\widehat\mu_x)}^2,
\quad \forall \Omega\!\subset\! G_x \ \text{open},\,\forall v\!\in\! C_\rc^\infty (\Omega).
\end{equation}
We recall that $Q$ denotes the homogeneous dimension of $G_x$; see \eqref{eq:defQ}. By~our assumption of equiregularity, $Q$ is independent of $x\in M$.

By definition, $c_x^\FK$ is the largest constant such that \eqref{eq:FKx1} holds. We~will prove momentarily the positivity of this constant, even uniformly in $x$.

\begin{remark}\label{rem:fkindepmeas}
Having in mind the proof of Pleijel's theorem, it~is important to write the above estimates using the appropriately normalized Lebesgue measure $\widehat \mu_x$ on $G_x$. Note that the Faber--Krahn constant $c_x^\FK$ depends both on $G_x$ and on a specific normalization constant determined by the measure $\mu$; see \cite[App.\,A.5.6]{CHT1} and our Appendix \ref{AppC}. When $\tilde\mu$ is a second measure on $M$ satisfying the same properties as $\mu$, then $d\tilde\mu = h\,d\mu$ for a smooth, positive function $h$ on $M$. Then, for any $x\in M$,
\[
d\widehat{\tilde\mu}_x = h(x) \,d\widehat\mu_x,
\]
which shows that
\[
\tilde c^\FK_x = h(x)^{\sfrac 2Q} \,c^\FK_x.
\]
As an important consequence, we~see that the integral
\[
\int_M \left(c^\FK_x\right)^{-\sfrac Q2}d\mu(x)
\]
depends on $M$ and $X_1,\ldots, X_p$, but is independent of the measure $\mu$.\footnote{We are grateful to Y.\,Colin de Verdière for pointing this out to us.}
\end{remark}

Further information about the Faber--Krahn constant is contained in the following lemma, whose proof we defer to the Subsection \ref{sec:fk}.

\begin{lemma}\label{faberkrahn}
We have
\begin{equation}\label{eq:fk1}
\inf_{x\in M} c_x^\FK >0.
\end{equation}
Moreover, the function $M\ni x\mto c_x^\FK$ is uniformly Hölder continuous.
\end{lemma}

\section{Main result for sub-Laplacians in the equiregular case}\label{sec:mainsubriem}

We continue to work in the setting of the previous section. In~particular, we~suppose that Assumptions \ref{CHOr} and \ref{equireg} are satisfied. Our main statement concerning nilpotent approximation is the following theorem:

\begin{theorem}\label{t:manifold}
Let $-\Delta_\Xb^{M,\mu} =\sum_{j=1}^p X_j^\star X_j$ be an equiregular sub-Riemannian Laplacian on a closed connected manifold $M$ with given measure $\mu$. Then
\begin{equation}\label{eq:mainmfd}
\limsup_{k\to+\infty} \frac{\nu_k}{k} \leq \biggl(\int_M \bigl(c^\FK_x \bigr)^{-\sfrac Q2} \,d\mu(x) \biggr) \biggl(\int_M c_x^\Weyl \,d\mu(x) \biggr)^{-1},
\end{equation}
where $\nu_k$ denotes the maximal number of nodal domains of an eigenfunction of $-\Delta_\Xb^{M,\mu}$ associated with eigenvalue $\lambda_k$.
\end{theorem}

\begin{corollary}\label{maincor}
If
\begin{equation}\label{eq:hyp}
\biggl(\int_M \bigl(c^\FK_x \bigr)^{-\sfrac Q2} \,d\mu(x) \biggr) \biggl(\int_M c_x^\Weyl \,d\mu(x) \biggr)^{-1} < 1,
\end{equation}
then Pleijel's theorem holds. In~particular, if
\[
\bigl(c^\FK_x \bigr)^{-\sfrac Q2} \,\bigl(c_x^\Weyl \bigr)^{-1} < 1
\quad\text{for all}\ x\in M,
\]
then \eqref{eq:hyp} and therefore Pleijel's theorem holds.
\end{corollary}

The first part of the corollary follows immediately from Theorem \ref{t:manifold}. For the second part, we~note that both $x\mto c^\FK_x$ and $x\mto c_x^\Weyl$ are continuous, so under the assumption of the second part of the corollary there is a constant $\gamma<1$ such that $\bigl(c^\FK_x \bigr)^{-\sfrac Q2} \,\bigl(c_x^\Weyl \bigr)^{-1}\leq\gamma$ for all $x\in M$. This implies that the left side of \eqref{eq:hyp} is at most $\gamma$.

Theorem \ref{t:manifold} and Corollary \ref{maincor} are improvements due to Y.\,Colin de Verdière of bounds that appeared in a preprint version of this paper \cite{FrHe1}; see also the announcement \cite{FrHe2}. There we had $(\inf_{x\in M} c_x^\FK)^{-\sfrac Q2} \mu(M)$ instead of $\int_M \bigl(c^\FK_x \bigr)^{-\sfrac Q2} \,d\mu(x)$. The usefulness of this improvement can be seen, for instance, in the examples in Section~\ref{AppA}. We~are very grateful to Y.\,Colin de Verdière for allowing us to incorporate his ideas into our paper.

\begin{remark} \label{remdom}
Theorem \ref{t:manifold} remains valid when we consider the Dirichlet realization $-\Delta^{\Omega,\mu}_\Xb$ of an equiregular sub-Riemannian Laplacian in a relatively compact open set~$\Omega$ in a manifold $M$. The proof relies on Remark \ref{rem:weylsubsets}. It~is enough to have Assumptions~\ref{CHOr} and~\ref{equireg} satisfied in a neighborhood of $\overline{\Omega}$.
Examples are discussed in Section~\ref{AppA}. In~particular, we~can consider an open, relatively compact set $\Omega$ in a fixed stratified group~$G$, where $\Xb$ is a basis of $\mathcal G_1$
and $\mu$ is (in the exponential coordinates) the Lebesgue measure. In~this case, the function $x\mto c^\Weyl_x$ is constant. Situations of this type are further discussed in Part \ref{part2}.
\end{remark}

\begin{remark}\label{remcourantasymp}
The weak version of Courant's theorem (that is, the first part of Theorem~\ref{courant}) yields
\begin{equation}\label{eq:weakcourant}
\limsup_{k\to+\infty} \frac{\nu_k}k \leq 1.
\end{equation}
Here we use the fact that Weyl's formula implies $\lim_{k\to\infty} \sfrac{\mult(\lambda_k)}{k}=0$. In~fact, the asymptotic bound \eqref{eq:weakcourant} remains valid under assumptions much weaker than the equiregularity assumption (Assumption \ref{equireg}), which we assume in this paper. It~suffices that the singular set $\mathcal S$ is Whitney stratified by equisingular smooth submanifolds and that $\mathcal D_1$ is $\mathcal S$-nilpotentizable, because under these assumptions a Weyl law was established in \cite{CHT1}, from which it follows that $\lim_{k\to\infty} \sfrac{\mult(\lambda_k)}{k}=0$. We~refer to~\cite{CHT1} for a definition of the notions used here. Particular cases where these assumptions are satisfied are the Baouendi--Grushin case \cite[Th.\,7.2]{CHT1} and the Martinet case \cite[Th.\,7.3]{CHT1}. In~this connection we emphasize that our proof of Theorem \ref{t:manifold} relies heavily on Assumptions \ref{CHOr} and \ref{equireg}. Whether or not this theorem can be extended to some non-equiregular situations (for instance those under which Weyl's law was established in \cite{CHT1}) remains an \emph{open problem}.
\end{remark}

In view of the previous remark, Theorem~\ref{t:manifold} is only interesting for its application in the corollary, that is, when \eqref{eq:hyp} holds. In~Part \ref{part2} we will investigate the validity of this condition in the case of an open subset of $\mathbb H_n\times\R^k$. Other instances where one might be able to prove the validity of \eqref{eq:hyp} are, for example, in the form $M_3\times \mathbb T^k$ with $k$ large enough and with $M_3$ a compact $3$-dimensional contact manifold.

\begin{remark}
In the Riemannian case (that is, when $p=n$ and when $\mu$ is the Riemannian volume measure) the assertion of Theorem \ref{t:manifold} reduces to the theorem of Bérard and Meyer \cite{BerMe}. Indeed, in this case we have $Q=n$ and for every $x\in M$ the operator $-\widehat\Delta_x$ is the ordinary Laplacian on $G_x=\R^n$. Consequently, in the notation of the introduction and of Part \ref{part2},
\[
c^\FK_x = C^\FK(\R^n)
\quad\text{and}\quad
c_x^\Weyl = \mathcal W(\R^n).
\]
Both quantities are independent of $x$ and we arrive at the same bound with constant
\[
\gamma(\R^n) = (C^\FK(\R^n))^{-\sfrac Q2} \,(\mathcal W(\R^n))^{-1}
\]
as in the case of domains in Euclidean space; see \eqref{eq:pleijelintro}.
\end{remark}

\begin{remark}\label{rem:indepmeas}
The upper bound in \eqref{eq:mainmfd} depends on the manifold $M$ and on the vector fields $X_1,\ldots, X_p$ (up to orthogonal transformations), but it does \emph{not} depend on the measure $\mu$. This follows from Remarks \ref{rem:weylindepmeas} and \ref{rem:fkindepmeas}. Indeed, according to these remarks both integrals on the right side of \eqref{eq:mainmfd} are independent of $\mu$.
\end{remark}

\section{Proof of the main sub-Riemannian results}
\label{sec:mainsubriemproof}

\subsection{Comparing Laplacians}\label{sec:comparinglaplacians}
Throughout this section we choose a Riemannian structure on $M$ that is compatible with its smooth structure. (This is always possible, first in local coordinates and then globally via a partition of unity.) This Riemannian structure allows us to consider (open) geodesic balls $B(x,r)$ at $x\in M$ of radius $r>0$.

Recall that our nilpotent approximation in Section~\ref{s1} was carried out on open subsets $W\subset M$ where \eqref{eq:loclinindep} and \eqref{eq:loclinindeplapl} are satisfied. Also in this subsection we will work locally. More precisely, for an open set $W\subset M$ as before we choose a compact subset $K \subset W$. By~compactness, there is an $\epsilon_*>0$ such that for any $x\in K$ the privileged coordinates at $x$ are well defined in $B(x,\epsilon_*)$.

In the formulation of the following lemma we identify functions $v$ on $M$ with support in $B(x,\epsilon_*)$ with functions $v\circ\theta_x^{-1}$ on $G_{x}$ with support in $\theta_x(B(x,\epsilon_*))$.

\begin{lemma}\label{l:remainder2}
Fix $W$ and $K$ as above. Then there are constants $C$, $\epsilon_0>0$ and $s>0$ such that for any $x\in K$, any $0<\epsilon\leq\epsilon_0$ and any function $v\in C^\infty_\rc(B(x,\epsilon))$ one has
\begin{align*}
\left| \langle \widehat{\Delta}_{x} v, v \rangle_{L^2(G_{x,\widehat\mu_{x}})} - \langle \Delta_\Xb^{M,\mu} v, v \rangle_{L^2(M,\mu)}\right| \leq C \epsilon^s \langle - \widehat{\Delta}_{x} v, v \rangle_{L^2(G_{x},\widehat\mu_{x})}.
\end{align*}
\end{lemma}

\begin{proof}
Throughout the proof, we~will make use of the uniformity with respect to~$x$ of several geometric constructions around a point $x\in K$. This is discussed and proved in \cite{Me,HN, HN1,Rot} and will be used freely in what follows.

We give the proof in three steps.

\subsubsection*{Step 1: Change of the measure}
We denote by $\tilde\nabla$ the sub-Riemannian gradient, so
\[
-\langle\Delta_\Xb^{M,\mu} v,v\rangle_{L^2(M,\mu)} = \sum_{j=1}^p \|X_j v\|_{L^2(M,\mu)}^2 = \sum_{j=1}^{n_1} \|\tilde X_j v\|_{L^2(M,\mu)}^2 = \| \tilde\nabla v \|_{L^2(M,\mu)}^2.
\]
Due to the localization of the support of $v$, we~have
\begin{align*}
(1-C\varepsilon) \| \tilde\nabla v \|_{L^2(G_x,\widehat\mu_x)}^2
\leq \| \tilde\nabla v \|_{L^2(M,\mu)}^2
\leq (1+C\varepsilon) \| \tilde\nabla v \|_{L^2(G_x,\widehat\mu_x)}^2.
\end{align*}

\subsubsection*{Step 2: Comparing $\tilde X_{i,x}$ and $\widehat X_{i,x}$}
Due to the compactness of $M$, there exists $\epsilon_0>0$ such that, for each $x\in M$ we have a localization function $\tilde \chi_x$ such that $\tilde \chi_x = 1$ on $B(x, 2\epsilon_0)$, $\supp(\tilde \chi_x) \subset B (x,4\epsilon_0)$ and all the estimates on the derivatives are controlled uniformly. After replacing $\epsilon_0$ by $\min\{\epsilon_0,\frac14\epsilon_*\}$ if necessary, we~may assume that for each $x\in K$ the privileged coordinates at $x$ are well defined in $B(x,4\varepsilon_0)$.

For technical reasons we also have to introduce another cut-off function $\widehat \chi_x$ of the same type, viz.\ such that $\widehat \chi_x = 1$ on $B(x, \epsilon_0)$, $\supp(\widehat \chi_x) \subset B (x,2\epsilon_0)$ and with uniform bounds on the derivatives.

We observe that in the privileged coordinates centered at $x$ we have\footnote{See either Helffer--Nourrigat \cite[Prop.\,5.1]{HN}, Métivier \cite{Me}, \cite[\S 4.2.4]{EL}, or \cite[\S 1]{Rot}.} by Lemma~\ref{lemmaMet} that
$
\tilde \chi_x (\tilde X_{i,x} -\widehat X_{i,x}) \tilde\chi_x
$
is of degree $0$.

There is also a notion of type\footnote{The authors introduce first a notion of function of type $\lambda$ corresponding to a function that is homogeneous of degree $-Q +\lambda$ with respect to the dilation (with an addition condition when $\lambda=0$). By~integration against test functions, this defines a distribution of type $\lambda$. More generally, assuming that $W\subset M$ is such that $\theta_x(y)$ is defined for all $x,y\in W$, then a function $K$ on $W\times W$ is a kernel of type $\lambda$ if for any $\ell \geq 0$, we~have
\[
K(x,y) = \sum_{i=1}^s a_i(x) k_x^{(i)} (\theta_x(y)) b_i(y) + E^\ell (x,y),
\]
where $a_i, b_i \in C_\rc^\infty(W)$, where $k_x^{(i)}$ is a kernel of type $\geq \lambda$ with $(x,u)\mto k_x^{(i)}(u)$ smooth away from $u=0$, and where $E^\ell \in C^\ell(W \times W)$. An operator of type $\lambda$ is a mapping originally defined on $C_\rc^\infty(W)$ whose distribution kernel is a kernel of type $\lambda$.}
in \cite{RS} (see \cite[pp.\,654--655]{Rot} for the adaptation to the equiregular situation), which roughly speaking corresponds to the operator of ``nonpositive'' degree.

Note that (after localization around $0$ in the privileged coordinates) $(-\widehat \Delta_x)^{-1}$ is an operator of type 2 and that $\widehat X_{m,x}(-\widehat \Delta_x)^{-1} $ is an operator of type $1$ \cite[Th.\,8]{RS}. Then $\tilde \chi_x (\tilde X_{i,x} -\widehat X_{i,x}) \tilde\chi_x$ is an operator of type 0. Consequently, for each $m=1,\ldots,n_1$,
\[
K_{im,x} := \widehat \chi_x \tilde \chi_x (\tilde X_{i,x} -\widehat X_{i,x}) \tilde\chi_x (-\widehat \Delta_x)^{-1} \widehat X_{m,x}
\]
is an operator of type 1.

For $v \in C_\rc^\infty (G_x)$, we~have the identity
\begin{align*}
\widehat \chi_x (\tilde X_{i,x} -\widehat X_{i,x}) \tilde\chi_x v & = - \sum_m \widehat \chi_x \tilde \chi_x (\tilde X_{i,x} -\widehat X_{i,x}) \tilde\chi_x (-\widehat \Delta_x)^{-1} \widehat X_{m,x} \widehat X_{m,x} v \\
& = \sum_m K_{im,x} \widehat X_{m,x} v.
\end{align*}

To estimate the $L^2$-norm of $ K_{im,x} \widehat X_{m,x} v$, we~first estimate the $L^q$-norm with $\sfrac 1q = \sfrac 12 - \sfrac 1Q$. To do so, we~use \cite[Th.\,7]{RS} and obtain for $w$ with support in a fixed compact subset in Euclidean space
\[
\| K_{im,x} w \|_q \leq C \| w \|_2.
\]

Assume now that $v\in C^\infty_\rc(B(x,\epsilon_0))$ and note that, identifying $v$ with a function on $G_x$, we~have
\[
\widehat \chi_x (\tilde X_{i,x} -\widehat X_{i,x}) \tilde\chi_x v
= (\tilde X_{i,x} -\widehat X_{i,x}) v.
\]
Thus, applying the above inequality to $w= \widehat X_{m,x} v$, which has support in the fixed compact set $B(x,\epsilon_0)$, we~obtain
\begin{align*}
\| (\tilde X_{i,x} -\widehat X_{i,x})v \|_q & \leq \sum_m \| K_{im,x} \widehat X_{m,x} v \|_q \leq C \sum_m \| \widehat X_{m,x} v \|_2 \\
& \leq C \sqrt{n_1} \Bigl(\sum_m \| \widehat X_{m,x} v \|_2^2 \Bigr)^{\sfrac12}.
\end{align*}
If, moreover, we~have $v\in C^\infty_\rc(B(x,\epsilon))$ with $\epsilon\leq\epsilon_0$, then, by~Hölder's inequality, since $(\tilde X_{i,x} -\widehat X_{i,x})v$ is supported in $B(x,\epsilon)$,
\[
\| (\tilde X_{i,x} -\widehat X_{i,x})v \|_2 \leq \tilde C \epsilon^{\sfrac n Q} \Bigl(\sum_m \| \widehat X_{m,x} v \|_2^2 \Bigr)^{\sfrac12}
\]
for some $\tilde C$. Here we used the fact that the Lebesgue measure of $B(x,\epsilon)$ is a constant times $\epsilon^n$.

In the above inequality, integration was with respect to Lebesgue measure on Euclidean space, but, by~the properties of the measure $\widehat\mu_x$, the same inequality holds when integrating with respect to the latter measure.

\subsubsection*{Step 3: End of the proof of the lemma}
We can now finish the proof in the following way. We~write
\begin{align*}
& \| \tilde \nabla v \|^2_{L^2(G_x,\widehat\mu_{x})} +
\langle \widehat \Delta_{x} v, v \rangle_{L^2(G_x,\widehat\mu_{x})} \\
& = \sum_i \! \left(2 \langle (\tilde X_{i,x}-\widehat X_{i,x}) v, \widehat X_{i,x} v \rangle_{L^2(G_x,\widehat\mu_{x})} \! + \! \langle (\tilde X_{i,x}-\widehat X_{i,x}) v, (\tilde X_{i,x}-\widehat X_{i,x}) v\rangle_{L^2(G_x,\widehat\mu_{x})}\right)
\end{align*}
and bound each term on the right side using the inequality from Step 2. By~Step 1, we~have the identity $\| \tilde \nabla v \|^2_{L^2(M,\mu)} = -\langle\Delta_\Xb^{M,\mu} v,v\rangle_{L^2(M,\mu)}$, as well as a control over the difference between this and $\| \tilde \nabla v \|^2_{L^2(G_x,\widehat\mu_{x})}$. Combining these ingredients, we~arrive at the inequality in the lemma.
\end{proof}

\subsection{On the Faber--Krahn constant}\label{sec:fk}
Our goal in this subsection is to prove Lem\-ma~\ref{faberkrahn}, which contains the basic properties of the local Faber--Krahn constant~$c^\FK_x$.

\begin{proof}[Proof of Lemma \ref{faberkrahn}]
We divide the proof into two steps.

\subsubsection*{Step 1: Positivity of the Faber--Krahn constant}
We shall prove that $c_x^\FK>0$ for every $x\in M$. The uniform positivity asserted in the lemma then follows from our continuity arguments in Step 2.

Thus, let $x\in M$ and, using Lemma \ref{locindep}, choose an open neighborhood $W\subset M$ of $x$ on which the nilpotentization procedure in Section~\ref{s1} can be carried out. In~\cite[Th.\,4.24]{Rot} it is shown that for any $1<p<\sfrac Q2$ there is a $C_{p} >0$ such that for any $f \in C_\rc^\infty (\mathbb R^n)$, we~have
\begin{equation}
\label{eq:rothschild}
\| (-\widehat \Delta_x)^{-1} f \|_q \leq C_{p} \| f \|_p
\end{equation}
with $\sfrac 1q=\sfrac 1p -\sfrac 2 Q$.
(This inequality is actually used as an intermediate result toward the proof that $-\widehat \Delta_x$ has a $(-Q+2)$-homogeneous fundamental solution $k_x(u)$ that depends smoothly on $x$; see \cite[Th.\,3.6]{Rot}.)\footnote{In \cite{Rot} it is also shown that $C_p$ can be chosen locally bounded with respect to the point $x\in W$, which can be used to give a direct proof of the claimed uniform positivity of $c^\FK_x$, independent of Step 2.}

Assuming first that $Q>2$, we~can apply \eqref{eq:rothschild} with $p=\spfrac{2Q}{Q+2}$ and $q=\spfrac{2Q}{Q-2}$ and obtain
\begin{align*}
\| (-\hat\Delta_x)^{-\sfrac12} f\|_2^2 =
\int_{\R^d} f (-\widehat \Delta_x)^{-1} f \,d\widehat\mu_x &\leq \| f\|_{\spfrac{2Q}{Q+2}} \|(-\hat\Delta_x)^{-1} f\|_{\spfrac{2Q}{Q-2}}\\
&\leq C_{\spfrac{2Q}{Q+2}} \| f \|_{\spfrac{2Q}{Q+2}}^2.
\end{align*}
Since $(-\hat\Delta_x)^{-\sfrac12}$, considered as an operator $L^{\spfrac{2Q}{Q+2}}\to L^2$, and its adjoint, considered as an operator $L^2\to L^{\spfrac{2Q}{Q-2}}$, coincide on $L^{\spfrac{2Q}{Q+2}}\cap L^2$ (as a consequence of the selfadjointness of $(-\hat\Delta)^{-\sfrac12}$ on $L^2$), we~obtain
\[
\| (-\hat\Delta_x)^{-\sfrac12} g\|_{\spfrac{2Q}{Q-2}}^2 \leq C_{\spfrac{2Q}{Q+2}} \| g \|_{2}^2.
\]
Substituting $u=(-\hat\Delta_x)^{-\sfrac12} g$ we obtain the Sobolev inequality
\[
C_{\spfrac{2Q}{Q+2}}^{-1} \| u \|_{\spfrac{2Q}{Q-2}}^2 \leq \langle (-\hat\Delta_x) u,u \rangle.
\]
By Hölder's inequality, as in the proof of Proposition \ref{fksob} below, this implies that $c_x^\FK\geq C_{\spfrac{2Q}{Q+2}}^{-1}$. This proves the claimed positivity when $Q>2$.

To deal with the cases $Q=1,2$ we follow an idea of Helffer and Nourrigat already used in \cite{Rot}. We~study the operator $-\hat\Delta_x - \sum_{i=1}^s \partial_i^2$ on $G_x\times\R^s$ with $s\in\N$ chosen such that $Q+s>2$. Proceeding exactly as before, we~obtain the Sobolev inequality
\[
C_{\spfrac{2Q}{Q+2}}^{-1} \| U \|_{\spfrac{2(Q+s)}{Q+s-2}}^2 \leq \langle (-\hat\Delta_x - \textstyle\sum_{i=1}^s \partial_i^2) U,U \rangle.
\]
for functions $U$ on $G_x\times\R^s$. Applying this to a product function $U = u\otimes \phi$ with a fixed function $\phi\in C^1_\rc(\R^s)$, we~obtain
\[
c\,C_{\spfrac{2Q}{Q+2}}^{-1} \| u \|_{\spfrac{2(Q+s)}{Q+s-2}}^2 \leq \langle (-\hat\Delta_x) u, u \rangle + \|u\|_2^2.
\]
with $c>0$ depending only on $\phi$. By~a simple scaling argument in $G_x$ (see, e.g., \cite[Rem.\,2.47]{FLW}), this inequality can be brought in the form of a Sobolev interpolation inequality where on the right side a geometric mean of $\langle (-\hat\Delta_x) u, u \rangle$ and $\|u\|_2^2$ appears. In~this form one can again use Hölder's inequality, similarly as in the proof of Proposition \ref{fksob} below, to obtain the desired positive lower bound on the Faber--Krahn constant.

An alternative way of proving this step could rely on Varopoulos's proof \cite{Va} of the Sobolev inequality.

\subsubsection*{Step 2: Continuity of the Faber--Krahn constant\protect\footnotemark}\footnotetext{The statement and the following proof was suggested to us by Y.\,Colin de Verdière.}
Let $W$ and $K$ be as in Subsection~\ref{sec:comparinglaplacians}. Our aim is to prove that $x\mto c_x^\FK$ is Hölder continuous on $K$. More precisely, we~will show that there are positive constants $C,\epsilon',\alpha$ (depending on $K$) such that for all $\epsilon\in(0,\epsilon']$ we have
\begin{equation}
\label{eq:fkholder}
\left| \frac{c^\FK_x}{c^\FK_{x_0}} - 1\right| \leq C \epsilon^\alpha
\quad\text{for all}\ x,x_0\in K \ \text{with}\ x\in B(x_0,\epsilon).
\end{equation}
Once we have shown this, we~can deduce the asserted Hölder continuity on $M$. Indeed, according to Lemma \ref{locindep} we can choose an open neighborhood $W$ around each point and then we can choose a slightly smaller neighborhood $W'$ with $K:=\overline{W'}\subset W$. By~compactness we can cover $M$ by finitely many sets $W'$ and then the Hölder continuity on each $K$ implies the Hölder continuity on $M$.

We also note that the pointwise positivity of $c^\FK_x$, proved in Step 1, together with~\eqref{eq:fkholder} implies its uniform positivity on $K$. This implies the uniform positivity on $M$ asserted in the lemma by compactness of $M$.

We turn now to the proof of \eqref{eq:fkholder}. Let $K\subset M$ and recall that there is an $\epsilon_*>0$ such that, for all $x\in K$, the map $\theta_x$ is defined in $B(x,\epsilon_*)$. It~maps a neighborhood of~$x$ in $M$ to a neighborhood of $0$ in $\R^n$, which we will identify with $G_x$. We~fix $x_0\in K$ and restrict our attention to $x\in B(x_0,\frac12\epsilon_*)$. Note that for such $x$ the map $\theta_x$ is defined in $B(x_0,\frac12\epsilon_*)$.

Let $U_x :=\theta_{x}(B(x_0,\frac12\epsilon_*))\subset G_x$. For a set $\Omega\subset U_{x_0}\subset G_{x_0}$ and a function $u\in C^\infty_\rc(\Omega)$, let
\[
\Omega_x:= \theta_x\circ \theta_{x_0}^{-1}(\Omega)
\quad\text{and}\quad
u_x := u \circ \theta_{x_0} \circ \theta_x^{-1}.
\]
Then $u_x\in C^\infty_\rc(\Omega_x)$. Let $\epsilon_0$ be as in Lemma \ref{l:remainder2}. In~the following we consider $x\in B(x_0,\epsilon)$ with $\epsilon\leq\min\{\frac12\epsilon_*,\epsilon_0\}$. It~follows from Lemma \ref{l:remainder2} that
\begin{align*}
(1-C \epsilon^s) \langle -\widehat\Delta_{x_0} u, u \rangle_{L^2(G_{x_0},\widehat\mu_{x_0})} & \leq
\langle -\widehat\Delta_{x} u_x, u_x \rangle_{L^2(G_{x},\widehat\mu_{x})} \\
& \leq (1+C \epsilon^s) \langle -\widehat\Delta_{x_0} u, u \rangle_{L^2(G_{x_0},\widehat\mu_{x_0})}
\end{align*}
with constants $C$ and $s$ that are independent of $x$, $x_0$, $\epsilon$, $u$ and $\Omega$. Moreover,
\[
(1-C \epsilon) \| u \|_{L^2(G_{x_0},\widehat\mu_{x_0})}^2 \leq
\| u_x \|_{L^2(G_{x},\widehat\mu_{x})}^2
\leq (1+C \epsilon) \| u \|_{L^2(G_{x_0},\widehat\mu_{x_0})}^2
\]
with a (possibly different) constant $C$, but again independent of $x$, $x_0$, $\epsilon$, $u$ and $\Omega$. Note that the mapping $u\mto u_x$ is a bijection from $C^\infty_\rc(\Omega)$ to $C^\infty_\rc(\Omega_x)$. Therefore, combining the above bounds with the variational characterization of the first eigenvalue, we~infer that
\begin{equation}
\label{eq:fkcontproof1}
\left| \frac{\lambda_1(\Omega_x)}{\lambda_1(\Omega)} - 1\right| \leq C \epsilon^{\min\{s,1\}}.
\end{equation}
Here $\lambda(\Omega_x)$ and $\lambda(\Omega)$ denote the first eigenvalues of the Dirichlet realizations of $-\widehat\Delta_x$ in $L^2(\Omega_x,\widehat\mu_x)$ and of $-\widehat\Delta_{x_0}$ in $L^2(\Omega,\widehat\mu_{x_0})$, respectively. Clearly, we~also have
\begin{equation}
\label{eq:fkcontproof2}
\left| \frac{\widehat\mu_x(\Omega_x)}{\widehat\mu_{x_0}(\Omega)} - 1\right| \leq C \epsilon.
\end{equation}
We deduce that
\[
\left(\widehat\mu_x(\Omega_x)\right)^{\sfrac2Q} \lambda_1(\Omega_x)
\geq \bigl(1- C \epsilon^{\min\{s,1\}}\bigr) \left(\widehat\mu_{x_0}(\Omega)\right)^{\sfrac2Q} \lambda_1(\Omega) \geq \bigl(1- C \epsilon^{\min\{s,1\}}\bigr) c^\FK_{x_0}.
\]
Since the map $\Omega\mto\Omega_x$ is a bijection from open subsets of $U_{x_0}$ to open subsets of $U_x$, we~obtain the inequality
\begin{equation}
\label{eq:fkcontproof}
\left(\widehat\mu_x(\omega)\right)^{\sfrac2Q} \lambda_1(\omega)
\geq \bigl(1- C \epsilon^{\min\{s,1\}} \bigr) c^\FK_{x_0}
\end{equation}
for any open $\omega\subset U_x$.

Recall that we have dilations on $G_x$. Under a dilation of $\omega$, the eigenvalue $\lambda_1(\omega)$ of $-\widehat\Delta_x$ is homogeneous of degree $-2$, while the measure $\widehat\mu_x(\omega)$ is homogeneous of degree $Q$. Therefore the product $\left(\widehat\mu_x(\omega)\right)^{\sfrac2Q} \lambda_1(\omega)$ is homogeneous of degree zero. Since $0\in U_x$ (as $x\in B(x_0,\frac12\epsilon_*)$), we~can dilate any bounded set $\omega\subset G_x$ (where boundedness is understood, for instance, with respect to the Euclidean metric on $G_x$ identified with $\R^n$) so that it becomes a subset of $U_x$ and we deduce that \eqref{eq:fkcontproof} holds for any bounded open set $\omega\subset G_x$. Finally, the boundedness assumption on $\omega$ can be relaxed to a finite-measure assumption if we recall that the variational quotient defining $\lambda_1(\omega)$ only needs to be considered for functions in $C^\infty_\rc(\omega)$. To summarize, we~have shown that \eqref{eq:fkcontproof} holds for any open set $\omega\subset G_x$ of finite measure. This proves that
\[
c_x^\FK \geq \bigl(1- C \epsilon^{\min\{s,1\}} \bigr) c^\FK_{x_0}.
\]
The analogous inequality where the roles of $x$ and $x_0$ are interchanged is deduced from \eqref{eq:fkcontproof1} and \eqref{eq:fkcontproof2} in essentially the same way. This concludes the proof of (Hölder continuity) of $x\mto c^\FK_x$ at $x_0$. Moreover, the constants in this Hölder continuity bound only depend on $K$, as claimed.
\end{proof}

\subsection{The Faber--Krahn inequality on small sets}
We now come to the main step in the proof of Theorem \ref{t:manifold}. As~in the work of Bérard and Meyer~\cite{BerMe}, the idea is to prove a Faber--Krahn inequality where the constant is ``almost'' the ``good'' constant, provided the sets on which the inequality is applied are ``small''. For us, the ``good'' constant is in fact a function on $M$, namely $x\mto c^\FK_x$, and we capture the variation of this constant in terms of the measure $(c_x^\FK)^{-\sfrac Q2}\mu$. The ``smallness'' of sets is understood with respect to their $\mu$-measure. The precise statement is the following:

\begin{proposition}\label{faberkrahnbm}
For any $\theta>0$ there is an $\eta>0$ such that for any open set $\Omega\subset M$ with $\mu(\Omega)\leq\eta$ and any $v\in C^\infty_\rc(\Omega)$,
\begin{equation}
\langle -\Delta^{M,\mu}_\Xb v,v \rangle \geq (1-\theta) \left(\int_\Omega \left(c_x^\FK\right)^{-\sfrac Q2} d\mu(x)\right)^{-\sfrac2Q} \| v\|^2.
\end{equation}
\end{proposition}

\begin{proof}
Recall that by Lemma \ref{locindep} any point $a\in M$ has a neighborhood $W_a$ where the nilpotentization procedure in Section~\ref{s1} can be carried out. We~fix an open neighborhood $W_a'$ of $a$ with $\overline{W_a'}\subset W_a$. By~compactness there are finitely many points $a_1,\ldots,a_L\in M$ such that $\bigcup_{\ell=1}^L W_{a_\ell}' = M$. We~apply Lemma \ref{l:remainder2} with $W_{a_\ell}$ and $\overline{W_{a_\ell} '}$ in place of $W$ and $K$ and obtain constants $C$, $\epsilon_0$ and $s$ such that the conclusion of that lemma holds. We~may and will assume that these constants are independent of $\ell$.

For each $\varepsilon\in(0,\epsilon_0]$, we~introduce a family of smooth cut-off functions $\chi_j:M\to \R$ such that
\begin{itemize}
\item $\sum_j \chi_j^2= 1$ everywhere,
\item for each $j$ there exists $x_j = x_j(\epsilon)\in M$ with $\supp (\chi_j) \subset B(x_j(\epsilon),\epsilon)$,
\item there exists $C>0$ (independent of $\varepsilon>0$) such that everywhere in $M$,
\[
\sum_j |\nabla \chi_j|^2 \leq C \varepsilon^{-2}.
\]
\end{itemize}
For any $v\in C^\infty(M)$ we have the identity
\begin{align}\label{e:identity}
\langle -\Delta^{M,\mu}_\Xb v,v\rangle_{L^2(M)} & = - \sum_j \langle \chi_j \Delta^{M,\mu}_\Xb v,\chi_j v \rangle_{L^2(M)} \nonumber \\
&= \sum_j \left(- \langle [\chi_j,\Delta^{M,\mu}_\Xb]v,\chi_jv \rangle_{L^2(M)} - \langle \Delta^{M,\mu}_\Xb (\chi_jv), \chi_j v\rangle_{L^2(M)}\right) \nonumber \\
&=\sum_j \Bigl(- \| v \tilde \nabla\chi_j \|_{L^2(M)}^2 - \langle \Delta^{M,\mu}_\Xb (\chi_j v),\chi_j v \rangle_{L^2(M)} \Bigr).
\end{align}
Here, as in the previous proof, $\tilde \nabla$ is not the Euclidean gradient, but the sub-Riemannian gradient, and
\[
\| v\tilde \nabla\chi_j \|_{L^2(M)}^2 = \sum_{\ell} \| v X_\ell \,\chi_j \|_{L^2(M)}^2.
\]
Identity \eqref{e:identity} is a sub-Riemannian version of the IMS localization formula in mathematical physics; see, for instance, \cite[\S 3.1]{CyFrKiSi}.

Note first that, by~our construction of the $\chi_j$ and the compactness of $M$, there exists a constant $C_1$ such that
\[
\sum_{\ell} \| v X_\ell \,\chi_j \|_{L^2(M)}^2 \leq C_1 \epsilon^{-2} \|v\|_{L^2(M)}^2.
\]
Let $\eta>0$ be a parameter that will be chosen later depending on $\epsilon$. Assuming that $v\in C^\infty_\rc(\Omega)$ with $\mu(\Omega)\leq\eta$, we~use the previous bound to get
\begin{multline}
\label{eq:remainder1}
\sum_{\ell} \| v X_\ell \,\chi_j \|_{L^2(M)}^2\\[-8pt]
\leq C_1 \epsilon^{-2}
\eta^{\sfrac2Q} \left(\inf_{x\in M} c_x^\FK\right)^{-1}
\left(\int_\Omega \left(c^\FK_x\right)^{-\sfrac Q2} d\mu(x)\right)^{-\sfrac 2Q}
\|v\|_{L^2(M)}^2.
\end{multline}
Note that the infimum on the right side is finite by Lemma \ref{faberkrahn}.

We turn our attention to the last term in \eqref{e:identity}. Note that Lemma \ref{l:remainder2} is applicable since $\epsilon\leq\epsilon_0$ and since for any $j$ there is an $\ell$ with $x_j\in W_{a_\ell}'$. We~infer that
\begin{multline*}
\left| \langle \Delta^{M,\mu}_\Xb \chi_j v, \chi_j v \rangle_{L^2(M,\mu)} - \langle \widehat{\Delta}_{x_j} \chi_j v, \chi_j v \rangle_{L^2(G_{x_j,\widehat\mu_{x_j}})}\right| \\
\leq C_2 \varepsilon^{s} \langle - \widehat{\Delta}_{x_j} (\chi_j v), \chi_j v \rangle_{L^2(G_{x_j},\widehat\mu_{x_j})}.
\end{multline*}
We combine this bound with the local Faber--Krahn inequality \eqref{eq:FKx1},
\begin{equation}
- \langle \widehat \Delta_{x_j}\chi_jv, \chi_jv \rangle_{L^2(G_{x_j},\widehat\mu_{x_j})} \geq c^\FK_{x_j} \ \widehat\mu_{x_j}(\supp (\chi_j v))^{-\sfrac2Q} \,\|\chi_j v \|_{L^2(G_{x_j},\widehat\mu_{x_j})}^2.
\end{equation}
By the Hölder continuity of $x\mto c^\FK_x$ (Lemma \ref{faberkrahn}) we have
\[
c_{x_j}^\FK \geq (1-C_3 \epsilon^t) \,c_x^\FK
\quad\text{for all} \ x\in B(x_j,\epsilon)
\]
with a certain exponent $t>0$. This, together with the smoothness of the measure, implies that
\begin{align*}
c^\FK_{x_j} \ \widehat\mu_{x_j}(\supp (\chi_j v))^{-\sfrac2Q}
& \geq (1-C_4 \epsilon^t) \biggl(\int_{\supp(\chi_j v)} \left(c^\FK_x\right)^{-\sfrac Q2} d\mu(x) \biggr)^{-\sfrac 2Q} \\
& \geq (1-C_4 \epsilon^t) \biggl(\int_\Omega \left(c^\FK_x\right)^{-\sfrac Q2} d\mu(x) \biggr)^{-\sfrac 2Q}.
\end{align*}
Moreover, again by the smoothness of the measure,
\[
\|\chi_j v \|_{L^2(G_{x_j},\widehat\mu_{x_j})}^2 \geq (1-C_5 \varepsilon) \|\chi_j v \|_{L^2(M)}^2.
\]
Thus, we~have proved that
\begin{multline*}
\langle -\Delta^{M,\mu}_\Xb (\chi_j v),\chi_j v \rangle_{L^2(M)}\geq (1-C_2 \epsilon^s)(1-C_4 \epsilon^t) (1-C_5\epsilon) \\
\times
\left(\int_\Omega \left(c^\FK_x\right)^{-\sfrac Q2} d\mu(x)\right)^{-\sfrac 2Q}
\|\chi_j v \|_{L^2(M)}^2.
\end{multline*}

Summing over $j$, and inserting the resulting bound together with \eqref{eq:remainder1} into \eqref{e:identity}, we~obtain
\begin{multline*}
- \langle \Delta v, v \rangle_{L^2(M)}\geq \Bigl((1-C_2 \epsilon^s)(1-C_4 \epsilon^t) (1-C_5\epsilon) - C_1 \epsilon^{-2} \bigl(\inf_{x\in M} c_x^\FK \bigr)^{-1} \eta^{\sfrac2Q} \Bigr) \\
\times \left(\int_\Omega \left(c^\FK_x\right)^{-\sfrac Q2} d\mu(x)\right)^{-\sfrac 2Q} \| v \|_{L^2(M)}^2.
\end{multline*}
Now given $\theta>0$ we first choose $\epsilon\in(0,\epsilon_0]$ such that
\[
(1-C_2 \epsilon^s)(1-C_4 \epsilon^t) (1-C_5\epsilon) \geq 1- \frac{\theta}{2}
\]
and then $\eta>0$ such that
\[
C_1 \epsilon^{-2} \bigl(\inf_{x\in M} c_x^\FK \bigr)^{-1} \eta^{\sfrac2Q} \leq \frac{\theta}{2}.
\]
In this way we obtain the claimed inequality.
\end{proof}

\subsection{Proof of Theorem \ref{t:manifold}}
We are finally in position to give the proof of our main result for sub-Laplacians in the equiregular case.

\begin{proof}[Proof of Theorem \ref{t:manifold}]
Let $u_k$ be an eigenfunction associated with $\lambda_k$. Given \hbox{$\theta>0$}, let $\eta$ be as in Proposition \ref{faberkrahnbm} and consider any nodal domain $D_{k\ell}$ of~$u_k$ with \hbox{$\mu(D!_{k\ell})\leq\eta$}. We~denote by $\underline{u}_{k\ell}$ the restriction of $u_k$ to $D_{k\ell}$, extended by~$0$ out\-side~$D_{k\ell}$. Then, $\underline{u}_{k\ell}$ is not necessarily in the operator domain of $-\Delta$, but it is in the form domain, as shown in Theorem \ref{restriction}. Therefore it can be approximated with respect to the form-norm by $C^\infty_\rc(D_{k\ell})$ functions. Therefore we deduce from Proposition \ref{faberkrahnbm} that
\[
\langle -\Delta \underline{u}_{k\ell},\underline{u}_{k\ell}\rangle_{L^2(M)}
\geq (1-\theta) \left(\int_{D_{k\ell}} \left(c_x^\FK\right)^{-\sfrac Q2} d\mu(x)\right)^{-\sfrac 2Q} \| \underline{u}_{k\ell} \|_{L^2(M)}^2.
\]
Here we are slightly abusing notation by writing $\langle -\Delta \underline{u}_{k\ell},\underline{u}_{k\ell}\rangle_{L^2(M)}$ instead of the more precise $\| \tilde \nabla \underline{u}_{k\ell}\|^2_{L^2(M)}$. Similarly, using the weak formulation of the eigenvalue equation and the fact that $\underline{u}_{k\ell}=u_k$ on $D_{k\ell}$ we find
\[
\langle -\Delta \underline{u}_{k\ell},\underline{u}_{k\ell}\rangle_{L^2(M)} = \lambda_k \|\underline{u}_{k\ell}\|_{L^2(M)}^{2}.
\]
Combining the two previous equations and noting that $\underline{u_{k\ell}}$ does not vanish identically, we~obtain the inequality
\begin{equation}\label{e:almostthere}
\lambda_k \geq (1-\theta) \left(\int_{D_{k\ell}} \left(c_x^\FK\right)^{-\sfrac Q2} d\mu(x)\right)^{-\sfrac 2Q}.
\end{equation}
We denote by $\mathcal A_\eta$ the family of nodal sets $D_{k\ell}$ of $u_k$ satisfying $\mu(D_{k\ell})\leq \eta$. Raising \eqref{e:almostthere} to the power $Q/2$ and summing over $\ell$ we obtain
\[
\lambda_k^{\sfrac Q2} \int_M \left(c_x^\FK\right)^{-\sfrac Q2} d\mu(x)
\geq \lambda_k^{\sfrac Q2} \sum_{D_{k\ell}\in \mathcal{A}_\eta} \int_{D_{k\ell}} \left(c_x^\FK\right)^{-\sfrac Q2} d\mu(x) \geq (1-\theta)^{\sfrac Q2} \,(\#\mathcal{A}_\eta).
\]
Clearly, for the number of nodal sets of $u_k$ with $\mu$-measure exceeding $\eta$ we have
\[
\eta \#\{ \ell:\ D_{k\ell}\not\in \mathcal A_\eta \} \leq \sum_{D_{k\ell}\not\in \mathcal{A}_\eta} \mu(D_{k\ell}) \leq \mu(M).
\]
The two previous relations imply that
\[
\frac{\nu_k}{k} \leq (1-\theta)^{-\sfrac Q2} \left(\int_M \left(c_x^\FK\right)^{-\sfrac Q2} d\mu(x)\right) \frac{\lambda_k^{\sfrac{Q}{2}}}{k} + \eta^{-1} \mu(M) \frac{1}{k}.
\]
Combining this bound with the Weyl law from Theorem \ref{weylmet}, we~obtain
\[
\limsup_{k\to\infty} \frac{\nu_k}{k}\leq (1-\theta)^{-\sfrac Q2} \biggl(\int_M\bigl(c_x^\FK\bigr)^{-\sfrac Q2} d\mu(x) \biggr) \biggl(\int_M c_x^\Weyl \,d\mu(x) \biggr)^{-1}.
\]
Since $\theta>0$ is arbitrary, we~obtain the bound claimed in Theorem \ref{t:manifold}.
\end{proof}

\section{Basic examples}\label{AppA}

In this section we give examples of the applicability of Theorem \ref{t:manifold}. More precisely, we~will use the version from Remark \ref{remdom}, which concerns the result on an open subset with Dirichlet boundary conditions. We~will present the same computation in various forms which, we~hope, illustrates the different techniques that can be applied for a concrete operator.

\subsection{Set-up of the example}
We denote coordinates on $\R^3$ by $(x,y,z)$. In~an open subset $\Omega\subset\R^3$, we~consider the vector fields
\[
X_1= \frac{\partial}{\partial x} + K_1(x,y) \frac{\partial}{\partial z},
\quad
X_2= \frac{\partial}{\partial y} + K_2 (x,y) \frac{\partial}{\partial z},
\]
under the assumption that
\[
\curl \vec K = \frac{\partial K_2}{\partial x} - \frac{\partial K_1}{\partial y} >0
\quad\text{in}\ \Omega.
\]
In view of
\[
[X_1,X_2] = \curl \vec K \,\frac{\partial}{\partial z},
\]
we see that Assumption \ref{CHOr} is satisfied with $r=2$ and Assumption \ref{equireg} is satisfied with $n_1=2$ and $n_2=3$.

As measure $\mu$ we take the Lebesgue measure $dx\,dy\,dz$ restricted to $\Omega$. Then $X_j^\star = - X_j$ for $j=1,2$ and
\[
-\Delta_\bX^\Omega = - X_1^2- X_2^2.
\]
We consider this operator with Dirichlet boundary conditions on $\partial\Omega$.

We investigate the validity of Pleijel's theorem in this example. Our conclusion will be that Pleijel's theorem is valid in the example, provided it is valid on the Heisenberg group $\mathbb H$. Remarkably, this is independent of the choice of $\vec K$. While currently we do not know whether Pleijel's theorem is valid on the Heisenberg group, we~will show in the second part of this paper that it is valid assuming Pansu's conjecture about the isoperimetric inequality on $\mathbb H=\mathbb H_1$. Also, we~will discuss some modifications of the above example where we are able to prove unconditionally the validity of Pleijel's theorem.

In the remainder of this subsection we will describe the nilpotent approximation in this concrete example. The following two subsections are devoted to the computational details on how to perform the approximation. The reader interested in the conclusions may take the computations for granted and jump directly to the conclusions in Subsection~\ref{sec:exampleconclusion}.

Let $(x_0,y_0,z_0)\in\Omega$. Privileged coordinates $(u_1,u_2,u_3)$ at $(x_0,y_0,z_0)$ are given (modulo higher order terms if we use the canonical privileged coordinates\footnote{We refer to \cite{ChoPo1,ChoPo2} for a complete description of possible privileged coordinates}) by
\begin{equation} \label{eq:cdv}
u_1= x-x_0,\quad u_2 = y-y_0,\quad u_3 = \frac{1}{\curl \vec{K} (x_0,y_0)} \,(z-z_0) + P (x-x_0,y-y_0),
\end{equation}
where $P$ is a polynomial of order $2$. With respect to these coordinates the nilpotentizations of the vector fields $X_1,X_2$ at $(x_0,y_0,z_0)$ are
\begin{equation}\label{eq:chapeau}
\widehat X_{1,(x_0,y_0,z_0)} = \frac{\partial}{\partial u_1} + \frac 12 u_2 \frac{\partial}{\partial u_3},\quad
\widehat X_{2,(x_0,y_0,z_0)} = \frac{\partial}{\partial u_2} - \frac 12 u_1 \frac{\partial}{\partial u_3}
\end{equation}
and the nilpotentization of the Lebesgue measure at $(x_0,y_0,z_0)$ is
\begin{equation}\label{eq:exmeas}
d\widehat\mu_{(x_0,y_0,z_0)} = \curl \vec{K} (x_0,y_0)\,du_1 \,du_2 \,du_3.
\end{equation}
In particular, we~see that the nilpotent Lie group $G_{x_0,y_0,z_0}$ is the Heisenberg group~$\mathbb H$.

\subsection{Nilpotent approximation via a direct change of variables}
In this subsection we will describe one possible way of arriving at formulas \eqref{eq:chapeau} and \eqref{eq:exmeas}. While we have introduced the nilpotent approximation in Section~\ref{s1} in terms of canonical privileged coordinates of the first kind, in this subsection we will work with a different system of coordinates. The point we want to convey is that for practical computations it is often preferable to use coordinates that are similar, but different from the canonical privileged coordinates of the first kind. The fact that this still gives that ``correct'' nilpotentization will be discussed at the end of this subsection. For comparison, in the next subsection, we~perform the nilpotent approximation in terms of canonical privileged coordinates of the first kind.

After a translation, we~may assume that the point $(x_0,y_0,z_0)\in\Omega$ at which we are performing the nilpotent approximation is the origin, that is, we~assume $(x_0,y_0,z_0)= (0,0,0) \in \Omega$.

As will be justified later, since our example has step 2 (that is, only brackets of length $\leq 2$ are relevant), to compute the nilpotent approximation we may replace the functions $K_1$ and $K_2$ in the vector fields $X_1$ and $X_2$ by their linear approximation around $(0,0,0)\in\Omega$. That is, we~will consider
\[
X_1^\lin := \frac{\partial}{\partial x} + (K_1(0,0)) + \alpha x + \beta y) \frac{\partial}{\partial z},\quad
X_2^\lin := \frac{\partial}{\partial y} +(K_2 (0,0) + \gamma x + \delta y) \frac{\partial}{\partial z}
\]
with
\[
\alpha := \frac{\partial K_1}{\partial x} (0,0),
\quad
\beta := \frac{\partial K_1}{\partial y} (0,0),
\quad
\gamma := \frac{\partial K_2}{\partial x} (0,0),
\quad
\delta := \frac{\partial K_2}{\partial y} (0,0).
\]
Moreover, we~set
\[
\widehat \delta := \curl \vec K (0,0)
\]

We will now change variables $(x,y,z)\mto (u_1,u_2,u_3)$ as in \eqref{eq:cdv} to bring the vector fields $X_1^\lin, X_2^\lin$ into the form \eqref{eq:chapeau}. For pedagogical reasons we proceed step by step, each change of variables having the form described above (except that the factor $\widehat\delta^{-1}$ appears only in the last step).
\begin{itemize}
\item We make a first change of variables to replace $K_1(0,0)$ and $K_2(0,0)$ by $0$.
For this, we~make a change of variables of the form
\[
\tilde x = x,\quad
\tilde y =y,\quad \tilde z = z - K_1(0,0) x - K_2(0,0) y,
\]
and get in the new coordinates
\[
X_1^\lin= \frac{\partial}{\partial \tilde x} + (\alpha \tilde x + \beta \tilde y) \frac{\partial}{\partial \tilde z},\quad
X_2^\lin= \frac{\partial}{\partial \tilde y} +(\gamma \tilde x + \delta \tilde y) \frac{\partial}{\partial \tilde z}.
\]
From now on, we~omit the tilde.
\item We make a second change of variables to replace $\alpha$ and $\delta$ by $0$.
For this, we~make a change of variables of the form
\[
\tilde x = x,\quad \tilde y =y,\quad \tilde z = z - \frac 12 \alpha x^2 - \frac 12 \delta y^2,
\]
and get in the new coordinates
\[
X_1^\lin= \frac{\partial}{\partial \tilde x} + \beta \tilde y \frac{\partial}{\partial \tilde z},\quad
X_2^\lin= \frac{\partial}{\partial \tilde y} + \gamma \tilde x \frac{\partial}{\partial \tilde z}.
\]
Again, from now on, we~omit the tilde.
\item We make a third change of variables to have $\beta^\new = - \gamma^\new$.
For this, we~make a change of variables of the form
\[
\tilde x = x,\quad \tilde y =y,\quad \tilde z = z - \frac{\beta+\gamma}2\,xy,
\]
and get in the new coordinates (recalling $\widehat\delta = \beta - \gamma$)
\[
X_1^\lin= \frac{\partial}{\partial \tilde x} + \frac{\widehat\delta}{2}\,\tilde y\,\frac{\partial}{\partial \tilde z},\quad
X_2^\lin= \frac{\partial}{\partial \tilde y} - \frac{\widehat\delta}{2}\,\tilde x\,\frac{\partial}{\partial \tilde z}.
\]
Again, from now on, we~omit the tilde.
Note that all three changes of variables until now have respected the Lebesgue measure.
\item We make a fourth and final change of variables to remove the dependence on $\widehat\delta$ from the vector fields.
For this, we~make a change of variables of the form
\[
\tilde x = x,\quad \tilde y =y,\quad \tilde z = \widehat \delta^{-1} z,
\]
and get in the new coordinates
\[
X_1^\lin= \frac{\partial}{\partial \tilde x} + \frac{1}{2}\,\tilde y\,\frac{\partial}{\partial \tilde z},\quad
X_2^\lin= \frac{\partial}{\partial \tilde y} - \frac{1}{2}\,\tilde x\,\frac{\partial}{\partial \tilde z}.
\]
This is the claimed form \eqref{eq:chapeau} of the vector fields.
This time, the change of variables does not respect the Lebesgue measure. It~transforms Lebesgue measure $dx\,dy\,dz$ into the measure $\widehat\delta\,d\tilde x\,d\tilde y\,d\tilde z$, which is the claimed form \eqref{eq:exmeas} of the nilpotentized measure.
\end{itemize}

We end this subsection by addressing two issues that are left open by the above arguments. First, we~used the vector fields $X_j^\lin$ instead of $X_j$ and, second, the coordinates that we constructed are not the canonical privileged coordinates of the first kind.

To address the first issue, we~make the same change of variables as above, but with the original vector fields $X_1,X_2$ instead of $X_1^\lin,X_2^\lin$. We~find that in the new coordinates $(u_1,u_2,u_3)$ we have
\begin{align*}
X_{1,(0,0,0)} & = \widehat X_{1,(0,0,0)} + \widehat\delta^{-1} \,r_1(u_1,u_2) \,\frac{\partial}{\partial u_3},\\
X_{2,(0,0,0)} & = \widehat X_{2,(0,0,0)} + \widehat\delta^{-1} \,r_2(u_1,u_2) \,\frac{\partial}{\partial u_3},
\end{align*}
Here $\widehat X_{j,(0,0,0)}$ are the vector fields from \eqref{eq:chapeau} and the functions $r_1,r_2$ are defined by $K_j(x,y)= K_j(0,0) + \nabla K_j(0,0) \cdot (x, y) + r_{j} (x,y)$. It~follows that
\[
r_{j}(u_1,u_2) = \mathcal O (u_1^2 + u_2^2).
\]

This brings us to the second issue, which can be resolved in two ways. The first way is to note that with our change of variables (although it is not to canonical privileged coordinate of the first kind) we have obtained the conclusion of Lemma \ref{lemmaMet}. Noting that our proof of Theorem \ref{t:manifold}, via the reference to \cite{Rot}, only relied on the conclusion of this lemma, we~see that Theorem \ref{t:manifold} remains valid with the vector fields in \eqref{eq:chapeau} and the measure in \eqref{eq:exmeas}.

A second way to resolve this issue is to appeal to a general result saying that the nilpotent approximation can be performed in any system of so-called privileged coordinates and that the resulting nilpotent approximations are sub-Riemannian isometric. We~refer to \cite[\S A.5.4]{CHT1} for a discussion of these notions and assertions.

\Subsection{Computation of the canonical privileged coordinates and application}
In this subsection we perform the nilpotent approximation according to its definition in Section~\ref{s1} in terms of the canonical privileged coordinates of the first kind.

Again, after a translation we may assume that $(x_0,y_0,z_0)= (0,0,0)\in\Omega$.
We start from the vector fields
\[
Y_1:=X_1,\quad
Y_2:=X_2,\quad
Y_3:=[X_1,X_2] = \curl \vec K\,\frac{\partial}{\partial z}.
\]
which are adapted to the flag at $(0,0,0)\in\Omega$. According to \eqref{eq:pc} the canonical privileged coordinates of the first kind, $(u_1,u_2,u_3)$, are given by
\begin{equation}\label{eq:pc3}
(x,y,z)= \exp (\sum_{i=1}^3 u_{i}Y_i) \cdot (0,0,0).
\end{equation}
We observe that
\[
\sum_{i=1}^3 u_{i}Y_i = u_1 \frac{\partial}{\partial x} + u_2 \frac{\partial}{\partial y} +(\curl \vec K\,u_3 + K_1 u_1 + K_2 u_2) \frac{\partial}{\partial z}.
\]
To compute the right hand side of \eqref{eq:pc3} we have to compute the solution at $t=1$ of the differential system
\begin{align*}
\frac{dx}{dt} &=u_1,\\
\frac{dy}{dt} &=u_2, \\
\frac{dz}{dt} &= \curl \vec K (x,y)\,u_3 + K_1(x,y) u_1 + K_2 (x,y)u_2.
\end{align*}
with initial condition $(x(0),y(0),z(0))=(0,0,0)$.
This is easy to solve. We~have $x(t)= tu_1$, $y(t) =t u_2$ and
\[
z(t) = \int_0^t \bigl(\curl \vec K (su_1,su_2)\,u_3 + K_1(su_1,su_2) u_1 + K_2 (su_1,su_2)u_2\bigr)\,ds.
\]
Hence we get
\begin{align*}
(x,y,z) & =\exp (\sum u_{i}Y_i) \cdot (0,0,0)\\
&= \biggl(u_1,u_2, \int_0^1\bigl(\curl \vec K (su_1,su_2)\,u_3 + K_1(su_1,su_2) u_1 + K_2 (su_1,su_2)u_2\bigr) \,ds\biggr).
\end{align*}
The map $(u_1,u_2,u_3) \mto (x,y,z)$ is a local diffeomorphism since its differential at $(u_1,u_2,u_3)=(0,0,0)$ is
\[\arraycolsep3.5pt
\begin{pmatrix}
1&0&0\\
0&1&0\\
K_1(0,0)&K_2(0,0)&\curl \vec K (0,0)
\end{pmatrix}
\]
with determinant $\curl \vec K (0,0)>0$.

We note (for comparison with the first approach) that
\begin{align*}
\int_0^1\bigl(\curl \vec K &(su_1,su_2)\,u_3 + K_1(su_1,su_2) u_1 + K_2 (su_1,su_2)u_2\bigr)\,ds \\
& = \curl \vec K (0,0) u_3 + K_1(0,0) u_1+ K_2(0,0) u_2\\
& \quad
+ \frac12 \Big[ \frac{\partial K_1}{\partial x}(0,0) u_1^2+ \frac{\partial K_1}{\partial y}(0,0) u_1u_2 + \frac{\partial K_2}{\partial x}(0,0) u_1u_2 + \frac{\partial K_2}{\partial y}(0,0) u_2^2\Big]\\
& \quad + r(u_1,u_2,u_3)
\end{align*}
with $r$ of degree $\leq -3$ for the dilation.

Note that if we neglect the remainder $r$, this change of variables is exactly the one in the previous subsection. Indeed, for the inverse of the map $(u_1,u_2,u_3)\mto (x,y,z)$ we get $u_1=x$, $u_2=y$ and
\begin{multline*}
u_3 = \frac{1}{\curl \vec K (0,0)}\\
\times\biggl(z - \frac12 \Bigl(\frac{\partial K_1}{\partial x}(0,0) x^2+ \Bigl(\frac{\partial K_1}{\partial y}(0,0) + \frac{\partial K_2}{\partial x}(0,0) \Bigr) xy
+ \frac{\partial K_2}{\partial y}(0,0) y^2\Bigr) \biggr) \\
+ \tilde r(x,y,z).
\end{multline*}

With this formula for the canonical privileged coordinates at hand, the computations are essentially the same as in the previous subsection and we arrive again at the formulas \eqref{eq:chapeau} and \eqref{eq:exmeas}. We~omit the details.

\subsection{Conclusion in the setting of the example}\label{sec:exampleconclusion}
In order to apply Theorem \ref{t:manifold} we need to compute the local Weyl constant $c^\Weyl_{(x_0,y_0,z_0)}$, as well as the local Faber--Krahn constant $c_{(x_0,y_0,z_0)}^\FK$. We~will express them in terms of the corresponding constants on the Heisenberg group $\mathbb H = \mathbb H_1$, which will be studied in more detail in Part \ref{part2} of this paper.

There is, however, a slight notational inconsistency coming from different normalizations. In~the present section we arrived at the vector fields \eqref{eq:chapeau}, whereas in Part~\ref{part2} we will find it more convenient to work in the formulation \eqref{Otherformulation}. These two formulations are equivalent via scaling. More precisely, if~$\mathcal W(\mathbb H_1)$ and $C^\FK(\mathbb H_1)$ denote the Weyl and Faber--Krahn constants in the normalization of \eqref{Otherformulation} and if
$\widetilde{\mathcal W}(\mathbb H_1)$ and $\widetilde C^\FK(\mathbb H_1)$ denote the corresponding constants in the normalization of \eqref{eq:chapeau}, then
\[
\widetilde{\mathcal W}(\mathbb H_1) = 4\,\mathcal W(\mathbb H_1)
\quad\text{and}\quad
\widetilde C^\FK(\mathbb H_1) = 2^{-1}\,C^\FK(\mathbb H_1).
\]
As a consequence, the combination
\begin{equation}
\label{eq:pleijelheisenexample}
\bigl(\widetilde C^\FK(\mathbb H_1)\bigr)^{-2} \bigl(\widetilde{\mathcal W}(\mathbb H_1)\bigr)^{-1} = \left(C^\FK(\mathbb H_1)\right)^{-2} \left(\mathcal W(\mathbb H_1)\right)^{-1} = \gamma(\mathbb H_1)
\end{equation}
is independent of the normalization. The right equality in \eqref{eq:pleijelheisenexample} is a definition; see~\eqref{eq:defgamma}. The number $\gamma(\mathbb H_1)$ plays the role of the Pleijel constant on $\mathbb H_1$.

\begin{lemma}\label{exampleconst}
In our example we have for every $(x_0,y_0,z_0)\in\Omega$,
\[
c^\Weyl_{(x_0,y_0,z_0)} = \frac{\widetilde{\mathcal W}(\mathbb H_1)}{\mathrm{curl}\,\vec{K}(x_0,y_0)}
\quad\text{and}\quad
c_{(x_0,y_0,z_0)}^\FK = \sqrt{\mathrm{curl}\,\vec{K}(x_0,y_0)} \ \widetilde C^\FK(\mathbb H_1).
\]
\end{lemma}

\begin{proof}
The assertions follow by a simple scaling argument. According to \eqref{eq:chapeau} the differential expression for the nilpotent approximation $\widehat\Delta_{(x_0,y_0,z_0)}$ coincides with the Laplacian on the Heisenberg group, the only difference is that the measure $\widehat\mu_{(x_0,y_0,z_0)}$ is Lebesgue measure multiplied by the positive constant $\mathrm{curl}\,\vec{K}(x_0,y_0)$; see \eqref{eq:exmeas}. This factor leads to the corresponding expressions for the local Weyl and Faber--Krahn constants.
\end{proof}

It follows from Lemma \ref{exampleconst} that
\[
\int_\Omega c^\Weyl_{(x,y,z)} \,dx\,dy\,dz = \widetilde{\mathcal W}(\mathbb H_1) \ \int_\Omega \frac{dx\,dy\,dz}{\text{curl}\,\vec{K}(x,y)}
\]
and
\[
\int_\Omega (c_{(x,y,z)}^\FK)^{-2} \,dx\,dy\,dz = \widetilde C^\FK(\mathbb H_1) \,\int_\Omega \frac{dx\,dy\,dz}{\text{curl}\,\vec{K}(x,y)}.
\]
Therefore, from Theorem \ref{t:manifold} (or, more precisely, its version for the sub-Laplacian on an open set with Dirichlet boundary conditions in Remark \ref{remdom}) we deduce the bound
\[
\limsup_{k\to+\infty} \frac{\nu_k}{k} \leq \bigl(\widetilde C^\FK(\mathbb H_1)\bigr)^{-2} \bigl(\widetilde{\mathcal W}(\mathbb H_1)\bigr)^{-1} = \gamma(\mathbb H_1).
\]
Here we used \eqref{eq:pleijelheisenexample}. Our conclusion is that \emph{if the Pleijel constant $\gamma(\mathbb H_1)$ on the Heisenberg group is $<1$, then Pleijel's theorem holds in the present example.}

We will investigate the validity of the inequality $\gamma(\mathbb H_1)<1$ in Part \ref{part2}. Currently we have no proof of this bound, but we will prove that it holds provided a well-known conjecture by Pansu concerning the isoperimetric inequality on the Heisenberg group is true.

\subsection{Variations of the example}
In this subsection we present two variations of the above example. The point of these modified examples is that there we can prove the validity of Pleijel's theorem.

In the first variation we consider the same vector fields $X_1,X_2$ as before in this section and we consider the operator
\[
-X_1^2 - X_2^2 -\Delta_w
\]
on $\Omega \subset \R^3 \times \mathbb R^k$ with $k\in\N$. Here $\Delta_w$ denotes the Laplacian with respect to the variable $w\in\R^k$.

Repeating the above analysis, we~arrive at a certain constant $\gamma (\mathbb H \times \mathbb R^k)$ that is the analogue of the constant $\gamma (\mathbb H) = \bigl(\widetilde C^\FK(\mathbb H_1)\bigr)^{-2} \bigl(\widetilde W (\mathbb H_1)\bigr)^{-1}$. As~will be discussed in Part \ref{part2}, the assumption $k\geq 3$ guarantees that $\gamma (\mathbb H \times \mathbb R^k)<1$; see Theorem \ref{maincomp} together with \eqref{eq:defgamma}. Therefore Pleijel's theorem holds in this example for all $\vec K$ with $\curl\vec K>0$. We~do not carry out the details.

In the second variation, in an open subset of $\R^n\times\R^n\times\R$, $n\in\mathbb N$, with coordinates $(\vec x,\vec y,z)$, $\vec x =(x_1,\dots,x_n)$ and $\vec y=(y_1,\cdots, y_n)$, we~consider the following vector fields for $j=1,\ldots, n$:
\begin{align*}
X'_j&= \partial_{x_j} + K_1^j (x_j,y_j) \partial_z,\\
X''_j& = \partial_{y_j} + K_2^j (x_j,y_j) \partial_z.
\end{align*}
We assume that $\curl \vec K{}^j >0$ for $j=1,\ldots, n$.
For the nilpotentization at $(\vec x,\vec y,z)$ in coordinates $(\vec u,\vec v,w)$ this leads to the following vector fields on $\mathbb H_n$,
\begin{align*}
\hat X'_j & = \curl\vec K{}^j(x_j,y_j)^{1/2} \ (\partial _{u_j} +\tfrac 12 v_j \partial_w),\\
\hat X''_j& = \curl\vec K{}^j(x_j,y_j)^{1/2} \ (\partial _{v_j} - \tfrac 12 u_j \partial_w),
\end{align*}
together with the Haar measure
\[
\biggl(\prod_{j=1}^n \curl \vec K{}^j(x_j,y_j) \biggr)^{-1} \,du\,dv\,dw.
\]
For the Faber--Krahn part, one has to consider on $\mathbb H_n$
\[
\widehat \Delta_{\vec x,\vec y,z}= \sum_j \curl \vec K{}^j(x_j,y_j) \bigl((\partial _{u_j} + \tfrac{ 1}{2} v_j \partial_w)^2+(\partial _{v_j} - \tfrac {1}{2}u_j\partial_w)^2 \bigr).
\]

Assuming that there is a constant $\kappa>0$ such that $\curl \vec K{}^j \equiv \kappa$ for $j=1,\ldots,n$, it~will be clear from our analysis in Part \ref{part2} (Theorem \ref{maincomp}) that Pleijel's theorem holds for $n\geq 4$ in any open set of $\Omega$. Since the Weyl and Faber--Krahn constants depend continuously on the numbers $\curl \vec K{}^j(x_j,y_j)$, we~see that for any $n\geq 4$ there is a $\delta>0$ such that Pleijel's theorem holds provided that
\[
\biggl| \frac{\curl \vec K{}^j(x_j,y_j)}{\curl \vec K{}^k(x_k,y_k)} - 1\biggr| \leq \delta
\quad\text{for all}\ j, k, x_j,y_j, x_k,y_k.
\]

\part{Pleijel's bound for \texorpdfstring{$\mathbb H_n\times\R^k$}{HnRk}}\label{part2}

\section{The Pleijel argument for \texorpdfstring{$\mathbb H_n\times\R^k$}{HnRk} }\label{sec:mainhnrk}

We work on $\mathbb H_n\times\R^k$, where $n\in\N$, $k\in\N_0$. The case $k=0$ corresponds to the Heisenberg group $\mathbb H_n$. (Everything remains valid for $n=0$ as well, that is, for $\R^k$, but in this case the results below are well-known.) Typically, we~will denote coordinates in $\mathbb H_n$ by $(x,y,z)$ with $x,y\in\R^n$ and $z\in\R$, and we will denote coordinates in $\R^k$ by $w$. The measure $dx\,dy\,dz\,dw$ is the Lebesgue measure on $\R^{2n+1+k}$. For the vector fields we will use the following normalization,\footnote{In this and the remaining sections it is more convenient to use another normalization than in the first sections. This simply corresponds to a scaling of the $z$-variable.}
\begin{subequations}\label{Otherformulation}
\begin{equation}
X_j = \partial_{x_j} + 2y_j \partial_z,\quad
Y_j = \partial_{y_j} - 2x_j \partial_z,\quad
W_j = \partial_{w_j}.
\end{equation}
The sub-Laplacian is
\begin{equation}
\Delta^{\mathbb H_n\times\R^k} = \sum_{j=1}^n (X_j^2 + Y_j^2) + \sum_{i=1}^k W_i^2.
\end{equation}
\end{subequations}

If $\Omega\subset\mathbb H_n\times\R^k$ is an open set of finite measure, then the spectrum of the Dirichlet realization of $ -\Delta^{\mathbb H_n\times\R^k}_\Omega $ is discrete and we can denote its eigenvalues, in nondecreasing order and repeated according to multiplicities, by~$\lambda_\ell(\Omega)$, $\ell\in\N$. We~know that eigenfunctions are $C^\infty$ in $\Omega$ (\cite{Ho}) and therefore the nodal domains are well defined as the connected components of the complement of their zero set in $\Omega$. We~denote by $\nu_\ell(\Omega)$ the maximal number of nodal domains of eigenfunctions corresponding to eigenvalue $\lambda_\ell(\Omega)$. We~are interested in an upper bound on
\[
\limsup_{\ell\to\infty} \frac{\nu_\ell(\Omega)}{\ell}
\]
that depends only on $n$ and $k$. The simplest such upper bound is given by one, as~mentioned in Remark \ref{remcourantasymp}.

In the spirit of Pleijel's theorem, here we try to improve upon the upper bound by one. Just as Pleijel's bound, our bound depends on two constants, namely the constant in the Weyl asymptotics and the Faber--Krahn constant. Let us introduce these constants. The Weyl asymptotics in the case with boundary (which was also established by G.\,Métivier in \cite[Th.\,1.3]{Me}) states
that, for any open set $\Omega\subset\mathbb H_n\times\R^k$ of finite measure,
\begin{equation}\label{Metwithbdry}
\lambda^{-\psfrac{2n+2+k}{2}} \#\{ \ell:\ \lambda_\ell(\Omega)<\lambda \} \to \mathcal W(\mathbb H_n\times\R^k) \,|\Omega|
\quad\text{as}\ \lambda\to\infty.
\end{equation}
We will give a (relatively) explicit expression for the constant $\mathcal W(\mathbb H_n\times\R^k)$ in the next section.

The Faber--Krahn constant $C^\FK(\mathbb H_n\times\R^k)$ is defined to be the largest constant such that for any open $\Omega\subset\mathbb H_n\times\R^k$ of finite measure and for any $u\in S^1_0(\Omega)$ one has
\begin{multline}
\label{eq:fk}
\int_\Omega \biggl(\sum_{j=1}^n \bigl((X_ju)^2 + (Y_ju)^2\bigr) + \sum_{i=1}^k (W_iu)^2\biggr)\,dx\,dy\,dz\,dw \\
\geq C^\FK(\mathbb H_n\times\R^k) |\Omega|^{-\spfrac{2}{2n+2+k}} \int_\Omega u^2 \,dx\,dy\,dz\,dw.
\end{multline}
Here $S^1_0(\Omega)$ (see Subsection \ref{ss2.2}) denotes the form domain of the Dirichlet realization of~$-\Delta$ on $\Omega$ or, equivalently, the completion of $C^1_\rc(\Omega)$ with respect to the quadratic for
\[u \mto \int_\Omega \biggl(\sum_{j=1}^n \bigl((X_ju)^2 + (Y_ju)^2 \bigr) + \sum_{i=1}^k (W_iu \biggr)^2 + u^2)\,dx\,dy\,dz\,dw.
\]
The defining inequality for the Faber--Krahn constant can also be stated as
\[
\lambda_1(\Omega) \geq C^\FK(\mathbb H_n\times\R^k) \,|\Omega|^{-\spfrac{2}{2n+2+k}}
\]
for all open $\Omega\subset\mathbb H_n\times\R^k$ of finite measure.

Let us set
\begin{equation}\label{eq:defgamma}
\gamma(\mathbb H_n\times\R^k) := \left(C^\FK(\mathbb H_n\times\R^k)\right)^{-\psfrac{2n+2+k}{2}} \left(\mathcal W(\mathbb H_n\times\R^k)\right)^{-1}.
\end{equation}
Here is our Pleijel-type bound.

\begin{theorem}\label{pleijelnk}
For any open $\Omega\subset\mathbb H_n\times\R^k$ of finite measure,
\[
\limsup_{\ell\to\infty} \frac{\nu_\ell(\Omega)}{\ell} \leq \gamma(\mathbb H_n\times\R^k).
\]
\end{theorem}

\begin{proof}[Proof of Theorem \ref{pleijelnk}]
We consider an eigenfunction $u$ corresponding to the eigenvalue $\lambda_\ell(\Omega)$. Let $(\omega_\alpha)_\alpha$ be its nodal domains and let $\nu_\ell(u)$ be their number. (We will see shortly that this number is finite.) By Theorem \ref{restriction} with $M=\mathbb H_n\times\R^k$,
we know that $\lambda_\ell(\Omega) = \lambda_1(\omega_\alpha)$ and that $u|_{\omega_\alpha}$ is the ground state of the Dirichlet realization on~$\omega_\alpha$. Thus,
\begin{align*}
\frac{\nu_\ell(u)}{\ell} & = \frac{\lambda_\ell(\Omega)^{\psfrac{2n+2+k}{2}}}{\ell} \sum_\alpha \lambda_1(\omega_\alpha)^{-\psfrac{2n+2+k}{2}} \\
& \leq \frac{\lambda_\ell(\Omega)^{\psfrac{2n+2+k}{2}}}{\ell} \left(C^\FK(\mathbb H_n\times\R^k)\right)^{-\spfrac{2}{2n+2+k}} \sum_\alpha |\omega_\alpha| \\
& \leq \frac{\lambda_\ell(\Omega)^{\psfrac{2n+2+k}{2}}}{\ell} \left(C^\FK(\mathbb H_n\times\R^k)\right)^{-\spfrac{2}{2n+2+k}} |\Omega|.
\end{align*}
Since this is true for any eigenfunction corresponding to $\lambda_\ell(\Omega)$, we~deduce that
\[
\frac{\nu_\ell(\Omega)}{\ell} \leq \frac{\lambda_\ell(\Omega)^{\psfrac{2n+2+k}{2}}}{\ell} \left(C^\FK(\mathbb H_n\times\R^k)\right)^{-\spfrac{2}{2n+2+k}} |\Omega|.
\]
Taking the limsup as $\ell\to\infty$ and with in mind the Weyl asymptotics, we~arrive at the claimed bound.
\end{proof}

We recall that by Theorem \ref{courant} we have $\limsup_{\ell\to\infty} \sfrac{\nu_\ell(\Omega)}{\ell} \leq 1$; see also Remark~\ref{remcourantasymp}. In~the remaining sections of this paper we will give sufficient conditions on $n$ and $k$ to have $\gamma(\mathbb H_n\times\R^k)<1$. We~recall that for $n=0$ this was shown to be the case for $k=2$ by Pleijel \cite{Pl} and for general $k$ by Bérard and Meyer \cite{BerMe}. Moreover, Helffer and Persson Sundqvist \cite{HPS} showed that, for $n=0$, the sequence $k\mto\gamma(\R^k)$ is decreasing. Here we shall prove, among other things, the following.

\begin{theorem}\label{maincomp}
Let $n\geq 1$ and $k\geq 0$ with $(n,k)\not\in\{(1,0),(2,0),(3,0),(1,1)\}$. Then $\gamma(\mathbb H_n\times\R^k)<1$.
\end{theorem}

The proof of Theorem \ref{maincomp} is somewhat long and spread out over several sections. Here is a guide. The part concerning $k=0$ is proved in Subsection \ref{sec:maincompk0} and that concerning $n\geq 3$ in Subsection \ref{sec:maincompk}. The part concerning $n=1$ and $n=2$ is proved in Subsection \ref{sec:maincompn12}.

There is a well-known conjecture, due to Pansu \cite{Pa2}, about the sharp isoperimetric constant on the Heisenberg group. It~is generally believed to be true and supported by several partial results. We~will discuss this in some detail in Section~\ref{sec:isoperimetric}. We~shall show that the validity of this conjecture implies Pleijel's bound.

\begin{proposition}\label{pansupleijel}
Let $n\in\N$ and assume that Pansu's conjecture \eqref{eq:pansusconj} holds. Then $\gamma(\mathbb H_n\times\R^k)<1$ for all $k\in\N_0$.
\end{proposition}

We will prove this proposition for $k=0$ and $n=1,2,3$ in Corollary \ref{pansupleijelk0} and for $k=1$ and $n=1$ in Subsection \ref{sec:maincompn12}. In~the remaining cases Theorem \ref{maincomp} applies.

\section{Computing the constant in the Weyl asymptotics}\label{sec:weylhnrk}

\subsection{The case of $\mathbb H_n$}\label{sec:weylheis1}
As discussed in Theorem \ref{weylmet} and \eqref{Metwithbdry} for the case with boundary, the Weyl asymptotics for the Dirichlet realization of the sub-Laplacian in open subsets $\Omega$ of $\mathbb H_n$ state that
\[
\lambda^{-\psfrac{2n+2}{2}} \#\{ \ell:\ \lambda_\ell(\Omega)<\lambda \} \to \mathcal W(\mathbb H_n) \ |\Omega|
\quad\text{as}\ \lambda\to\infty.
\]
Hansson--Laptev \cite{HL} have shown that these asymptotics hold under the sole assumption that $\Omega$ is an open set of finite measure.

Since we are interested in a relatively explicit expression of the constant $\mathcal W(\mathbb H_n)$ and since Hansson and Laptev use a different normalization from ours, we~repeat part of their argument; see also \cite{St} for a related computation. They show that
\begin{equation}\label{w1}
N(\lambda,-\Delta_\Omega^{\mathbb H_n}) \sim \int_\Omega \1(-\Delta^{\mathbb H_n} <\lambda)((x,y,z),(x,y,z)) \,dx\,dy\,dz
\quad\text{as}\ \lambda\to\infty,
\end{equation}
where $\1(-\Delta^{\mathbb H_n} <\lambda)((x,y,z),(x,y,z))$ is the on-diagonal spectral density of the sub-Laplacian on all of $\mathbb H_n$. (To be more precise, using coherent states Hansson and Laptev show \eqref{w1} when integrated over $\lambda$. Then a Tauberian theorem yields \eqref{w1} as stated.)

By translation invariance of the sub-Laplacian on $\mathbb H^n$, we~know that the diagonal of the spectral density $\1(-\Delta^{\mathbb H_n} <\lambda)((x,y,z),(x,y,z))$ is independent of the point $(x,y,z)$. Moreover, by~dilation covariance, we~know that it is proportional to $\lambda^{\sfrac{Q}{2}}$. Thus, there is a constant $\mathcal W(\mathbb H_n)>0$ such that
\begin{equation}\label{w2}
\1(-\Delta^{\mathbb H_n} <\lambda)((x,y,z),(x,y,z)) = \mathcal W(\mathbb H_n)\ \lambda^{\sfrac Q2},
\end{equation}
and we obtain the above form of the spectral asymptotics.

In the following we are interested in finding an explicit expression for $\mathcal W(\mathbb H_n)$. We~proceed by an explicit diagonalization of the operator $-\Delta^{\mathbb H_n}$. By~a Fourier transform with respect to $z$ one arrives at the family of operators
\[
- \sum_{j=1}^n \bigl((\partial_{x_j} + 2iy_j \zeta)^2 + (\partial_{y_j} - 2i x_j \zeta)^2 \bigr),
\]
where $\zeta\in\R$ is the Fourier variable dual to the variable $z$.

Now for each fixed $j$,
\[
- \bigl((\partial_{x_j} + 2i y_j \zeta)^2 + (\partial_{y_j} - 2ix_j \zeta)^2 \bigr),
\]
is a Landau Hamiltonian corresponding to constant magnetic field with intensity $4|\zeta|$. These $n$ Hamiltonians are independent of each other. The spectrum of each one is given by $4|\zeta|(2k_j +1)$, $k_j\in\N_0 =\{0,1,2,\ldots\}$, and each such eigenvalue contributes $4|\zeta|/(2\pi)$ to the spectral density.

The spectral function of the operator $-\Delta$ on $\mathbb H_n$ is then given by
\begin{align*}
\1(-\Delta<\lambda)((x,y,z)&,(x,y,z)) \\
& = \int_\R \frac{d\zeta}{2\pi} \Bigl(\frac{4|\zeta|}{2\pi}\Bigr)^n \sum_{k \in\N_0^n} \1\bigl(4|\zeta|(2(k_1+\ldots+k_n)+n)<\lambda\bigr) \\
& = \frac{4^n}{(2\pi)^{n+1}} \sum_{k\in\N_0^n} \frac{2}{n+1} \Bigl(\frac{\lambda}{4(2(k_1+\ldots+k_n)+n)}\Bigr)^{n+1} \\
& = \frac{1}{2(n+1)} \frac{\lambda^{n+1}}{(2\pi)^{n+1}} \sum_{k\in\N_0^n} \Bigl(\frac{1}{2(k_1+\ldots+k_n)+n}\Bigr)^{n+1} \\
& = \frac{1}{2(n+1)} \frac{\lambda^{n+1}}{(2\pi)^{n+1}} \sum_{m\in\N_0} \binom{m+n-1}{m} \Bigl(\frac{1}{2m+n}\Bigr)^{n+1}.
\end{align*}
Thus, we~have shown that
\begin{equation}\label{w4}
\mathcal W(\mathbb H_n) = \frac{1}{2(n+1)} \frac{1}{(2\pi)^{n+1}} \sum_{m\in\N_0} \binom{m+n-1}{m} \frac{1}{(2m+n)^{n+1}}.
\end{equation}
Note that
\begin{equation}\label{eq:W1}
\mathcal W(\mathbb H) = \frac{1}{4} \frac{1}{(2\pi)^2} \sum_{m\in\N_0} \frac{1}{(2m+1)^2} = \frac14 \frac1{(2\pi)^2} \frac{\pi^2}8 = \frac1{128}
\end{equation}
and
\begin{equation}\label{eq:W2}
\mathcal W(\mathbb H_2) = \frac{1}{6} \frac{1}{(2\pi)^3} \sum_{m\in\N_0} \frac{m+1}{(2m+2)^3} = \frac{1}{6} \frac{1}{(2\pi)^3} \frac{\pi^2}{48} = \frac{1}{48^2 \pi}.
\end{equation}
It seems like Hansson--Laptev \cite{HL} were not aware of a more explicit form of $\mathcal W(\mathbb H_n)$ for $n\geq 3$. We~will give explicit formulas below for $3\leq n \leq 13$ and propose a general conjecture.

Let us set
\begin{equation}\label{eq:defcn}
c_n := \sum_{m\in\N_0} \binom{m+n-1}{m} \frac{1}{(2m+n)^{n+1}}.
\end{equation}
We have seen above that
\[
c_1=\frac{\pi^2}8 \quad \text{and}\quad
c_2 = \frac{\pi^2}{48}.
\]
Numerical values for the $c_n$, $n=3,4,5,6$, are given in the Hansson--Laptev paper \cite{HL} and can be completed by using Wolfram Alpha or Mathematica.\footnote{Thanks to J.\,Viola and F.\,Nicoleau for their help.}
This program shows first that explicit formulas can be found for $c_n$ and then gives numerical values, which are consequently quite accurate. We~get
{\smaller\smaller
\[
\begin{aligned}
c_3 & = \frac{\pi^2(12-\pi^2)}{768} &&\approx 2.7378\cdot10^{-2}\\
c_4 & = \frac{\pi^2 (15-\pi^2)}{17280}&&\approx 2.9303 \cdot  10^{-3}\\
c_5& = \frac{\pi^2 (120-100\pi^2 + 9 \pi^4)}{368640}&&\approx 2.6027\cdot  10^{-4}\\
c_6&= \frac{\pi^2(315-105 \pi^2+ 8 \pi^4)}{29030400}&&\approx 1.9706 \cdot  10^{-5}\\
c_7&=\frac{\pi^2 (6720 - 19600 \pi^2 + 14504 \pi^4 - 1275 \pi^6)}{2477260800}&&\approx 1.2988\cdot 10^{-6}\\
c_8 &= \frac{\pi^2 (1575 - 1470 \pi^2 + 490 \pi ^4 - 36 \pi ^6)}{24385536000}&&\approx 7.5736\cdot10^{-8}\\
c_9&=\frac{\pi^2(40 320 - 282 240 \pi^2 + 663 264 \pi^4 - 439 144 \pi^6 + 37975 \pi^8)}{3329438515200}&&\approx 3.9589\cdot  10^{-9}\\
c_{10}&= \frac{\pi^2 (3465 - 6930 \pi^2 + 6006 \pi ^4 - 1804 \pi^6 + 128 \pi ^8)}{15450675609600}&&\approx1.8749\cdot10^{-10}
\end{aligned}
\]
}%
Continuing with Mathematica, we~get for the quotients:
{\smaller\smaller
\begin{align*}
\frac{c_{11}}{c_{10}} & = \frac{3 \left(1774080 - 24393600 \pi ^2 + 129773952 \pi ^4 - 258523760 \pi ^6 + 160227716 \pi ^8 - 13712895 \pi ^{10}\right)}{10240
\left(3465-6930 \pi ^2+6006 \pi ^4-1804 \pi ^6+128 \pi ^8\right)}, \\
\frac{c_{12}}{c_{11}} & = \frac{256 \left(2837835-10405395 \pi ^2 + 18432414 \pi ^4 - 13774761 \pi ^6 + 3835832 \pi ^8 - 265344 \pi ^{10}\right)}{27027
(1774080 \!-\! 24393600 \pi ^2 \!+\! 129773952 \pi ^4 \!-\! 258523760 \pi ^6 \!+\! 160227716 \pi ^8 \!-\! 13712895 \pi ^{10})},\\
\frac{c_{13}}{c_{12}} & = \frac{\hat C_{13}}{40960 \left(2837835 - 10405395 \pi ^2 + 18432414 \pi ^4 - 13774761 \pi ^6 + 3835832 \pi ^8 - 265344 \pi ^{10}\right)},
\end{align*}
}%
with
\begin{multline*}
\hat C_{13}= 7 \bigl(2075673600-49470220800 \pi ^2+497175719040 \pi ^4-2161554183360 \pi ^6\\
+3895229400920 \pi ^8-2314322017956 \pi ^{10}+196697984175
\pi ^{12}\bigr).
\end{multline*}
This leads to
\begin{equation*}
c_{11} \approx 8.1149\,\cdot\,10^{-12}, \quad
c_{12} \approx 3.23369\,\cdot\,10^{-13},\quad
c_{13} \approx 1.1938 \,\cdot\,10^{-14}.
\end{equation*}

Although not important for our applications, it~is nice to see that this leads to the following guess:
\begin{quotation}
{\it
For any $n$, there is a polynomial $P_n$ of degree $[\frac {n+1} 2]$ with rational coefficients such that
\[
c_n= P_n(\pi^2),
\]
where $ [x]$ denotes the largest integer satisfying $[x]\leq x$.
}
\end{quotation}
This is related to formulas concerning multi-zeta functions, which are recognized by Mathematica and, as communicated to us by F.\,Nicoleau, to the Lerch function \cite[p.\,32]{MOS} and its derivatives.

\subsection{The Weyl constant on $\mathbb H_n\times\R^k$ with $k\geq 1$}
We show that the constant in the Weyl formula on $\mathbb H_n\times\R^k$ can be expressed in terms of that on $\mathbb H_n$.

\begin{lemma}\label{wnk}
For any $n,k\in\N$,\vspace*{-3pt}
\begin{equation}
\label{eq:wnk}
\mathcal W(\mathbb H_n\times\R^k) = \mathcal W(\mathbb H_n)\ (4\pi)^{-\sfrac{k}{2}} \frac{\Gamma(n+2)}{\Gamma(\frac{2n+k+4}{2})}.
\end{equation}
\end{lemma}

\begin{proof}
By the same argument as in the previous subsection we know that $\mathcal W(\mathbb H_n\times\nobreak\R^k)$ is given by\vspace*{-3pt}
\begin{equation*}
\1\bigl(-\Delta_{\mathbb H_n\times \mathbb R^k}<\lambda\bigr)((x,y,z,w),(x,y,z,w))=\mathcal W(\mathbb H_n\times\R^k)\ \lambda^{\psfrac{2n+2+k}{2}}.
\end{equation*}
Since\vspace*{-3pt}
\begin{multline*}
\1\bigl(-\Delta_{\mathbb H_n\times \mathbb R^k}<\lambda\bigr)((x,y,z,w),(x,y,z,w)) \\
= \int_\R \frac{d\zeta}{2\pi} \int_{\R^k} \frac{d\tau}{(2\pi)^k} \Bigl(\frac{4|\zeta|}{2\pi}\Bigr)^n \sum_{k \in\N_0^n} \1\bigl(4|\zeta|(2(k_1+\ldots+k_n)+n)+ \tau^2<\lambda\bigr)
\end{multline*}
(which follows by Fourier transforming $z\mto \zeta$ and $w\mto \tau$), we~find\vspace*{-3pt}
\begin{align*}\label{sp2}
\mathcal W(\mathbb H_n\times\R^k) & = (2\pi)^{ - k} |\Sph^{k-1}| \biggl(\int_0^1(1-\rho^2)^{n+1} \rho^{k-1} d\rho\biggr) \mathcal W(\mathbb H_n) \\
& = (2\pi)^{ - k} |\Sph^{k-1}| \biggl(2^{-1} \int_0^1(1-\sigma)^{n+1} \sigma^{\psfrac{k-2}2} d\sigma \biggr) \mathcal W(\mathbb H_n) \\
& = (2\pi)^{ - k} |\Sph^{k-1}| 2^{-1} \frac{\Gamma(n+2)\,\Gamma(\frac k2)}{\Gamma(\frac{2n+4+k}{2})} \ \mathcal W(\mathbb H_n).
\end{align*}
Here we expressed the beta function integral appearing on the second line just above in terms of gamma functions. Inserting $|\Sph^{k-1}| = \Gamma(\frac k2)^{-1} 2 \pi^{\sfrac k2}$, we~arrive at the claimed formula.
\end{proof}

\section{A first bound on the Faber--Krahn constant}\label{sec:fksob}

\subsection{A bound via the Sobolev constant}
We obtain a bound on the Faber--Krahn constant in terms of the (critical) Sobolev inequality on $\mathbb H_n\times\R^k$. By~definition, $C^\Sob(\mathbb H_n\times\R^k)$ is the largest constant such that for all $u\in S^1_0(\mathbb H_n\times\R^k)$\vspace*{-3pt}
\begin{multline*}
\int_{\mathbb H_n\times\R^k} \biggl(\sum_{j=1}^n \bigl((X_ju)^2 + (Y_ju)^2\bigr) + \sum_{i=1}^k (W_iu)^2\biggr) dx\,dy\,dz\,dw \\[-5pt]
\geq C^\Sob(\mathbb H_n\times\R^k) \biggl(\int_{\mathbb H_n\times\R^k} |u|^{\spfrac{2(2n+2+k)}{2n+k}} \,dx\,dy\,dz\,dw \biggr)^{\pspfrac{2n+k}{2n+2+k}}.
\end{multline*}

\begin{lemma}\label{fksob}
$
C^\FK(\mathbb H_n\times\R^k) \geq C^\Sob(\mathbb H_n\times\R^k).
$
\end{lemma}

\begin{proof}
If $\Omega\subset\mathbb H_n\times\R^k$ is open with finite measure and if $u\in S^1_0(\Omega)$, then, by~Hölder,
\begin{multline*}
\int_\Omega u^2 \,dx\,dy\,dz\,dw\\[-8pt]
\leq |\Omega|^{\spfrac{2}{2n+2+k}} \left(\int_{\mathbb H_n\times\R^k} |u|^{\spfrac{2(2n+2+k)}{2n+k}}\,dx\,dy\,dz\,dw\right)^{\pspfrac{2n+k}{2n+2+k}}.
\end{multline*}
Bounding the right side by Sobolev, we~obtain a Faber--Krahn-type inequality with constant $C^\Sob(\mathbb H_n\times\R^k)$. This implies the claimed bound for the optimal constants.
\end{proof}

An explicit expression for $C^\Sob(\mathbb H_n)$ was found by Jerison and Lee \cite{JL}; for an alternative proof see also \cite{FrLi}. We~have
\begin{equation}\label{eq:jl}
C^\Sob(\mathbb H_n) = \frac{4\pi n^2}{(2^{2n} n!)^\spfrac{1}{n+1}}.
\end{equation}

\subsection{The Pleijel constant $\gamma(\mathbb H_n)$}\label{sec:maincompk0}
Our goal in this subsection is to prove the part of Theorem \ref{maincomp} for $k=0$, that is, we~are going to prove that $\gamma(\mathbb H_n)<1$ for $n\geq 4$. To bound $\gamma(\mathbb H_n)$, defined in \eqref{eq:defgamma}, we~use Lemma \ref{fksob} and the explicit expression for $C^\Sob(\mathbb H_n)$ from \eqref{eq:jl} to get
\begin{equation}
\label{eq:gammatilden}
\gamma(\mathbb H_n) \leq \left(C^\Sob(\mathbb H_n)\right)^{-n-1} \mathcal W(\mathbb H_n)^{-1} = \frac{2^n(n+1)!}{n^{2(n+1)}} \frac{1}{c_n} =: \tilde\gamma_n,
\end{equation}
where $c_n$ is defined in \eqref{eq:defcn}. Inspired by \cite{HPS} we will consider the quotients $\tilde\gamma_n/\tilde\gamma_{n-1}$.

In view of \eqref{eq:gammatilden}, the part of Theorem \ref{maincomp} for $k=0$ is an immediate consequence of the following assertion.

\begin{proposition}\label{gammatildemono}
The sequence $n\mto \tilde\gamma_n$ is decreasing for $n\geq 1$. Moreover, $\tilde\gamma_n<1$ for $n\geq 4$.
\end{proposition}

Our proof of this proposition relies on numerical computations for $n\leq 13$. Using the values of $c_n$ from the previous section, we~get
\[
\begin{aligned}
\tilde\gamma_1 & \approx 3.242&
\tilde\gamma_2 & \approx 1.824\\
\tilde\gamma_3 & \approx 1.069&
\tilde\gamma_4 &\approx 6.251 \cdot 10^{-1} \\
\tilde\gamma_5 &\approx 3.628 \cdot 10^{-1} &
\tilde\gamma_6 & \approx  2.088 \cdot 10^{-1} \\
\tilde\gamma_7 & \approx 1.195\cdot 10^{-1} &
\tilde\gamma_8 & \approx  6.808\cdot 10^{-2} \\
\tilde\gamma_9 & \approx  3.860\cdot10^{-2} &
\tilde\gamma_{10} & \approx 2.180 \cdot 10^{-2} \\
\tilde\gamma_{11} &\approx 1.227\cdot 10^{-2} &
\tilde\gamma_{12} &\approx 6.891\cdot 10^{-3}  \\
\tilde\gamma_{13} & \approx 3.859\cdot 10^{-3}&\hspace*{1cm}
\end{aligned}
\]
It is also instructive to look at the quotients $\tilde\gamma_n/\tilde\gamma_{n-1}$, for which we get
\[
\begin{aligned}
\tilde\gamma_2/\tilde\gamma_1 & =  0.5625 &\hspace*{1cm}
\tilde\gamma_3/\tilde\gamma_2 & \approx 0.5861 \\
\tilde\gamma_4/\tilde\gamma_3 & \approx 0.5848 & \tilde\gamma_5/\tilde\gamma_4 & \approx 0.5804 \\
\tilde\gamma_6/\tilde\gamma_5 & \approx 0.5757 &
\tilde\gamma_7/\tilde\gamma_6& \approx 0.5721 \\
\tilde\gamma_8/\tilde\gamma_7& \approx0.5697&
\tilde\gamma_9/\tilde\gamma_8& \approx0.5670\\
\tilde\gamma_{10}/\tilde\gamma_9& \approx 0.5648&
\tilde\gamma_{11}/\tilde\gamma_{10}& \approx0.5630\\
\tilde\gamma_{12}/\tilde\gamma_{11}& \approx0.5614&
\tilde\gamma_{13}/\tilde\gamma_{12}& \approx0.5601
\end{aligned}
\]
These computations suggest that the sequence $\tilde\gamma_n/\tilde\gamma_{n-1}$ is decreasing for $n\geq 3$ and convergent, although this remains unproved.

\begin{proof}
We are going to show that $\sfrac{\tilde\gamma_n}{\tilde\gamma_{n-1}}<1$ for $n\geq 13$. Since the same holds for $n\leq 12$ by the above numerical computations, we~will obtain the claimed monotonicity. This monotonicity, together with the numerical fact that $\tilde\gamma_4<1$, implies the corresponding inequality for all $n\geq 4$.

To show that $\tilde\gamma_n/\tilde\gamma_{n-1} <1$ for $n\geq 13$, we~note that
\begin{equation} \label{eq:b}
\frac{\tilde\gamma_n}{\tilde\gamma_{n-1}} = \frac{2(n+1)}{n^2} \,(1-1/n)^{2n} \,\frac{c_{n-1}}{c_n}.
\end{equation}
We write the definition \eqref{eq:defcn} of $c_n$ in the form
\[
c_n= \sum_{m\in\N_0} \frac{ (m+n-1)! }{(n-1)! m!} \frac{1}{(2m+n)^{n+1}}.
\]
In view of \eqref{eq:b} we are mainly interested in a lower bound on the quotient $c_n/c_{n-1}$. Observing that
\begin{multline*}
\frac{ (m+n-1)! }{(n-1)! m!} \frac{1}{(2m+n)^{n+1}}\\
= \frac{m+n-1}{(n-1)(2m+n)} (1- \frac{1}{2m+n})^n \frac{ (m+n-2)! }{(n-2)! m!} \frac{1}{(2m+n-1)^{n}},
\end{multline*}
we get
\begin{equation}\label{eq:a}
\frac{c_n}{c_{n-1}} \geq \frac{1}{n-1} \inf_m \frac{m+n-1}{2m+n} \Bigl(1- \frac{1}{2m+n}\Bigr)^n.
\end{equation}
Hence we have to analyze $\inf_m \theta_n(m)$ with
\begin{equation}\label{eq:theta}
\theta_n(m):= \frac{m+n-1}{2m+n} \Bigl(1- \frac{1}{2m+n}\Bigr)^n.
\end{equation}

We need to carefully analyze the sequence $m \mto \theta_n(m)$. Before doing this, let us provide some heuristics coming from limiting regimes. As~$m=0$, we~have \hbox{$\theta_n(0)=\frac{n-1}{n}(1- \sfrac{1}{n})^n$}, which tends to $e^{-1}$ as $n\to +\infty$. As~$m$ tends to $+\infty$, we~have $\theta_n(m)\to\nobreak\sfrac 12$, but we need uniform lower bounds with respect to $m$ and $n$.

Asymptotically as $n\to +\infty$, a lower bound for $\theta_n(m)$ is given by considering the infimum of the function (think of the change of variable $y=m/n$)
\[
(0,+\infty) \ni y \mto \frac{1+y}{1+2y}\,e^{-1/(2y+1)}.
\]
After the change of variable $u=y+\sfrac 12$, we~have to analyze
\[
(\sfrac 12,+\infty) \ni u \mto \frac{\sfrac12 +u}{2u} e^{-1/2u}= \Bigl(\frac 12 + \frac{1}{4u}\Bigr) e^{-1/2u}.
\]
This function is increasing and its minimum is at $u=\sfrac 12$ and equals $e^{-1}$. Hence we have
\begin{equation}
\frac{m+n}{2m+n}e^{- \spfrac{n}{n+2m}} \geq e^{-1}.
\end{equation}
This leads to
\begin{equation}\label{eq:expdecaytildegammanupper}
\limsup_{n\to +\infty} \frac{\tilde\gamma_n}{\tilde\gamma_{n-1}} \leq 2 e^{-1} \approx 0.735.
\end{equation}
This is closer to the guess.
For the lower bound and using an upper bound for $\limsup_{n\to +\infty} \theta_n(m)$, we~get
\[
\liminf_{n\to +\infty} \frac{\tilde\gamma_n}{\tilde\gamma_{n-1}} \geq 4 e^{-2} \approx 0.541.
\]
This is quite close to the numerics.

After having discussed these heuristics, we~will turn to the proof of rigorous bounds. Returning to $\theta_n(m)$ in \eqref{eq:theta}, we~write
\[
\theta_n(m)\geq \frac{m+n}{2m+n} \Bigl(1- \frac{1}{2m+n}\Bigr)^n - \frac 1n.
\]
Using
\[
- \log (1-x) = x + \int_0^x \frac{t}{1-t} \,dt,\quad x\in [0,1),
\]
we get
\[
x\leq - \log (1-x) \leq x + \frac{1}{1-x} \frac{x^2}{2},\quad x\in [0,1).
\]
With $x=\spfrac{1}{2m+n}$, we~obtain
\[
\Bigl(1- \frac{1}{2m+n}\Bigr)^n= e^{n \log (1- \spfrac{1}{2m+n})}\geq e^{- \spfrac{n}{n+2m} } e^{ - \sfrac{1}{2 (n-1)}}.
\]
Coming back to $\theta_n(m)$ and what we have done for the limsup
\[
\theta_n(m)\geq e^{-1} e^{- \sfrac {1}{2(n-1)}} -\sfrac 1n.
\]
So we finally get
\[
c_n \geq \frac{1}{n-1} \,(e^{-1} e^{- \sfrac 1{2(n-1)}} -\sfrac 1n) c_{n-1}.
\]
Coming back to \eqref{eq:b},
we get
\begin{align*}
\frac{\tilde\gamma_n}{\tilde\gamma_{n-1}}
& \leq \frac{2(n+1)(n-1)}{n^2}(1-1/n)^{2n}(e^{-1} e^{- \sfrac {1}{2(n-1)}} -\sfrac 1n)^{-1} \\
& \leq 2 e^{-1} \,\frac{(n+1)(n-1)}{n^2}(e^{- \sfrac 1{2(n-1)}} -\sfrac en)^{-1} \leq 2 e^{-1}(e^{- \sfrac {1}{2(n-1)}} -\sfrac en)^{-1}.
\end{align*}
So finally, we~have shown that
\begin{equation}\label{eq:quotientgammatilde}
\frac{\tilde\gamma_n}{\tilde\gamma_{n-1}} \leq 2 e^{-1}(e^{- \sfrac {1}{2(n-1)}} -\sfrac en)^{-1}.
\end{equation}
For $n=13$, we~have
\begin{equation} \label{eq:9.9}
e^{- \sfrac 1{2(n-1)}} -\sfrac en \approx 0.75009,
\end{equation}
and, consequently,
\[
\frac{\tilde\gamma_{13}}{\tilde\gamma_{12}} <1.
\]
Looking at the bound \eqref{eq:quotientgammatilde} and its monotonicity with respect to $n$, the bound $\sfrac{\tilde\gamma_n}{\tilde\gamma_{n-1}}<1$ holds for \emph{any} $n\geq 13$. This completes the proof of Proposition \ref{gammatildemono}.
\end{proof}

\subsection{A lower bound on the Sobolev constant on $\mathbb H_n\times\R^k$ for $k\geq 1$}
In this subsection we prove a lower bound on the Sobolev constant $C^\Sob(\mathbb H_n\times\R^k)$ in terms of the Sobolev constant $C^\Sob(\mathbb H_n)$ and the constant appearing in a certain Sobolev interpolation inequality on $\R^k$. Assume $2\leq q<\infty$ if $k\leq 2$ and $2\leq q\leq\spfrac{2k}{k-2}$ if $k>2$, and denote by $C^\GN_q(\R^k)$ the largest possible constant in the inequality, valid for $u\in H^1(\R^k)$,
\begin{subequations}\label{eq:interineq}\begin{equation}
\left(\int_{\R^k} |\nabla u|^2\,dw\right)^\theta \left(\int_{\R^k} | u|^2\,dw\right)^{1-\theta}
\geq C^\GN_q(\R^k) \biggl(\int_{\R^k} |u|^q \,dw\biggr)^{\sfrac2q},
\end{equation}
where
\begin{equation}\label{eq:interineqb}
\theta = k(\sfrac12 - \sfrac1q).
\end{equation}
(The value of $\theta$ is determined by scaling.)
\end{subequations}
For $k=1$ the explicit value of the constant $C^\GN_q(\R)$ is known from a work of Nagy \cite{Na}. For $k\geq 2$ its explicit value is not known, but we will still be able to derive some results in Subsection \ref{sec:maincompk}.

\begin{proposition}\label{lifting}
For all $n,k\in\N$, setting $Q=2n+2$ and $q=\frac{2(Q+k)}{Q+k-2}$,
\[
C^\Sob(\mathbb H_n\times\R^k) \geq C^\GN_{q}(\R^k) \ (C^\Sob(\mathbb H_n))^{\spfrac{Q}{Q+k}} \ \frac{Q+k}{Q^{\spfrac{Q}{Q+k}} \,k^{\spfrac{k}{Q+k}}}.
\]
\end{proposition}

The argument that follows is inspired by the Laptev--Weidl method of lifting in dimension \cite{LaWe} and similar to one used in \cite{FGL}.

\begin{proof}
We begin by applying the inequality
\[
ab \leq \frac1p a^p + \frac1{p'} b^{p'}
\]
with $p=\psfrac{Q+k}{k}$, so $p'=\psfrac{Q+k}{Q}$, and $a=\mu^{-\spfrac{k}{Q+k}} g^{\pspfrac{Q+k-2}{Q+k}}$, $b= \mu^{\spfrac{k}{Q+k}} g^{\spfrac{2}{Q+k}}$. We~get
\[
g \leq \frac{k}{Q+k}\,\mu^{-1} g^{\psfrac{Q+k-2}{k}} + \frac{Q}{Q+k}\,\mu^{\sfrac kQ} g^{\sfrac{2}{Q}}.
\]
We apply this with
\[
g = \int_{\R^k} |u(\zeta,w)|^q\,dw
\quad\text{and}\quad
\mu = c\biggl(\int_{\R^k} |u(\zeta,w)|^2\,dw\biggr)^{\sfrac{Q}{k}},
\]
where $\zeta\in\mathbb H_n$ is fixed and $c>0$ is a parameter. We~obtain
\begin{align*}
\int_{\R^k} |u(\zeta,w)|^q\,dw & \leq c^{-1} \frac{k}{Q+k} \frac{\left(\int_{\R^k} |u(\zeta,w)|^q\,dw\right)^{\psfrac{Q+k-2}{k}}}{\left(\int_{\R^k} |u(\zeta,w)|^2\,dw\right)^{\sfrac{Q}{k}}} \\
& \quad + c^{\sfrac{k}{Q}} \frac{Q}{Q+k} \left(\int_{\R^k} |u(\zeta,w)|^q\,dw\right)^{\sfrac2Q} \int_{\R^k} |u(\zeta,w)|^2\,dw \\
& \leq c^{-1} \frac{k}{Q+k} (C^\GN_{q}(\R^k))^{-\psfrac{Q+k}{k}} \int_{\R^k} |\nabla_w u(\zeta,w)|^2\,dw \\
& \quad
+ c^{\sfrac{k}{Q}} \frac{Q}{Q+k} \left(\int_{\R^k} |u(\zeta,w)|^q\,dw\right)^{\sfrac2Q} \int_{\R^k} |u(\zeta,w)|^2\,dw.
\end{align*}
Here we used the Sobolev interpolation inequality on $\R^k$.
We integrate this inequality with respect to $\zeta$ and obtain
\begin{multline*}
\|u\|_q^q \leq c^{-1} \frac{k}{Q+k} (C^\GN_{q}(\R^k))^{-\psfrac{Q+k}{k}} \| \nabla_w u\|_2^2 \\
+ c^{\sfrac{k}{Q}} \frac{Q}{Q+k} \int_{\mathbb H_n} \left(\int_{\R^k} |u(\zeta,w)|^q\,dw\right)^{\sfrac2Q} \int_{\R^k} |u(\zeta,w)|^2\,dw\,d\zeta.
\end{multline*}
Now by Fubini and Hölder
\begin{align*}
\int_{\mathbb H_n} \biggl(\int_{\R^k} |u(\zeta,w)|^q\,dw \biggr)^{\sfrac2Q}& \biggl(\int_{\R^k} |u(\zeta,w)|^2\,dw \biggr) \,d\zeta \\
& = \int_{\R^k} \biggl(\int_{\mathbb H_n} \left(\int_{\R^k} |u(\zeta,w')|^q\,dw'\right)^{\sfrac2Q} |u(\zeta,w)|^2\,d\zeta\biggr) dw \\
& \leq \| u \|_q^{\sfrac{2q}{Q}} \int_{\R^k} \left(\int_{\mathbb H_n} |u(\zeta,w)|^{\spfrac{2Q}{Q-2}} \,d\zeta\right)^{\psfrac{Q-2}{Q}}\,dw.
\end{align*}
Applying the Sobolev inequality on $\mathbb H_n$ on the right side and inserting the inequality in the above bound, we~get
\begin{multline*}
\|u\|_q^q \leq c^{-1} \frac{k}{Q+k} (C^\GN_{q}(\R^k))^{-\psfrac{Q+k}{k}} \| \nabla_w u\|_2^2 \\
+ c^{\sfrac{k}{Q}} \frac{Q}{Q+k} (C^\Sob(\mathbb H_n))^{-1} \| u \|_q^{\sfrac{2q}{Q}} \int_{\R^k} \sum_{j=1}^n \int_{\mathbb H_n} (|X_ju|^2 + |Y_ju|^2)\,d\zeta \,dw.
\end{multline*}

We finally choose $c$ such that
\[
c^{-1} \frac{k}{Q+k} (C^\GN_{q}(\R^k))^{-\psfrac{Q+k}{k}} = c^{\sfrac{k}{Q}} \frac{Q}{Q+k} (C^\Sob(\mathbb H_n))^{-1} \| u \|_q^{\sfrac{2q}{Q}},
\]
that is, we~choose
\[
c = \left(\sfrac{k}{Q}\right)^{\sfrac{Q}{Q+k}} (C^\GN_{q}(\R^k))^{-\sfrac{Q}{k}} (C^\Sob(\mathbb H_n))^{\spfrac{Q}{Q+k}} \| u \|_q^{-\spfrac{2q}{Q+k}}.
\]
In this way we get
\begin{multline*}
\| u \|_q^q \leq (C^\GN_{q}(\R^k))^{-1} (C^\Sob(\mathbb H_n))^{-\spfrac{Q}{Q+k}} \| u \|_q^{\spfrac{2q}{Q+k}}
\frac{Q^{\sfrac{Q}{Q+k}}\,k^{\spfrac{k}{Q+k}}}{Q+k}\\
\biggl(\| \nabla_w u\|_2^2
+ \int_{\R^k} \sum_{j=1}^n \int_{\mathbb H_n} (|X_ju|^2 + |Y_ju|^2)\,d\zeta \,dw\biggr).
\end{multline*}
This proves the assertion.
\end{proof}

From Proposition \ref{lifting} we obtain the following bound on the Pleijel constant. We~recall that $\tilde\gamma_n$ was defined in \eqref{eq:gammatilden}.

\begin{corollary}\label{pleijellifting}
For all $n,k\in\N$, setting $Q= 2n+2$ and $q=\frac{2(Q+k)}{Q+k-2}$,
\[
\gamma(\mathbb H_n\times\R^k) \leq
(C^\GN_{q}(\R^k))^{-\psfrac{Q+k}{2}} \,\frac{Q^{\sfrac{Q}{2}} \,k^{\sfrac{k}{2}}}{(Q+k)^{\psfrac{Q+k}2}} \,
(4\pi)^{\sfrac{k}{2}} \,\frac{\Gamma(\frac{Q+k+2}{2})}{\Gamma(\frac{Q+2}2)} \,\tilde\gamma_n.
\]
\end{corollary}

\begin{proof}
According to Lemma \ref{fksob}, Proposition \ref{lifting} and \eqref{eq:wnk} we have
\begin{align*}
\gamma(\mathbb H_n\times\R^k) & \leq \left(C^\Sob(\mathbb H_n\times\R^k)\right)^{-\psfrac{Q+k}{2}} \left(\mathcal W(\mathbb H_n \times\R^k)\right)^{-1} \\
& \leq (C^\GN_{q}(\R^k))^{-\psfrac{Q+k}{2}} \ (C^\Sob(\mathbb H_n))^{-\sfrac{Q}{2}} \,\frac{Q^{\sfrac{Q}{2}} \,k^{\sfrac{k}{2}}}{(Q+k)^{\psfrac{Q+k}2}} \\
& \hspace*{4cm} \times
\mathcal W(\mathbb H_n)^{-1} (4\pi)^{\sfrac{k}{2}} \frac{\Gamma(\frac{Q+k+2}{2})}{\Gamma(\frac{Q+2}2)}.
\end{align*}
This is the claimed inequality.
\end{proof}

\subsection{The Pleijel constant $\gamma(\mathbb H_n\times\R^k)$ for $k\geq 1$} \label{sec:maincompk}
Our goal in the present subsection is to prove the following result.

\begin{proposition}\label{largedimension}
If $n\geq 3$ and $k\geq 1$ then $\gamma(\mathbb H_n\times\R^k)<1$.
\end{proposition}

Our proof will rely on the bound on $\gamma(\mathbb H_n\times\R^k)$ from Corollary \ref{pleijellifting}. In~order to exploit this corollary we need a good lower bound on the constant $C^\GN_q(\R^k)$. As~we already mentioned, an explicit expression for this constant is known in dimension $k=1$ due to the work of Nagy \cite{Na}. The idea behind the following result is to reduce the case of $\R^k=\R^{k-1}\times\R$ to that of $\R^{k-1}$ and $\R$, and to iterate. This argument is most easily expressed in the parametrization
\[
q=\frac{2(\gamma+\sfrac k2)}{\gamma+\sfrac k2 -1}.
\]
The theorem will require the assumption $\gamma\geq\sfrac12$, which however will not present a restriction in our application. It~appears in the proof since the constant $C_{\tilde q}^\GN(\R)$ with $\tilde q= \frac{2(\gamma+\sfrac 12)}{\gamma+\sfrac 12 -1}$ is only defined for $\gamma\geq\sfrac12$.

\begin{lemma}
Let $k\geq 1$ and $\gamma\geq\sfrac12$. Then
\[
C_{\frac{2(\gamma+\sfrac k2)}{\gamma+\sfrac k2-1}}^\GN(\R^k) \geq \left(\left(\frac{k\pi}{2}\right)^{\sfrac k2} \frac{\Gamma(\gamma+\frac12)}{\Gamma(\gamma+\frac{k+1}2)} \,\frac{(\gamma+\frac{k-1}{2})^{\gamma+\psfrac{k-1}2}}{(\gamma-\frac12)^{\gamma-\sfrac12}}\right)^{\spfrac{1}{\gamma+\sfrac k2}}
\]
with the convention $(\gamma-\frac12)^{\gamma-\sfrac12}=1$ for $\gamma=\sfrac12$.
\end{lemma}

As the following proof will show, the inequality in the lemma is an equality if $k=1$.

\begin{proof}
For $d\in\N$, let us introduce the constant
\[
L_{\gamma,d}^{(1)} = \frac{\gamma^\gamma \left(\sfrac d2\right)^\sfrac d2}{(\gamma+\frac d2)^{\gamma+\sfrac d2}} \,\Bigl(C_\frac{2(\gamma+\sfrac d2)}{\gamma+\sfrac d2-1}^\GN(\R^d)\Bigr)^{-\gamma-\sfrac d2}.
\]
The main step of the proof will be to show that
\begin{equation}
\label{eq:ltproduct}
L_{\gamma,k}^{(1)} \leq \prod_{j=0}^{k-1} L_{\gamma+\sfrac j2,1}^{(1)}.
\end{equation}
Accepting this for the moment, let us complete the proof. Inserting Nagy's explicit expression for the constant $C_q^\GN(\R)$ \cite{Na} (see also \cite[Th.\,2.48]{FLW}), which reads
\begin{equation}\label{eq:Nag}
C_{q}^\GN(\R) = \left(\frac{(q+2)^{q+2}}{(q-2)^{q-2}2^{2(q+2)}}\right)^\sfrac{1}{2q} \left(\sqrt\pi \,\frac{\Gamma(\frac{q}{q-2})}{\Gamma(\frac{q}{q-2}+\frac12)}\right)^\psfrac{q-2}{q},
\end{equation}
into the definition of $L_{\gamma,1}^{(1)}$, we~obtain
\[
L_{\gamma,1}^{(1)} = \frac{1}{\sqrt\pi}\,\frac{\Gamma(\gamma+1)}{\Gamma(\gamma+\frac12)}\,\frac{(\gamma-\frac12)^{\gamma-\sfrac12}}{(\gamma+\frac12)^{\gamma+\sfrac12}};
\]
see also \cite[Cor.\,5.4]{FLW}. Inserting this into \eqref{eq:ltproduct} and simplifying we obtain the inequality in the lemma.

It remains to prove \eqref{eq:ltproduct}. According to \cite[Prop.\,5.3]{FLW}, the constant $L_{\gamma,d}^{(1)}$ coincides with the so-called one-particle Lieb--Thirring constant. (This observation is essentially due to Lieb and Thirring \cite{LiTh}.) We use now the inequality
\begin{equation}
\label{eq:ltproductstep}
L_{\gamma,k}^{(1)} \leq L_{\gamma,1}^{(1)} \,L_{\gamma+\sfrac12,k-1}^{(1)}
\end{equation}
from \cite{FGL}; see also \cite[Prop.\,5.14]{FLW}. This bound can be translated into the claimed lower bound on $C_{q_k}^\GN(\R^k)$ in terms of the constants $C_{q_1}^\GN(\R)$ and $C_{q_{k-1}}^\GN(\R^{k-1})$. It~is more convenient, however, to stay in the notation of the Lieb--Thirring constants. Iterating the above inequality, we~arrive at \eqref{eq:ltproduct}.

We emphasize that \eqref{eq:ltproductstep} can be translated into a lower bound on $C_{q_k}^\GN(\R^k)$ in terms of the constants $C_{q_1}^\GN(\R)$ and $C_{q_{k-1}}^\GN(\R^{k-1})$, where $q_d=\frac{2(\gamma+\sfrac d2)}{\gamma+\sfrac d2-1}$. This bound could also be proved directly using the method of proof of Proposition \ref{lifting}.
\end{proof}

In our application, we~need $C^\GN_q(\R^k)$ with $q=\frac{2(Q+k)}{Q+k-2}$ where $Q=2n+2$. We~see that $\frac{2(\gamma+\sfrac k2)}{\gamma+\sfrac k2 -1} = q$ if $\gamma=\sfrac Q2$. Note also that, since $Q\geq 4$, the assumption $\gamma\geq\sfrac12$ is satisfied. Therefore in our application we have
\begin{equation}
\label{eq:cgnlower}
C_q^\GN(\R^k) \geq \biggl(\Bigl(\frac{k\pi}{2}\Bigr)^{\sfrac k2} \frac{\Gamma(\frac{Q+1}2)}{\Gamma(\frac{Q+k+1}2)} \,\frac{(\frac{Q+k-1}{2})^{\psfrac{Q+k-1}2}}{(\frac{Q-1}2)^{\psfrac{Q-1}2}}\biggr)^{\spfrac{2}{Q+k}}.
\end{equation}

We will deduce our main result from this bound, Corollary \ref{pleijellifting}, the known facts about $\tilde\gamma_n$ and the monotonicity properties in the following lemma.

\begin{lemma}\label{mono}
For each $Q\geq 4$,
\[
k \mto 2^{2k} \,\frac{\Gamma(\frac{Q+k+1}2)\,\Gamma(\frac{Q+k+2}2)}{(Q+k-1)^{\psfrac{Q+k-1}2} \,(Q+k)^\psfrac{Q+k}2}
\]
is decreasing for $k\geq 0$.
\end{lemma}

\begin{proof}
Denoting
\[
\alpha_k := 2^{2k} \,\frac{\Gamma(\frac{Q+k+1}2)\,\Gamma(\frac{Q+k+2}2)}{(Q+k-1)^{\psfrac{Q+k-1}2} \,(Q+k)^\psfrac{Q+k}2},
\]
we have
\[
\frac{\alpha_k}{\alpha_{k-1}} = 2 \,\frac{(Q+k-2)^{\psfrac{Q+k-2}2}}{(Q+k)^\psfrac{Q+k-2}2},
\]
which we can write as $2(1- x_0^{-1})^{x_0-1}$ with $x_0=\psfrac{Q+k}{2}$. We~will prove momentarily that $x\mto(1-x^{-1})^{x-1}$ is decreasing for $x\geq 1$. Since it is equal to $1/2$ at $x=2$ and since $x_0> 2$ (provided $k-1\geq 0\geq 4-Q$), we~deduce that $2(1- x_0^{-1})^{x_0-1}< 1$, proving that $\alpha_k/\alpha_{k-1}< 1$.

It remains to prove that $x\mto(1-x^{-1})^{x-1}$ is decreasing for $x\geq 1$ or, what is the same, $x\mto (x-1)\ln(1-x^{-1})$ is decreasing for $x\geq 1$. This follows since the derivative of the latter function is $\ln(1-x^{-1}) + x^{-1}\leq 0$ for $x>1$.
\end{proof}

We finally prove the main result of this subsection.

\begin{proof}[Proof of Proposition \ref{largedimension}]
Inserting the bound \eqref{eq:cgnlower} into Corollary \ref{pleijellifting} we obtain
\begin{align}\label{eq:gnbound}
\gamma(\mathbb H_n\times\R^k) & \leq
2^{2k} \,\frac{\Gamma(\frac{Q+k+1}2)\,\Gamma(\frac{Q+k+2}2)}{\Gamma(\frac{Q+1}2)\,\Gamma(\frac {Q+2}2)} \,\frac{(Q-1)^{\psfrac{Q-1}2} \,Q^{\sfrac Q2}}{(Q+k-1)^{\psfrac{Q+k-1}2} \,(Q+k)^\psfrac{Q+k}2} \,\tilde\gamma_n.
\end{align}

First consider $n\geq 4$. According to Lemma \ref{mono}, the right side of \eqref{eq:gnbound} is decreasing with respect to $k$. For $k=0$ it is equal to $\tilde\gamma_n$. Since we know that $\tilde\gamma_n<1$ for $n\geq 4$, we~learn that $\gamma(\mathbb H_n\times\R^k)<1$ for all $n\geq 4$ and $k\geq 0$.

Now let $n= 3$. Again by Lemma \ref{mono}, the right side of \eqref{eq:gnbound} is decreasing with respect to $k$. For $k=1$ it is given by
\[
4 \,\frac{\Gamma(\frac{11}2)}{\Gamma(\frac{9}2)} \,\frac{7^{\sfrac{7}2}}{9^{\sfrac{9}2}} \,\tilde\gamma_3 = 2 \,\left(\frac 79\right)^{\sfrac{7}2} \tilde\gamma_3
\approx 8.872 \cdot 10^{-1} <1.
\]
Here we used the value $\tilde\gamma_3 \approx 1.069$ from Subsection \ref{sec:maincompk0}. This proves that $\gamma(\mathbb H_3\times\R^k)<1$ for all $k\geq 1$. This completes the proof.
\end{proof}

\section{Continuation by looking at a direct product}\label{s10}
This section is a short aside from the main topic of this part. We~will explain that most of the arguments in the previous sections generalize from $\mathbb H_n\times\R^k$ to the case~of
\[
G = G_1 \times G_2,
\]
where $G_1$ and $G_2$ are stratified, nilpotent simply connected Lie groups. If $\Delta^{G_j}$ is the sub-Laplacian on $G_j$, then
\[
\Delta^G = \Delta^{G_1} \otimes I + I\otimes \Delta^{G_2}
\]
is the sub-Laplacian on $G$.

\subsection{Spectral density}
If $Q_1$ and $Q_2$ denote the homogeneous dimensions of $G_1$ and $G_2$, respectively, then
\[
Q:= Q_1 + Q_2
\]
is the homogeneous dimension of $G$. We~recall that for $G_1$ and $G_2$ we have the analogues of formula \eqref{w2}
\[
\1(-\Delta^{G_1} <\mu)(\zeta_1,\zeta_1)=\mathcal W(G_1) \,\mu^{Q_1/2}
\quad\text{for all}\ \zeta_1 \in G_1
\]
and
\[
\1 (-\Delta^{G_2} <\mu)(\zeta_2,\zeta_2)=\mathcal W(G_2) \,\mu^{Q_2/2}
\quad\text{for all}\ \zeta_2 \in G_2,
\]
and also on $G$,
\begin{equation}
\1 (-\Delta^{G_1\times G_2} <\mu)(\zeta,\zeta)= \mathcal W(G) \,\mu^{Q/2} \quad\text{for all}\ \zeta\in G.
\end{equation}
An abstract version of the argument in Lemma \ref{wnk} shows that
\begin{equation}\label{eq:weylconstproduct}
\mathcal W(G) = \mathcal W(G_1) \,\mathcal W(G_2) \,\frac{\Gamma(\frac{Q_1}2 +1) \,\Gamma(\frac{Q_2}2 +1)}{\Gamma (\frac Q2+1)}.
\end{equation}
Indeed, by~the spectral theorem we have
\[
\1 (-\Delta^{G_1\times G_2} <\mu) = \int_0^\infty \left(\1(-\Delta^{G_1} + \lambda <\mu) \otimes I\right) d \left(I \otimes \1(-\Delta^{G_2}<\lambda)\right).
\]
Evaluating this identity for the corresponding integral kernels on the diagonal yields
\[
\mathcal W(G) \,\mu^{Q/2} = \mathcal W(G_1) \,\mathcal W(G_2) \int_0^\infty (\mu-\lambda)_+^{Q_1/2} \,d \left(\lambda^{Q_2/2}\right),
\]
which proves \eqref{eq:weylconstproduct}.

\subsection{Sobolev inequality}
We now discuss the Faber--Krahn part of Pleijel's proof. In~fact, as in the previous section it will be more convenient to work with the Sobolev inequality rather than the Faber--Krahn inequality. Thus, in this subsection we try to get a good lower bound on the Sobolev constant $C^\Sob(G)$ in
\[
\iint_{G_1\times G_2} |\nabla_G u|^2 d\zeta_1\,d\zeta_2 \geq C^\Sob(G) \left(\iint_{G_1\times G_2} |u|^q \,d\zeta_1\,d\zeta_2\right)^\sfrac{2}{q},
\]
where we abbreviate
\[
q= \frac{2(Q_1+Q_2)}{Q_1+Q_2-2}= \frac{2 Q}{Q-2}.
\]
For the validity of the inequality we require $Q\geq 3$.

We use the Sobolev inequality on $G_1$ (assuming $Q_1 \geq 3$),
\[
\int_{G_1} |\nabla_{G_1} u|^2 d\zeta_1 \geq C^\Sob(G_1) \,\left(\int_{G_1} |u|^\spfrac{2Q_1}{Q_1 - 2}\,d\zeta_1\right)^\psfrac{Q_1-2}{Q_1},
\]
as well as the following Sobolev interpolation inequality on $G_2$,
\[
\left(\int_{G_2} |\nabla_{G_2} u|^2\,d\zeta_2\right)^{\theta} \left(\int_{G_2} | u|^2\,d \zeta_2\right)^{1-\theta}
\geq C^\GN_q(G_2) \left(\int_{G_2} |u|^{q} \,d\zeta_2\right)^\sfrac2 {q},
\]
where the parameters $q$ and $\theta$ are related by
\[
\frac 1 q = \frac {1-\theta}{2} + \theta \,\frac{Q_2-2}{2Q_2}.
\]
When $Q_2\geq 3$, this inequality holds in the range $\theta\in[0,1]$, corresponding to $q\in [2,\frac{2Q_2}{Q_2-2}]$, while if $Q_2=1,2$, it~holds in the range $\theta\in[0,\frac{Q_2}2)$, corresponding to $q\in[2,\infty)$.

These inequalities are known to hold; see, for example \cite{Fo, Va} for the case of dimension $Q_1,Q_2\geq 3$. We~do not know of an explicit reference for the inequality in case $Q_2=1,2$, but it can easily be deduced from existing results in the literature. One way is to proceed as in the proof of Lemma \ref{faberkrahn} by `artificially' adding a factor of $\R^k$ (with its additive group structure) and applying the result in dimension $\geq 3$. Another way is to apply \cite[Cor.\,3.6]{BaCoLeSC}, which reduces the validity of the Sobolev inequality to the validity of another inequality. In~the present case we can use \cite[Th.\,8.1]{BaCoLeSC} in view of the validity of the isoperimetric inequality on $G_2$; see, for instance, \cite[Prop.\,1.17 and the following remarks]{HaKo}.

We will apply the Sobolev interpolation inequality with the parameter
\[
\theta_{12} = \frac{Q_2}Q,\quad \frac 1q = \frac 12 - \frac{1}{Q}.
\]
Note that this choice satisfies $\theta_{12}\in[0,1]$ if $Q_2\geq 3$ and $\theta_{12}\in[0,\frac{Q_2}{2})$ if $Q_2=1,2$ (recall that we assume $Q_1\geq 3$), so the Sobolev interpolation inequality is indeed valid.

Proceeding in the same way as in the proof of Proposition \ref{lifting}, we~obtain the following bound for the Sobolev inequality on $G$:
\begin{equation}\label{eq:asym}
C^\Sob(G)\geq C^\GN_q(G_2)\,\bigl(C^\Sob (G_1)\bigr)^{Q_1/Q} \frac{Q} {Q_1^{Q_1/Q} Q_2^{Q_2/Q}},
\end{equation}
where we recall that $q$ is determined according to $\sfrac 1q = \sfrac 12 - \sfrac{1}{Q}$.

When both $Q_1\geq 3$ and $Q_2 \geq 3$, we~can apply Hölder's inequality to deduce that $C^\GN_q(G_2)\geq \bigl(C^\Sob(G_2)\bigr)^{Q_2/Q}$ and obtain the more symmetric bound
\begin{equation}\label{eq:sym}
C^\Sob(G)\geq \bigl(C^\Sob(G_1)\bigr)^{Q_1/Q}\,\bigl(C^\Sob(G_2)\bigr)^{Q_2/Q} \frac{Q} {Q_1^{Q_1/Q} Q_2^{Q_2/Q}}.
\end{equation}

\subsection{Pleijel's Theorem}
The proof of Pleijel's bound in Theorem \ref{pleijelnk} extends with obvious changes to the case of general stratified nilpotent groups. Bounding the Faber--Krahn constant that appears in this bound by the corresponding Sobolev constant as in Lemma \ref{fksob}, we~arrive at a Pleijel-type bound with constant
\[
\tilde\gamma (G) := \bigl(C^\Sob(G)\bigr)^{-\sfrac{Q}{2}}\mathcal W(G)^{-1}.
\]

In particular, when both $Q_1\geq 3$ and $Q_2\geq 3$, we~can insert formula \eqref{eq:weylconstproduct} and inequality \eqref{eq:sym} and arrive at the bound
\begin{align}\label{eq:Pla}
\tilde\gamma (G) \leq \tilde\gamma (G_1) \,\tilde\gamma (G_2) \ \frac{Q_1^{\sfrac{Q_1}2}\,Q_2^{\sfrac{Q_2}2}}{Q^{\sfrac Q2}} \,
\frac{\Gamma (\frac Q2+1)}{\Gamma(\frac{Q_1}2 +1) \,\Gamma(\frac{Q_2}2 +1)}.
\end{align}

Specializing to the case $G_2=\R^k$ we can get

\begin{theorem}
Let $G_1$ a stratified nilpotent group of homogeneous dimension \hbox{$Q_1\geq3$}. Then there exists $k_0(G_1)$ such that for any $k\geq k_0(G_1)$, Pleijel's Theorem holds for the Dirichlet realization in any domain $\Omega\subset G_1\times\R^k$ of the canonical sub-Laplacian associated with the group $G_1 \times \mathbb R^k$.
\end{theorem}

More generally, assume that $G_1$ a stratified nilpotent group of homogeneous dimension $Q_1\geq 3$ and assume that $G_2^{(k)}$ is a sequence of nilpotent groups such that the homogeneous dimensions $Q_2^{(k)}$ tend to $\infty$ and the corresponding constants $\tilde\gamma(G_2^{(k)})$ tend to $0$. Then there exists $k_0$ such that for $k\geq k_0$ we have
\[
\gamma(G_1 \times G_2^{(k)}) <1.
\]

\begin{proof}
We prove the more general assertion described after the theorem. According to Stirling's approximation, we~have
\begin{equation}\label{eq:Stir}
\ln \Bigl(x^{-\psfrac{x+1}2} \Gamma\Bigl(\frac{x+2}2\Bigr) \Bigr) = - \frac x2 \ln(2e) - \frac12 \ln \pi + \mathcal O(x^{-1})
\quad\text{as}\ x\to\infty.
\end{equation}
Writing $Q^{(k)} = Q_1 + Q_2^{(k)}$, this implies that, as $k\to +\infty$,
\[
\frac{(Q_2^{(k)})^{Q_2^{(k)}/2}}{(Q^{(k)})^{Q^{(k)}/2}} \frac{\Gamma (Q^{(k)}/2+1)}{\Gamma(Q_2^{(k)}/2 +1)}\to \Bigl(\frac{1}{2e}\Bigr)^{Q_1/2}.
\]
The assertion then follows from \eqref{eq:Pla} and the assumption $\tilde\gamma(G_2^{(k)})\to 0$.
\end{proof}

\section{A second bound on the Faber--Krahn constant}\label{sec:isoperimetric}

\subsection{Preliminary discussion}
While in the previous two sections we have used the Sobolev constant to get bounds on the Faber--Krahn constant, we~will use in this section the isoperimetric constant for this purpose. In~those previous two sections we have faced the difficulty that the optimal Sobolev constant on $\mathbb H_n\times\R^k$ is not known when $k\geq 1$ and our main work consisted in getting bounds on this constant. Similarly, the optimal isoperimetric constant on $\mathbb H_n\times\R^k$ is not known when $k\geq 0$ and we will try to get bounds on it. We~emphasize that the isoperimetric constant is even unknown in the case of $\mathbb H_n$ (that is, $k=0$), despite a famous conjecture of Pansu \cite{Pa2} and several efforts to prove it, which we cite below.

It is known that, in a relatively general setting, the validity of an isoperimetric inequality implies the validity of a Faber--Krahn inequality; see \cite{Car,BaCoLeSC,Va}. These works also provide bounds on the Faber--Krahn constant in terms of the isoperimetric constant, but the numerical values in these bounds have not been the main concern in these works and, as far as we know, they are not sufficient to deduce Pleijel's theorem. In~contrast to these works, we~will use methods that are more specific to the problem at hand and use the particularities of the Heisenberg group and of Euclidean space. (For instance, in Lemma \ref{folland} we use the form of the Green's function on the Heisenberg group, and in Theorem \ref{isoperk} we use the isoperimetric inequality on Euclidean space.)

\subsection{Definitions}
We study the isoperimetric inequality in $\mathbb H_n\times\R^k$ from the point of view of the theory of sets of finite perimeter. For an introduction to this theory in the Euclidean setting we recommend \cite{Ma}. From this theory we recall, for instance, that the perimeter of a measurable set $E\subset\R^k$ is defined as
\[
\per_{\R^k} E := \sup \biggl\{ \int_E \div \phi \,dw :\ \phi\in C^1_\rc(\R^k,\R^k),\,|\phi|\leq 1\biggr\}
\]
and that for bounded sets with $C^1$ (or even Lipschitz) boundary this coincides with the standard surface area of the set; see \cite[Ex.\,12.5 \& 12.6]{Ma}. The definition of sets of finite perimeter goes back to Caccioppoli and their theory was developed by De~Giorgi.

There is a natural and well-known extension of this theory to the setting of stratified nilpotent groups, but we limit ourselves here to the case of $\mathbb H_n\times\R^k$. In~this case the (horizontal) perimeter of a measurable set $E\subset\mathbb H_n\times\R^k$ is defined to be
\begin{multline*}
\per_{\mathbb H_n\times\R^k}(E):=\sup\biggl\{ \int_E \biggl(\sum_{j=1}^n (X_j\phi+Y_j\phi) + \sum_{i=1}^k W_i\phi\biggr) dx\,dy\,dz\,dw : \\[-8pt]
\phi\in C^1_\rc(\mathbb H_n\times\R^k,\R^{2n+k}),\,|\phi|\leq 1\biggr\}.
\end{multline*}
We denote by $I(\mathbb H_n\times\R^k)$ the largest constant such that for every measurable set $E\subset\mathbb H_n\times\R^k$ of finite measure one has\footnote{More generally, for a nilpotent stratified group $G$ of homogeneous dimension $Q$ the isoperimetric inequality reads
\[
\mathrm{Per}_G (E) \geq I(G)\,|E|^{\psfrac{Q-1}{Q}}.
\]
}
\[
\per_{\mathbb H_n\times\R^k} (E) \geq I(\mathbb H_n\times\R^k) \,|E|^\pspfrac{2n+1+k}{2n+2+k}.
\]
For background and further details we refer to \cite{GaNh}, as well as the textbook \cite{CaDaPaTy} treating the model case of $\mathbb H_1$.

\subsection{Symmetrization}
In this subsection we show that a bound on the isoperimetric constant implies a bound on the Faber--Krahn constant. Arguments in this spirit are known in the Riemannian context and appear, for instance in \cite[Th.\,I.1.5]{BerMe}. As~usual, $j_{\nu,1}$ denotes the first positive zero of the Bessel function $J_\nu$. For the definition and properties of Bessel functions, we~refer to \cite{AbSt}.

\begin{proposition}\label{symmetrization}
For $n\geq 0$ and $k\geq 0$, we~have
\[
C^\FK(\mathbb H_n\times\R^k) \geq I(\mathbb H_n\times\R^k)^2 \,(2n+2+k)^{-2} j_{\psfrac{2n+k}{2},1}^2.
\]
\end{proposition}

The proof shows that equality holds in the Faber--Krahn inequality with the value given by the right side, provided that every superlevel set of the eigenfunction $u$ attains equality in the isoperimetric inequality. This happens in the Euclidean setting where $n=0$. In~this case, we~have
\begin{equation}
\label{eq:isofkexplicit}
\begin{aligned}
I(\R^k) &= k^\psfrac{k-1}{k} |\Sph^{k-1}|^\sfrac1k,\quad\text{and}\\
C^\FK(\R^k)&= k^{2\psfrac{k-1}{k}} |\Sph^{k-1}|^\sfrac2k k^{-2} j_{\psfrac{k-2}{2},1}^2 = \omega_k^{\sfrac2k} j_{\psfrac{k-2}{2},1}^2,
\end{aligned}
\end{equation}
where $\omega_k$ is the measure of the unit ball in $\R^k$; see, for instance, \cite[Th.\,14.1]{Ma} or \cite{Ta} for the first equality and \cite[Th.\,2.54]{FLW} or \cite[Th.\,5.1.2]{LMP} for the second one.

Incidentally, it~follows from \eqref{eq:isofkexplicit} that the inequality in Proposition \ref{symmetrization} can be stated as
\begin{equation}
\label{eq:symmetrization}
C^\FK(\mathbb H_n\times\R^k)\,I(\mathbb H_n\times\R^k)^{-2} \geq
C^\FK(\R^{2n+2+k}) \,I(\mathbb R^{2n+2+k})^{-2}.
\end{equation}

\begin{proof}
We abbreviate $D := 2n+2+k$. (This is the homogeneous dimension of $\mathbb H_n\times\R^k$, but we avoid the notation $Q$ which elsewhere we used for $2n+2$.)

Our goal is to prove inequality \eqref{eq:fk} for a given open set $\Omega\subset\mathbb H_n\times\R^k$ and a given function $u\in C^1_\rc(\Omega)$. (It will then extend to $S^1_0(\Omega)$ by density.) By truncation properties of Sobolev functions we may assume that $u\geq 0$. We~define $u_*$ to be a nonincreasing function on $(0,\infty)$ such that
\[
|\{ u >\tau \}| = \int_{\{u_*(r)>\tau\}} r^{D-1}\,dr
\quad\text{for all}\ \tau>0.
\]
(We could multiply the right side by a constant, for instance $|\Sph^{D-1}|$, to make it look more like $\R^D$, but this does not change the final outcome.) It follows from the layer cake formula that, for any $q>0$,
\begin{equation}
\label{eq:layercake}
\int_{\mathbb H_n\times\R^k} u^q\,dx\,dy\,dz\,dw = \int_0^\infty u_*(r)^q r^{D-1}\,dr.
\end{equation}

We are going to show that, for any $p\geq 1$,
\begin{equation}
\label{eq:rearrangegrad}
\int_{\mathbb H_n\times\R^k} |\nabla_{\mathbb H_n} u|^p\,dx\,dy\,dz\,dw
\geq I_{n,k}^p D^{-p\psfrac{D-1}{D}} \int_0^\infty (-u_*'(r))^p\,r^{D-1}\,dr,
\end{equation}
where, for short,
\[
I_{n,k}:=I(\mathbb H_n\times\R^k).
\]
To prove this, we~note that the co-area formula implies that
\begin{equation}
\label{eq:coarea}
|\{ u>\tau\}| = \int_\tau^\infty \left(\int |\nabla_{\mathbb H_n\times\R^k} u|^{-1} \,d|\nabla_{\mathbb H_n\times\R^k}\1_{\{u>t\}}|\right) dt.
\end{equation}
Let us explain the notation used here. First, we~abbreviated
\[
|\nabla_{\mathbb H_n\times\R^k} u| = \sqrt{ \sum_{j=1}^n ((X_j u)^2 + (Y_j u)^2) + \sum_{i=1}^k (W_i u)^2 }.
\]
Next, in the proof of the co-area formula one shows that the sets $\{ u> t\}$ have finite perimeter for almost every $t$, and $|\nabla_{\mathbb H_n\times\R^k}\1_{\{u>t\}}|$ denotes the corresponding perimeter measure. To recall the definition of the latter, let us note that, by~Riesz's theorem, if~$E$ is a set of finite perimeter, then there is an $\R^{2n+k}$-valued Radon measure $\mu_E$ on $\mathbb H_n\times\R^k$ such that
\[
\int_E \biggl(\sum_{j=1}^n (X_j\phi+Y_j\phi) + \sum_{i=1}^k W_i\phi\biggr) dx\,dy\,dz\,dw = - \int_E \phi\cdot d\mu_E
\]
for all $\phi\in C^1_\rc(\mathbb H_n\times\R^k,\R^{2n+k})$. The perimeter measure $|\nabla_{\mathbb H_n\times\R^k} \1_E|:= |\mu_E|$ is the corresponding total variation measure. Note that, in particular, we~have
\[
\per_{\mathbb H_n\times\R^k} E = \int d|\nabla_{\mathbb H_n\times\R^k}\1_E|.
\]
(Whenever $\{u=t\}$ is a sufficiently smooth hypersurface the perimeter measure $|\nabla_{\mathbb H_n\times\R^k}\1_{\{u>t\}}|$ is concentrated on $\{u=t\}$.)

For a textbook proof of the co-area formula in the Euclidean case we refer to \cite[Th.\,13.1]{Ma}. The corresponding formula in the case of stratified nilpotent groups appears in \cite[Th.\,5.2]{GaNh} and \cite[Cor.\,2.3.5]{FrSeSC}.

It follows from \eqref{eq:coarea} that
\[
\frac{d}{d\tau}|\{ u>\tau\}| = - \int |\nabla_{\mathbb H_n\times\R^k} u|^{-1} \,d|\nabla_{\mathbb H_n\times\R^k}\1_{\{u>\tau\}}|.
\]
By the assumed isoperimetric inequality and Hölder's inequality with \hbox{$\sfrac1p + \sfrac1{p'}=1$}, (in fact, only $p=2$ will be relevant), we~have
\begin{align*}
I_{n,k} |\{ u>\tau\}|^\psfrac{D-1}{D} & \leq \per_{\mathbb H_n\times\R^k} \{u>\tau\}
= \int d|\nabla_{\mathbb H_n\times\R^k}\1_{\{u>\tau\}}| \\
& \leq \left(\int |\nabla_{\mathbb H_n\times\R^k} u|^{-1} d|\nabla_{\mathbb H_n\times\R^k}\1_{\{u>\tau\}}|\right)^\sfrac1{p'} \\
& \hspace*{2cm} \times \biggl(\int |\nabla_{\mathbb H_n\times\R^k} u|^{p-1} d|\nabla_{\mathbb H_n\times\R^k}\1_{\{u>\tau\}}|\biggr)^\sfrac 1p.
\end{align*}
The previous two relations combined give
\[
\int |\nabla_{\mathbb H_n} u|^{p-1} d|\nabla_{\mathbb H_n}\1_{\{u>\tau\}}| \geq I_{n,k}^p |\{ u>\tau\}|^{p\psfrac{D-1}{D}} \Bigl(- \frac{d}{d\tau} |\{ u>\tau\}|\Bigr)^{-p+1}.
\]
By the co-area formula again, we~deduce that
\begin{align*}
\int_{\mathbb H_n\times\R^k} |\nabla_{\mathbb H_n} u|^p \,dx\,dy\,dz\,dw & = \int_0^\infty \int |\nabla_{\mathbb H_n\times\R^k} u|^{p-1} \,d|\nabla_{\mathbb H_n\times\R^k}\1_{\{u>\tau\}}| \,d\tau \\
& \geq I_{n,k}^p \int_0^\infty |\{ u>\tau\}|^{p\psfrac{D-1}{D}} \Bigl(- \frac{d}{d\tau} |\{ u>\tau\}|\Bigr)^{-p+1}\,d\tau.
\end{align*}
Note that the right side depends on $u$ only through the function $\tau\mto|\{ u>\tau\}|$. We~express it through the function $u_*$.

For the sake of simplicity we assume that $u_*$ is strictly monotone. The general case can be treated as well, but the idea becomes clearer in this special case. The general case can be handled along the lines of \cite[Th.\,15.29]{Leo}. Under the strict monotonicity assumption we can then define a function $[0,\|u\|_\infty]\ni\tau \mto R_\tau\in (0,\infty)$ by
\[
u_*(R_\tau) = \tau.
\]
It follows that
\[
|\{ u>\tau \}| = \int_{\{u_*(r)>\tau\}} r^{D-1}\,dr = D^{-1} R_\tau^D.
\]
Moreover, one can show that, if~$u$ is weakly differentiable, then $u_*$ is absolutely continuous (an argument of this type in the Euclidean case can be found in \cite[Th.\,15.23]{Leo}) and we have
\[
u_*'(R_\tau) \dot R_\tau = 1
\]
where the dot denotes differentiation with respect to $\tau$. It~follows that
\begin{align*}
\int_0^\infty |\{ u>\tau\}|^{p\psfrac{D-1}{D}}& \Bigl(- \frac{d}{d\tau} |\{ u>\tau\}|\Bigr)^{-p+1}\,d\tau \\
& = - D^{-p\psfrac{D-1}{D}} \int_0^\infty R_\tau^{D-1} \left(- u_*'(R_\tau)\right)^p \dot R_\tau \,d\tau \\
& = D^{-p\psfrac{D-1}{D}} \int_0^\infty r^{D-1} (-u_*'(r))^p\,dr.
\end{align*}
Thus, we~have completed the proof of \eqref{eq:rearrangegrad}.

At this point we specialize to the case $p=2$ and recall the one-dimensional inequality
\[
\int_0^R U'(r)^2 r^{D-1}\,dr \geq j_{\psfrac{D-2}{2},1}^2 \,R^{-2} \int_0^R U(r)^2 r^{D-1}\,dr,
\]
valid for all $R>0$ and all absolutely continuous functions $U$ on $(0,R]$ with \hbox{$U(R)=0$}. This follows from the Faber--Krahn inequality in $\R^D$ when restricted to radial functions; see, e.g., \cite[Th.\,2.54]{FLW}. We~apply the above one-dimensional inequality to $U=u_*$ and $R=R_0$. If we insert
\[
|\Omega| = |\{ u>0\}| = D^{-1} R_0^D,
\]
we obtain the claimed inequality.
\end{proof}

\subsection{Pansu's conjecture and its consequences}
The validity of the isoperimetric inequality in $\mathbb H=\mathbb H_1$ (that is, the fact that $I(\mathbb H)>0$) is due to Pansu \cite{Pa}. In~\cite{Pa2} Pansu made a conjecture about the set that realizes the optimal constant $I(\mathbb H)$. In~the references given below this conjecture is generalized to $\mathbb H_n$ with $n\geq 2$ and, after a computation, we~find:
\begin{equation}
\label{eq:pansusconj}
I(\mathbb H_n) = \frac{2n}{2n+1} \,\frac{(2n+2)^{\frac{2n+1}{2n+2}} \,\Gamma(\tfrac{2n+3}{2})^\frac 1{2n+2}\, \pi^\frac{2n+1}{2(2n+2)}\, 2^\frac 1{n+1}}{\Gamma(n+1)^\frac 1{n+1}} \qquad\begin{array}{l} \text{(Pansu's}\\\text{conjecture);}\end{array}
\end{equation}
see, e.g., \cite{DGN2}. (The extra factor of $2^\spfrac1{n+1}$ compared to their formula comes from the fact that they use $\partial_{x_j} +\frac12 y_j\partial_z$ etc., while we use $\partial_{x_j} + 2y_j\partial_z$ etc.)

Pansu's conjecture is generally believed to be true and has been verified under certain additional assumptions \cite{RR,DGN2,RR2,Ri,Mon,MoRi}; see also \cite{LeRi,LeMa,ChChHwYa,ChHwMaYa} for a sample of contributions to this problem, as well as the textbook \cite{CaDaPaTy} for an introduction to the isoperimetric problem on the Heisenberg group.

We recall that in Subsection \ref{sec:maincompk0} we have proved Pleijel's theorem in $\mathbb H_n$ when $n\geq 4$. Now we will prove it in the remaining dimensions, assuming the validity of Pansu's conjecture.

\begin{corollary}\label{pansupleijelk0}
Assume that Pansu's conjecture holds for some $n\in\{1,2,3\}$. Then the Pleijel constant $\gamma(\mathbb H_n)$ satisfies for the corresponding value of $n$
\begin{equation}\label{eq:panpl1}
\gamma(\mathbb H) \leq 0.406114,
\quad
\gamma(\mathbb H_2) \leq 0.155327,
\quad
\gamma(\mathbb H_3) \leq 0.0641172.
\end{equation}
\end{corollary}

\begin{proof}
We abbreviate $Q=2n+2$. The conjectured value for $I(\mathbb H_n)$ given above leads, via Proposition \ref{symmetrization}, to the bound on the Faber--Krahn constant
\begin{align*}
C^\FK(\mathbb H_n)
& \geq \frac{Q^{2\,\frac{Q-1}{Q}}(Q-2)^2 \,\Gamma(\tfrac{Q+1}{2})^{\frac 2Q} \pi^\frac{Q-1}{Q} 4^{\frac 2Q}}{(Q-1)^2\,\Gamma(\frac Q2)^\frac 4Q}\,
Q^{-2} j_{\frac{Q-2}{2},1}^2 \\
& =
\frac{(Q-2)^2 \,\Gamma(\tfrac{Q+1}{2})^{\frac 2Q} \pi^\frac{Q-1}{Q} 4^{\frac 2Q}}{Q^{\frac 2Q}\,(Q-1)^2\,\Gamma(\frac Q2)^\frac 4Q} j_{\frac{Q-2}{2},1}^2.
\end{align*}
This leads to the bound on the Pleijel constant
\[
\gamma(\mathbb H_n) \leq (C^\FK(\mathbb H_n))^{-\sfrac Q2} \mathcal W(\mathbb H_n)^{-1}
= \frac{Q (Q-1)^Q\,\Gamma(\frac Q2)^2} {(Q-2)^Q \,\Gamma(\tfrac{Q+1}{2}) \pi^{\frac{Q-1}2} 4} \ j_{\frac{Q-2}{2},1}^{-Q} \ \mathcal W(\mathbb H_n)^{-1}.
\]
Using Subsection \ref{sec:weylheis1}, which gives
\[
\mathcal W(\mathbb H) = \frac1{128},\quad
\mathcal W(\mathbb H_2) = \frac{1}{48^2\pi},\quad
\mathcal W(\mathbb H_3) = \frac{1}{2^7 \cdot 768}\,\frac{12 - \pi^2}{\pi^2},
\]
and the values \cite[Table 9.5]{AbSt}
\[
j_{1,1}\sim 3.8317,\quad
j_{2,1}\sim 5.1356,\quad
j_{3,1}\sim 6.3802,
\]
we get \eqref{eq:panpl1} and the corollary.
\end{proof}

\subsection{A bound on the isoperimetric constant in $\mathbb H_n$}
In this subsection we prove a lower bound on $I(\mathbb H_n)$ which is reasonably good in small dimensions $n=1$ and $n=2$. To state this bound, we~need two constants,
\begin{equation}\label{eq:sobrepconst}
\mathcal C_n := \frac{2^{n-3} n\Gamma(\frac n2)^2}{\pi^{n+1}}
\end{equation}
and, setting $Q=2n+2$,
\begin{align}\label{eq:bathtubconst}
\mathcal C_n' & := Q^\sfrac1Q \left(\pi^n \frac{\Gamma(\frac{2Q-1}{2(Q-1)})}{\Gamma(\frac{2Q-1}{2(Q-1)}+\frac12)} \frac{\Gamma(\frac n2+\frac{Q}{4(Q-1)})}{\Gamma(\frac n2+\frac{Q}{4(Q-1)}+\frac12)}
\frac{\Gamma(1+\frac{Q}{2(Q-1)})}{\Gamma(n+\frac{Q}{2(Q-1)})}\right)^\psfrac{Q-1}{Q}.
\end{align}
For $n=1$ we can write the latter constant without Gamma functions, see \eqref{eq:bathtubconst1} below. In~any dimension it can easily be evaluated numerically.

These two constants appear in our bound on the isoperimetric constant. More precisely, we~will prove the following theorem.

\begin{theorem}\label{isoper}
For any $n\in\N$,
\[
I(\mathbb H_n) \geq \mathcal C_n^{-1} (\mathcal C_n')^{-1}.
\]
\end{theorem}

We will prove that for any measurable set $E\subset\mathbb H_n$
\begin{equation}
\label{eq:isopersob}
\int_E u(\zeta)\,d\zeta \leq \mathcal C_n \mathcal C_n' |E|^\sfrac1Q \int_{\mathbb H_n} |\nabla_{\mathbb H_n} u(\zeta)|\,d\zeta.
\end{equation}
Here and below, we~denote variables in $\mathbb H_n$ by $\zeta$ and we write
\[
|\nabla_{\mathbb H_n} u(\zeta)| := \sqrt{ \sum_{j=1}^n ((X_ju)(\zeta)^2 +(Y_ju)(\zeta)^2) }.
\]
If we prove this inequality for, say, Lipschitz functions $u$ with compact support, then, by~a standard approximation argument, it~will extend, with the same constant, to all functions $g$ of bounded variation satisfying $|\{|u|>\lambda\}|<\infty$ for all $\lambda>0$. (The bounded variation condition is understood in the Heisenberg sense, as discussed, for instance, in \cite[\S 5.1]{CaDaPaTy}.) In particular \eqref{eq:isopersob} is valid for the characteristic function of a set of finite (horizontal) perimeter. Since
\[
\int_{\mathbb H_n} |\nabla_{\mathbb H_n} \1_E(\zeta)|\,d\zeta = \per_{\mathbb H_n} E
\quad\text{and}\quad
\int_E \1_E(\zeta)\,d\zeta = |E|,
\]
we obtain the inequality stated in the theorem. Thus, it~remains to prove \eqref{eq:isopersob}.

Our analysis is based on the following well-known representation formula, where we use the notation
\[
\| (x,y,z) \| = ((|x|^2+|y|^2)^2 + z^2)^\sfrac14
\quad\text{for}\ (x,y,z)\in\mathbb H_n.
\]

\begin{lemma}\label{folland}
With $\mathcal C_n$ from \eqref{eq:sobrepconst},
\[
u(\zeta) = \mathcal C_n \sum_{j=1}^n \int_{\mathbb H_n} \frac{(X_j\|{\cdot}\|)(\zeta^{-1} \eta) (X_ju)(\eta) + (Y_j\|{\cdot}\|)(\zeta^{-1}\eta) (Y_ju)(\eta)}{\|\zeta^{-1} \eta\|^{Q-1}}\,d\eta.
\]
Here $\zeta^{-1}\eta$ denotes the inverse of $\zeta$ composed with $\eta$ in the sense of the Heisenberg group.
\end{lemma}

\begin{proof}
As shown by Folland \cite{Fo0},
\[
u(\zeta) = - \tilde{\mathcal C}_n \sum_{j=1}^n \int_{\mathbb H_n} \|\zeta^{-1} \eta\|^{2-Q} ((X_j^2+Y_j^2) u)(\eta)\,d\eta
\]
with
\[
\tilde{\mathcal C}_n = \frac{2^{n-4}\Gamma(\frac n2)^2}{\pi^{n+1}}.
\]
The explicit expression for $\tilde{\mathcal C}_n$ can be found, for instance, in \cite[Eq.~between (1.2) and (1.3)]{BrFoMo} or \cite[Th.\,1.2]{CoLu}. (Meanwhile, there seems to be a computational error in the constant in \cite[Th.\,5.15]{CaDaPaTy}.)

Integrating by parts, we~find
\begin{multline*}
- \int_{\mathbb H_n} \|\zeta^{-1} \eta\|^{2-Q} ((X_j^2+Y_j^2) u)(\eta)\,d\eta \\
= \int_{\mathbb H_n} \left((X_j \|\zeta^{-1} \eta\|^{2-Q}) (X_j u)(\eta) + (Y_j \|\zeta^{-1} \eta\|^{2-Q}) (Y_j u)(\eta)\right) d\eta.
\end{multline*}
On the right side, the vector fields $X_j$ and $Y_j$ act with respect to the $\eta$ variable. By~left invariance, we~find
\[
X_j \|\zeta^{-1} \eta\|^{2-Q} = (X_j \|{\cdot}\|^{2-Q})(\zeta^{-1}\eta) = (2-Q)\|\zeta^{-1}\eta\|^{1-Q} (X_j\|{\cdot}\|)(\zeta^{-1}\eta)
\]
and similarly for $Y_j$. This proves the claimed formula.
\end{proof}

We deduce from Lemma \ref{folland} that for any measurable set $E\subset\mathbb H_n$,
\[
\int_E u(\zeta)\,d\zeta = \mathcal C_n \int_{\mathbb H_n} \sum_{j=1}^n \int_E \frac{(X_j\|{\cdot}\|)(\zeta^{-1} \eta) (X_ju)(\eta) + (Y_j\|{\cdot}\|)(\zeta^{-1}\eta) (Y_ju)(\eta)}{\|\zeta^{-1} \eta\|^{Q-1}}\,d\zeta \,d\eta.
\]
Here we used Fubini's theorem. The arguments given in Lemma \ref{bathtub} below show that the integrals are absolutely convergent, so the interchange of integrals is justified.

Our goal is to bound
\[
\sum_{j=1}^n \int_E \frac{(X_j\|{\cdot}\|)(\zeta^{-1} \eta) (X_ju)(\eta) + (Y_j\|{\cdot}\|)(\zeta^{-1}\eta) (Y_ju)(\eta)}{\|\zeta^{-1} \eta\|^{Q-1}}\,d\zeta
\]
from above, for any fixed $\eta$. Abbreviating $a_j:=(X_ju)(\eta)$, $b_j:=(Y_ju)(\eta)$ and $F:=E^{-1}\eta=\{ \zeta^{-1} \eta:\ \zeta\in E \}$, we~need to bound
\[
\sum_{j=1}^n \int_F \frac{a_j (X_j\|{\cdot}\|)(\xi) + b_j (Y_j\|{\cdot}\|)(\xi)}{\|\xi\|^{Q-1}}\,d\xi
\]
from above. Note that
\[
|F|=|E^{-1} \eta|
\quad\text{and}\quad
|a|^2+|b|^2 = |\nabla_{\mathbb H_n} g(v)|^2.
\]

\begin{lemma}\label{bathtub}
For any $a,b\in\R^n$ and any measurable $F\subset\mathbb H_n$,
\[
\sum_{j=1}^n \int_F \frac{a_j (X_j\|{\cdot}\|)(\xi) + b_j (Y_j\|{\cdot}\|)(\xi)}{\|\xi\|^{Q-1}}\,d\xi
\leq \mathcal C_n' |F|^\sfrac{1}{Q} (|a|^2+|b|^2)^{\sfrac12}
\]
with $\mathcal C_n'$ from \eqref{eq:bathtubconst}.
\end{lemma}

Let us accept this lemma for the moment and use it to complete the proof of \eqref{eq:isopersob}. According to the above argument, we~find, for any $\eta\in\mathbb H_n$,
\[
\sum_{j=1}^n \int_E \frac{(X_j\|{\cdot}\|)(\zeta^{-1} \eta) (X_ju)(\eta) + (Y_j\|{\cdot}\|)(\zeta^{-1}\eta) (Y_ju)(\eta)}{\|\zeta^{-1} \eta\|^{Q-1}}\,d\zeta
\leq \mathcal C_n' |E|^\sfrac1Q |\nabla_{\mathbb H_n} u(\eta)|.
\]
Consequently, we~have shown that
\[
\int_E u(\zeta)\,d\zeta \leq \mathcal C_n \mathcal C_n' |E|^\sfrac1Q \int_{\mathbb H_n} |\nabla_{\mathbb H_n} u(\eta)| \,d\eta.
\]
This is the claimed inequality \eqref{eq:isopersob}. It~remains to prove Lemma \ref{bathtub}.

\begin{proof}[Proof of Lemma \ref{bathtub}]
Let us abbreviate
\[
I[F] := \sum_{j=1}^n \int_F \frac{a_j (X_j\|{\cdot}\|)(\xi) + b_j (Y_j\|{\cdot}\|)(\xi)}{\|\xi\|^{Q-1}}\,d\xi.
\]
We think of the inequality in the lemma as an optimization problem, where we want to maximize $I[F]$ among all sets $F$ of a given measure. By~homogeneity the value of this fixed measure is irrelevant.

We know by the bathtub principle (see, for instance, \cite[Th.\,1.14]{LiLo}) that there is an optimal set $F_*$ for the inequality and that this optimal set is given by
\[
F_* = \biggl\{ \xi \in \mathbb H_n :\ \sum_{j=1}^n \frac{a_j (X_j\|{\cdot}\|)(\xi) + b_j (Y_j\|{\cdot}\|)(\xi)}{\|\xi\|^{Q-1}} >\kappa \biggr\}
\]
for some $\kappa>0$. In~fact, as already mentioned, by~homogeneity, the value of $\kappa$ is immaterial and we will set it equal to $1$ in what follows.
Thus, we~have, for any set $F\subset\mathbb H_n$,
\[
\frac{I[F]}{|F|^\sfrac1Q} \leq \frac{I[F_*]}{|F_*|^\sfrac1Q},
\]
and our task is to compute the right side. More precisely, we~want to show that
\[
\frac{I[F_*]}{|F_*|^\sfrac1Q} = \mathcal C_n' (|a|^2 + |b|^2)^{\sfrac12}
\]
with $\mathcal C_n'$ from \eqref{eq:bathtubconst}.

We begin by bringing $I[F]$ in a more explicit form. Writing $\xi=(x,y,t)$, we~compute
\begin{align*}
(X_j\|{\cdot}\|)(\xi) & = \|\xi\|^{-3} \bigl((|x|^2+|y|^2) x_j + y_j t\bigr),\\
(Y_j\|{\cdot}\|)(\xi) & = \|\xi\|^{-3} \bigl((|x|^2+|y|^2) y_j - x_j t\bigr),
\end{align*}
so
\begin{multline*}
\frac{a_j (X_j\|{\cdot}\|)(\xi) + b_j (Y_j\|{\cdot}\|)(\xi)}{\|\xi\|^{Q-1}} \\[-8pt]
= \frac{(|x|^2+|y|^2)(a\cdot x) + (a\cdot y)t + (|x|^2+|y|^2)(b\cdot y) - (b\cdot x)t}{\|\xi\|^{Q+2}}.
\end{multline*}
In particular,
\begin{align*}
I[F] = \int_F \frac{(|x|^2+|y|^2)(a\cdot x) + (a\cdot y)t + (|x|^2+|y|^2)(b\cdot y) - (b\cdot x)t}{\|\xi\|^{Q+2}}\,d\xi
\end{align*}
and
\[
F_* = \biggl\{ \xi\in\mathbb H_n:\ \frac{(|x|^2+|y|^2)(a\cdot x) + (a\cdot y)t + (|x|^2+|y|^2)(b\cdot y) - (b\cdot x)t}{\|\xi\|^{Q+2}} > 1\biggr\}.
\]

We claim that $|F_*|$ and $I[F_*]$ depend on $a$ and $b$ only through $|a|^2+|b|^2$. Indeed, from the above expression it is obvious that the two quantities do not change if $a$ and $b$ are simultaneously rotated by the same rotation matrix. This implies that the two quantities depend on $a$ and $b$ only through $|a|$, $|b|$ and $a\cdot b$. We~also see that the two quantities do not change if we apply a two-dimensional rotation to $(a_j,b_j)$. This implies that the two quantities depend on $a_j$ and $b_j$ only through $a_j^2+b_j^2$. This proves the claim.

As a consequence, we~may set $\alpha:=\sqrt{|a|^2+|b|^2}$ and assume that $a=\alpha e_n$ and $b=0$. Thus, we~have
\[
F_* = \Big\{ \xi \in\mathbb H_n:\ \alpha \frac{(|x|^2+|y|^2)x_n + y_n t}{\|\xi\|^{Q+2}} >1 \Big\}
\]
and we need to compute
\[
|F_*|
\quad\text{and}\quad
I[F_*] = \int_{F_*} \alpha \frac{(|x|^2+|y|^2)x_n + y_n t}{\|\xi\|^{Q+2}} \,d\xi.
\]

To perform these computations, we~introduce coordinates
\[
x_j = r \sqrt{\sin\theta} \cos\phi_j \,\omega_j,
\quad
y_j = r \sqrt{\sin\theta} \sin\phi_j \,\omega_j,
\quad
t = r^2 \cos\theta,
\]
where
\[
r>0,\ \theta\in(0,\pi),\ \phi=(\phi_1,\ldots,\phi_n)\in(-\pi,\pi)^n=:T^n \ \text{and}\ \omega \in\Sph^{n-1}\cap(0,\infty)^n=:\Sigma.
\]
We~claim that in this parametrization, the measure is given by
\begin{equation}
\label{eq:measure}
dx\,dy\,dt = r^{Q-1}dr\,\sin^{n-1}\theta\,d\theta\,d\mu(\omega)\,d\phi,
\end{equation}
where $d\phi$ is standard Lebesgue measure on $T^n$ and where $d\mu(\omega) = \omega_1\cdots\omega_n\,d\omega$ with the standard surface measure $d\omega$ on $\Sph^{n-1}$. More explicitly, if~we parametrize \hbox{$\omega\in\Sigma$}~by
\begin{align*}
\omega_1 & = \sin\theta_{n-1} \sin\theta_{n-2} \cdots \sin\theta_1,\\
\omega_2 & = \sin\theta_{n-1} \sin\theta_{n-2} \cdots \cos\theta_1,\\
\ldots \\
\omega_{n-1} & = \sin\theta_{n-1} \cos\theta_{n-2},\\
\omega_n & = \cos\theta_{n-1},
\end{align*}
with $0<\theta_j<\sfrac\pi 2$, then
\[
d\mu(\omega) = \omega_1\cdots\omega_n\,d\omega = \prod_{j=1}^{n-1} \sin^{2j-1} \theta_j\cos\theta_j\,d\theta_j.
\]
The latter formula follows immediately from $d\omega = \prod_{j=1}^{n-1} \sin^{j-1}\theta_j\,d\theta_j$.

We point out that when $n=1$, then $\Sigma=\{1\}$ and $\mu$ is trivial in the sense that $\mu(\{1\})=1$.

Let us provide the details for formula \eqref{eq:measure}. First, for each $j$, we~introduce polar coordinates
\[
x_j = r_j \cos\phi_j,
\quad
y_j = r_j \sin\phi_j,
\]
and note that
\[
dx_j\,dy_j = r_j dr_j \,d\phi_j.
\]
Next, we~consider $(r_1,\ldots,r_n)$ as an element of $\R^n$ and introduce hyperspherical coordinates in $\R^n$,
\[
r_j = \rho \omega_j
\]
with $\rho>0$ and $\omega=(\omega_1,\ldots,\omega_n)\in\Sph^{n-1}\cap(0,\infty)^n=\Sigma$. We~have
\[
dr_1\cdots dr_n = \rho^{n-1}d\rho\,d\omega,
\]
so
\[
dx\,dy = \rho^{2n-1} \,d\rho\,\omega_1\cdots\omega_n \,d\omega = \rho^{2n-1}\,d\rho\,d\mu(\omega).
\]
Finally, we~set
\[
\rho = r \sqrt{\sin\theta},
\quad
t= r^2 \cos\theta
\]
and compute easily
\[
d\rho\,dt = r^2 (\sin\theta)^{-\sfrac12}\,dr\,d\theta.
\]
Inserting this into the above formula for $dx\,dy$, we~obtain the claimed formula \eqref{eq:measure}.

In these coordinates we have
\[
\frac{(|x|^2+|y|^2)x_n + y_n t}{\|\xi\|^{Q+2}} = \sqrt{\sin\theta}\ \omega_n\,\frac{\sin\theta \cos\phi_n + \cos\theta \sin\phi}{r^{Q-1}} = \frac{\sqrt{\sin\theta}\ \omega_n\,\sin(\theta+\phi_n)}{r^{Q-1}}.
\]
In particular, the constraint $\alpha \psfrac{(|x|^2+|y|^2)x_n + y_n t}{\|\xi\|^{Q+2}}>1$ can be written as
\[
r< \left(\alpha \sqrt{ \sin\theta}\,\omega_n \sin_+(\theta+\phi_n)\right)^\spfrac{1}{Q-1}.
\]
This allows us to carry out the $r$-integration and to obtain
\begin{align*}
|F_*| & = \int_{T^n} d\phi \int_\Sigma d\mu(\omega) \int_0^\pi d\theta \,\sin^{n-1}\theta
\int_0^\infty dr \,r^{Q-1} \1_F \\
& = \alpha^\spfrac{Q}{Q-1} \frac1Q \int_{T^n} d\phi \int_\Sigma d\mu(\omega) \int_0^\pi d\theta \,\sin^{n-1}\theta \left(\sqrt{ \sin\theta}\,\omega_n \sin_+(\theta+\phi_n)\right)^\spfrac{Q}{Q-1}
\end{align*}
and
\begin{align*}
I[F_*] & = \alpha \int_{T^n} d\phi \int_\Sigma d\mu(\omega) \int_0^\pi d\theta \,\sin^{n-1}\theta \int_0^\infty dr \,r^{Q-1} \1_F \frac{\sqrt{\sin\theta}\,\omega_n \sin(\theta+\phi_n)}{r^{Q-1}} \\
& = \alpha \int_{T^n} d\phi \int_\Sigma d\mu(\omega) \int_0^\pi d\theta \,\sin^{n-1}\theta \\
& \quad\quad\quad \times \sqrt{\sin\theta}\,\omega_n \sin(\theta+\phi_n)
\left(\alpha \sqrt{ \sin\theta}\,\omega_n \sin_+(\theta+\phi_n)\right)^\spfrac{1}{Q-1} \\
& = \alpha^\spfrac{Q}{Q-1} \int_{T^n} d\phi \int_\Sigma d\mu(\omega) \int_0^\pi d\theta \,\sin^{n-1}\theta \left(\sqrt{ \sin\theta}\,\omega_n \sin_+(\theta+\phi_n)\right)^\spfrac{Q}{Q-1}.
\end{align*}
Thus,
\begin{multline*}
\frac{I[F_*]}{|F_*|^\sfrac1Q} = \alpha Q^\sfrac1Q\\[-5pt]
\times \biggl(\int_{T^n} d\phi \int_\Sigma d\mu(\omega) \int_0^\pi d\theta \,\sin^{n-1}\theta \left(\sqrt{ \sin\theta}\,\omega_n \sin_+(\theta+\phi_n)\right)^\spfrac{Q}{Q-1}\biggr)^\psfrac{Q-1}{Q}.
\end{multline*}

It remains to compute the integral on the right side. The $(\phi_1,\ldots,\phi_{n-1})$-integral can be computed immediately. Moreover, for fixed $\theta$, we~can compute the $\phi_n$-integral by translation invariance. We~obtain
\begin{multline*}
\int_{T^n} d\phi \int_\Sigma d\mu(\omega) \int_0^\pi d\theta \,\sin^{n-1}\theta \left(\sqrt{ \sin\theta}\,\omega_n \sin_+(\theta+\phi_n)\right)^\spfrac{Q}{Q-1} \\
= (2\pi)^{n-1} \int_0^\pi d\phi\,\sin^\spfrac{Q}{Q-1}\phi \int_0^\pi d\theta \,\sin^{n-1+\sfrac{Q}{2(Q-1)}}\theta \int_\Sigma d\mu(\omega)\,\omega_n^\spfrac{Q}{Q-1}.
\end{multline*}
Using the beta function identity\vspace*{-3pt}
\[
2\int_0^\sfrac\pi2 dt\,(\sin^{2z_1-1}t)\,(\cos^{2z_2-1}t) = B(z_1,z_2) = \frac{\Gamma(z_1)\,\Gamma(z_2)}{\Gamma(z_1+z_2)},
\]
we compute\vspace*{-3pt}
\begin{align*}
\int_0^\pi d\phi\,(\sin^\spfrac{Q}{Q-1}\phi) &= \sqrt\pi\,\frac{\Gamma(\frac{2Q-1}{2(Q-1)})}{\Gamma(\frac{2Q-1}{2(Q-1)}+\frac12)},\\
\int_0^\pi d\theta \,(\sin^{n-1+\sfrac{Q}{2(Q-1)}}\theta) &= \sqrt\pi\,\frac{\Gamma(\frac n2+\frac{Q}{4(Q-1)})}{\Gamma(\frac n2+\frac{Q}{4(Q-1)}+\frac12)}
\end{align*}
and, when $n\geq 2$,\vspace*{-3pt}
\begin{align*}
\int_\Sigma d\mu(\omega)\,\omega_n^\spfrac{Q}{Q-1} & = \biggl(\prod_{j=1}^{n-2} \int_0^\sfrac\pi2 d\theta_j\,\sin^{2j-1} \theta_j \cos\theta_j\biggr) \\[-8pt]
& \hspace*{3cm} \times \int_0^\sfrac\pi2 d\theta_{n-1} \,\sin^{2n-3} \theta_{n-1} \cos^{1+\frac{Q}{Q-1}} \theta_{n-1} \\[3pt]
& = \biggl(\prod_{j=1}^{n-2} \frac12 \frac{\Gamma(j)}{\Gamma(j+1)}\biggr) \frac12 \frac{\Gamma(n-1)\,\Gamma(1+\frac{Q}{2(Q-1)})}{\Gamma(n+\frac{Q}{2(Q-1)})} \\
& = 2^{-n+1} \frac{\Gamma(1+\frac{Q}{2(Q-1)})}{\Gamma(n+\frac{Q}{2(Q-1)})}.
\end{align*}
For $n=1$ the same formula remains valid, for in this case the integral is trivially equal to one. Thus, we~have shown that\vspace*{-3pt}
\begin{align*}
\frac{I[F_*]}{|F_*|^\sfrac1Q} &= \alpha Q^\sfrac1Q \left(\pi^n \frac{\Gamma(\frac{2Q-1}{2(Q-1)})}{\Gamma(\frac{2Q-1}{2(Q-1)}+\frac12)} \frac{\Gamma(\frac n2+\frac{Q}{4(Q-1)})}{\Gamma(\frac n2+\frac{Q}{4(Q-1)}+\frac12)}
\frac{\Gamma(1+\frac{Q}{2(Q-1)})}{\Gamma(n+\frac{Q}{2(Q-1)})}\right)^\psfrac{Q-1}{Q}\\
&= \alpha \mathcal C_n'.
\end{align*}
This completes the proof of the lemma.
\end{proof}

At this point the proof of Theorem \ref{isoper} is complete. We~now bring the lower bound in the theorem in a more explicit form in dimensions $n=1$ and $n=2$.

\begin{corollary}\label{isopercor}
For $n=1,2$ one has\vspace*{-3pt}
\[
I(\mathbb H)\geq 8 \cdot 3^{-\sfrac 98} \pi^{\sfrac14}
\quad\text{and}\quad
I(\mathbb H_2) \geq \pi^\sfrac43 6^{-\sfrac16} \left(\frac{\Gamma(\frac{13}{5})}{\Gamma(\frac{11}{10})} \frac{\Gamma(\frac{9}{5})}{\Gamma(\frac{13}{10})}\right)^{\sfrac{5}{6}}.
\]
\end{corollary}

Let us discuss the value of the lower bound of $I(\mathbb H)$. Numerically, one has\vspace*{-3pt}
\[
8 \cdot 3^{-\sfrac 98} \pi^{\sfrac14} \approx 3.09468.
\]
This should be compared with Pansu's conjecture \eqref{eq:pansusconj}, which gives the value\vspace*{-3pt}
\begin{equation}
\label{eq:pansun1}
2^{\sfrac52} 3^{-\sfrac34} \pi^{\sfrac12} \approx 4.39854.
\end{equation}
Thus, our value is still quite a bit away from the conjectured sharp value. On the other hand, it~improves over Pansu's original value\vspace*{-3pt}
\[(
8\pi/3)^\sfrac14\approx1.7013.
\]
See \cite{Pa} and also \cite[(7.15)]{CaDaPaTy}. (The discrepancy to the latter formula by a factor of~$4^\sfrac14$ comes from the different normalization of the vector fields used in that reference.)

\begin{proof}
The claimed formula for $n=1$ follows from Theorem \ref{isoper}, the value $\mathcal C_1=(4\pi)^{-1}$ and
\begin{equation}\label{eq:bathtubconst1}
\mathcal C_1' = 2^{-1} 3^\sfrac 98 \pi^\sfrac34.
\end{equation}
To prove the latter formula, we~note that, since $Q=4$,
\[
\mathcal C_1' = 4^\sfrac14 \left(\pi \frac{\Gamma(\frac 76) \,\Gamma(\frac 56)}{\Gamma(\frac{5}{3})\,\Gamma(\frac 43)}\right)^\sfrac34.
\]
By the reflection formula for the gamma function we deduce
\[
\frac{\Gamma(\frac 76) \,\Gamma(\frac 56)}{\Gamma(\frac{5}{3})\,\Gamma(\frac 43)}
= \frac{\frac16\,\Gamma(\frac16)\,\Gamma(\frac56)}{\frac23 \Gamma(\frac23)\,\frac13\,\Gamma(\frac13)} = \frac{3}{4} \,\frac{\Gamma(\frac16)\,\Gamma(\frac56)}{\Gamma(\frac13)\,\Gamma(\frac23)} = \frac34 \,\frac{\frac{\pi}{\sin\frac\pi 6}}{\frac{\pi}{\sin\frac\pi3}} = \frac34\,\frac{\sin\frac\pi6}{\sin\frac\pi3} = \frac34\,\frac{\frac{\sqrt 3}2}{\frac12} = \frac{3^\sfrac32}{4}.
\]
This proves the claimed formula \eqref{eq:bathtubconst1}.

The formula for $n=2$ follows directly from Theorem \ref{isoper}, using the definition of $\mathcal C_2$ and $\mathcal C_2'$.
\end{proof}

\begin{remark}
If one uses the lower bounds on $I(\mathbb H_n)$ from Corollary \ref{isopercor} for $n=1,2$ and inserts them into Proposition \ref{symmetrization}, one obtains a lower bound on $C^\FK(\mathbb H_n)$ that is better than the bound given by Lemma \ref{fksob} and the Jerison--Lee value \eqref{eq:jl}. Explicitly, we~obtain $C^\FK(\mathbb H) \geq 8.78829 \approx 2\pi \times 1.3987$ and $C^\FK(\mathbb H_2) \geq 17.9011 \approx 2^{7/3} \pi \times 1.13064$. This leads to the bounds
\[
\gamma(\mathbb H) \leq 1.65737
\quad\text{and}\quad
\gamma(\mathbb H_2) \leq 1.26183,
\]
which are both unsatisfactory. We~do see, however, that these values are better than the bounds from \eqref{eq:gammatilden}, which are stated after Proposition \ref{gammatildemono}, viz.
\[
\gamma(\mathbb H) \leq 3.2423
\quad\text{and}\quad
\gamma(\mathbb H_2) \leq 1.8238.
\]
When $n=3$, the lower bound from Theorem \ref{isoper}, when inserted into Proposition~\ref{symmetrization}, does not improve over the bound on $C^\FK(\mathbb H_3)$ given by Lemma \ref{fksob} and~\eqref{eq:jl}.
\end{remark}

Despite the negative results in the previous remark, we~will see that Theorem \ref{isoper} will be useful when dealing with $\mathbb H_n\times\R^k$ with $n=1$ and $n=2$.

\subsection{A bound on the isoperimetric constant in $\mathbb H_n\times\R^k$}\label{sec:maincompn12}
We recall that $I(\mathbb H_n\times\nobreak\R^k)$ denotes the isoperimetric constant on $\mathbb H_n\times\R^k$. Here we prove a lower bound on this constant in terms of $I(\mathbb H_n)$ (and the known isoperimetric constants in Euclidean space).

\begin{theorem}\label{isoperk}
For any $n,k\in\N$,
\[
I(\mathbb H_n\times\R^k) \geq I(\R^{2n+2+k}) \Bigl(\frac{I(\mathbb H_n)}{I(\R^{2n+2})}\Bigr)^\pspfrac{2n+2}{2n+2+k}.
\]
\end{theorem}

The proof of the theorem relies on the following lemma that concerns functions $v$ of bounded variations; see, e.g., \cite[Chap.\,2 \& 7]{Leo}. By~definition, the distributional derivative of $v$ is a measure. By~Lebesgue's decomposition theorem we can write this derivative as the sum of a singular measure $Dv^{(\mathrm{s})}$ and a measure that is absolutely continuous with respect to Lebesgue measure and whose density we denote by $v'$.

\begin{lemma}\label{isocartesian}
Let $N\geq 2$ and let $v\in L^\sfrac{N-2}{N-1}(\R)$ be a function of bounded variation. Then, for any $\beta>0$,
\begin{multline*}
\int_\R \sqrt{ (\beta v^\sfrac{N-2}{N-1})^2 + (v')^2}\,dt + |Dv^{(\mathrm{s})}|(\R)\\[-8pt]
\geq I(\R^N) \Bigl(\frac{\beta}{I(\R^{N-1})}\Bigr)^\sfrac{N-1}{N} \biggl(\int_\R |v| \,dt\biggr)^\sfrac{N-1}{N}.
\end{multline*}
\end{lemma}

If $N=2$, we~understand $v^\sfrac{N-2}{N-1}$ as the characteristic function of $\{|v|\neq 0\}$, and we understand the assumption $v\in L^\sfrac{N-2}{N-1}(\R)$ as requiring that this set has finite measure.

If in the first formula in \eqref{eq:isofkexplicit} we replace $|\Sph^{k-1}|$ by $2\pi^{\sfrac k2} \Gamma(\frac k2)^{-1}$, then we see that $I(\R^k)$ can be defined for any (not necessarily integer) real number $k>0$. With this definition, Lemma \ref{isocartesian} remains valid for (not necessarily integer) $N\geq 2$. However, for us only the case where $N$ is an integer case is relevant and we only provide the \emph{proof} in this case. We~note that by scaling, if~the inequality holds for one value of $\beta>0$, then it holds for any such value. Thus, we~may assume that $\beta = I(\R^{N-1})$. In~this case, the lemma follows from the isoperimetric inequality in $\R^N$ applied to the set $\{ (x',x_N)\in\R^{N-1}\times\R:\ \omega_{N-1} |x'|^{N-1}<v(x_N)\}$, where $\omega_{N-1}$ is the volume of the unit ball in $\R^{N-1}$.

\begin{proof}[Proof of Theorem \ref{isoperk}]
Let $E\subset\mathbb H_n\times\R^k$ have finite measure and finite perimeter and set
\begin{align*}
v(t) & := |\{ \zeta\in\mathbb H_n\times\R^{k-1} :\ (\zeta,t)\in E\}|, \\
p(t) & := \per_{\mathbb H_n\times\R^{k-1}} \{ \zeta\in\mathbb H_n\times\R^{k-1} :\ (\zeta,t)\in E\}.
\end{align*}
It follows from Fubini's theorem that $v$ is an integrable function with
\[
\int_\R v(t)\,dt = |E|.
\]
Moreover, one can show that $v$ is of bounded variation and that $p$ is integrable and~that
\[
\per_{\mathbb H_n\times\R^k} E \geq \int_\R \sqrt{ p^2 + (v')^2}\,dt + |Dv^{(\mathrm{s})}|(\R).
\]
(This requires some work. Similar facts appear in the review paper \cite{Ta} on the isoperimetric inequality on $\R^k$.) We bound
\[
p(t) \geq I(\mathbb H_n\times\R^{k-1})\,v(t)^\pspfrac{2n+k}{2n+k+1}
\quad\text{for all}\ t\in\R
\]
and apply the lemma with $\beta = I(\mathbb H_n\times\R^{k-1})$ and $N=2n+2+k$ to deduce that
\[
\per_{\mathbb H_n\times\R^k} E \geq
I(\R^{2n+2+k}) \Bigl(\frac{I(\mathbb H_n\times\R^{k-1})}{I(\R^{2n+1+k})}\Bigr)^\pspfrac{2n+1+k}{2n+2+k} |E|^\pspfrac{2n+1+k}{2n+2+k}.
\]
Thus,
\[
I(\mathbb H_n\times\R^k) \geq
I(\R^{2n+2+k}) \Bigl(\frac{I(\mathbb H_n\times\R^{k-1})}{I(\R^{2n+1+k})} \Bigr)^\pspfrac{2n+1+k}{2n+2+k}.
\]
Dropping $n$ from the notation and abbreviating
\[
\iota_k:= \left(I(\mathbb H_n\times\R^k)/I(\R^{2n+2+k})\right)^{2n+2+k},
\]
we can write this bound as $\iota_k\geq\iota_{k-1}$. Thus, $\iota_k\geq\iota_0$, which is the assertion of the theorem.
\end{proof}

As an immediate consequence of Theorem \ref{isoperk} and Proposition \ref{symmetrization} (in the form \eqref{eq:symmetrization}) we obtain the following lower bound on the Faber--Krahn constant.

\begin{corollary}\label{isoperkcor}
For any $n,k\in\N$,
\begin{equation}
\label{eq:isoperkcoralt}
C^\FK(\mathbb H_n\times\R^k) \geq C^\FK(\R^{2n+k+2}) \Bigl(\frac{I(\mathbb H_n)}{I(\R^{2n+2})} \Bigr)^\spfrac{2(2n+2)}{2n+2+k}.
\end{equation}
\end{corollary}

One can bring the bound \eqref{eq:isoperkcoralt} into a somewhat more explicit form by using the explicit expressions for $C^\FK(\R^{2n+k+2})$ and $I(\R^{2n+2})$ from \eqref{eq:isofkexplicit}.

With the help of Corollary \ref{isopercor} we are ready to prove that $\gamma(\mathbb H_n\times\R^k)<1$ if $n=1$ and $k\geq 2$ or if $n=2$ and $k\geq 1$.

\skpt
\begin{proof}[Proof of Theorem \ref{maincomp}. Cases $n\!=\!1$ and $n\!=\!2$]
We recall the expression for $\mathcal W(\mathbb H_n\times\R^k)$ from \eqref{eq:wnk} and write it as
\[
\mathcal W(\mathbb H_n\times\R^k) = \mathcal W(\R^{2n+2+k}) \mathcal W(\mathbb H_n) (4\pi)^{n+1} \Gamma(n+2).
\]
Using Corollary \ref{isoperkcor} we find that
\begin{equation}
\label{eq:gammankbound}
\gamma(\mathbb H_n\times\R^k) \leq \gamma(\R^{2n+2+k})\ \Bigl(\frac{I(\R^{2n+2})}{I(\mathbb H_n)}\Bigr)^{2n+2} \mathcal W(\mathbb H_n)^{-1} (4\pi)^{-n-1} \Gamma(n+2)^{-1}.
\end{equation}
Note that the upper bound depends on $k$ only through $\gamma(\R^{2n+2+k})$. We~recall from \cite[Th.\,5.1]{HPS} that the sequence $d\mto \gamma(\R^d)$ is decreasing. Therefore, if~we can show that the upper bound is $<1$ for $(n,k)=(1,2)$ and $(n,k)=(2,1)$, then the assertions of the theorem for $n=1$ and $n=2$ will follow.

\subsubsection*{Case $n=1$} Recalling from \eqref{eq:W1} that $\mathcal W(\mathbb H) = \sfrac1{128}$, we~get
\begin{align}\label{eq:gammankboundn1}
\gamma(\mathbb H\times\R^k)
& \leq 4^4 \Bigl(\frac{\Gamma(\frac{6+k}{2})}{\Gamma(3)}\Bigr)^2
\left(I(\mathbb H)\right)^{-4} \left(\sfrac{1}{j_{\psfrac{2+k}{2},1}}\right)^{4+k}
\left(\mathcal W(\mathbb H)\right)^{-1} 2^k \notag \\
& = 2^{13+k} \Gamma(\tfrac{6+k}2)^2 \left(I(\mathbb H)\right)^{-4} \left(\sfrac{1}{j_{\psfrac{2+k}{2},1}}\right)^{4+k}.
\end{align}
Specializing further to $k=2$, this bound becomes
\begin{align*}
\gamma(\mathbb H\times\R^2)
& \leq 2^{17} 3^2 \left(I(\mathbb H)\right)^{-4} \left(\sfrac{1}{j_{2,1}}\right)^6.
\end{align*}
We have $j_{2,1}\approx 5.13562$ \cite[Table\,9.5]{AbSt} and, using the lower bound on $I(\mathbb H)$ from Corollary \ref{isopercor}, we~obtain
\[
\gamma(\mathbb H\times\R^2) \leq 0.701019.
\]
This is $<1$, as desired.

\subsubsection*{Case $n=2$} Recalling from \eqref{eq:W2} that $\mathcal W(\mathbb H_2) = \sfrac{1}{48^2\,\pi}$, we~get
\begin{align*}
\gamma(\mathbb H_2\times\R^k)
& \leq
6^6 \Bigl(\frac{\Gamma(\frac{8+k}{2})}{\Gamma(4)}\Bigr)^2
\left(I(\mathbb H_2)\right)^{-6} \left(\sfrac{1}{j_{\psfrac{4+k}{2},1}}\right)^{6+k}
\left(\mathcal W(\mathbb H_2)\right)^{-1} 2^k \\
& = 2^{12+k} 3^6 \pi \Gamma(\tfrac{8+k}{2})^2 \left(I(\mathbb H_2)\right)^{-6} \left(\sfrac{1}{j_{\psfrac{4+k}{2},1}}\right)^{6+k}.
\end{align*}
Specializing further to $k=1$, this bound becomes
\begin{align*}
\gamma(\mathbb H_2\times\R)
& \leq 2^{5} 3^8 5^2 7^2 \pi^2 \left(I(\mathbb H_2)\right)^{-6} \left(\sfrac{1}{j_{\sfrac{5}{2},1}}\right)^{7}.
\end{align*}
We have $j_{\sfrac52,1}\approx 5.76346$ \cite[Table\,10.6]{AbSt} and, using the lower bound on $I(\mathbb H_2)$ from Corollary \ref{isopercor}, we~obtain
\[
\gamma(\mathbb H_2\times\R) \leq 0.823715.
\]
This is $<1$, as desired.
\end{proof}

As another application of Corollary \ref{isoperkcor} we now show that Pansu's conjecture implies Pleijel's theorem.

\begin{proof}[Proof of Proposition \ref{pansupleijel}]
In view of Theorem \ref{maincomp} and Corollary \ref{pansupleijelk0} it suffices to prove that Pansu's conjecture in $\mathbb H$ implies $\gamma(\mathbb H\times\R)<1$. Arguing as in the previous proof, we~want to show that the right side of \eqref{eq:gammankboundn1} is $<1$ when $k=1$ and when $I(\mathbb H)$ is replaced by the conjectured value on the right side of \eqref{eq:pansusconj}, which is $2^{\sfrac 52} 3^{-\sfrac 34}\pi^{\sfrac 12}$; see \eqref{eq:pansun1}. Using $j_{\sfrac32,1}\approx 4.49341$ \cite[Table\,10.6]{AbSt}, we~obtain
\[
2^{14} \,\Gamma(\tfrac 72)^2 \bigl(2^{\sfrac 52} 3^{-\sfrac 34}\pi^{\sfrac 12}\bigr)^{-4} j_{\sfrac32,1}^{-5} \approx 2.639 \cdot 10^{-1} <1,
\]
as claimed.
\end{proof}

\appendix

\part{Appendix}

\section{Local linear independence}\label{AppA2}

In our description of the nilpotent approximation in Section~\ref{s1} and in the proof of our main theorem in Section~\ref{sec:mainsubriemproof} we made use of results by Rothschild \cite{Rot} that were established under the assumption that the vectors $X_1(x),\ldots,X_p(x)$ are linearly independent at one (and hence, by~equiregularity) any point $x\in M$. In~this appendix we show that locally around any given point we may always reduce ourselves to the case where this is satisfied. We~emphasize that this is well-known in sub-Riemannian geometry and we briefly sketch the proof in that language in Subsection \ref{sec:subriemannianortho}. We~think, however, that it is beneficial to also give an elementary proof ``by hand''. Finally, in~Subsection \ref{ssA.4} we discuss topological obstructions to a global version of this statement.

To be precise, in this appendix we work under Assumptions \ref{CHOr} and \ref{equireg} and we recall that $n_j=\dim\mathcal D_j(x)$. We~shall prove the following.

\begin{lemma}\label{locindep}
For any $x_0\in M$ there is an open neighborhood $W\subset M$ of $x_0$ and smooth vector fields $\tilde X_1,\ldots,\tilde X_{n_1}$ defined in $W$ such that
\begin{equation}
\label{eq:loclinindepapp}
\Span\{ \tilde X_1(x),\ldots \tilde X_{n_1}(x) \} = \mathcal D_1(x)
\quad\text{for all}\ x\in W
\end{equation}
and
\[
-\Delta_\bX^{M,\mu} f = \sum_{j=1}^{n_1} \tilde X_j^\star \tilde X_j f
\quad\text{for all} \ f\in C^2(W).
\]
The vector fields $\tilde X_1,\ldots,\tilde X_{n_1}$ still satisfy Assumptions \ref{CHOr} (with the same $r$) and \ref{equireg} (with the dimensions $n_j$).
\end{lemma}

\subsection{A first example}
To start, we~consider the example where
\[
-\Delta_\Xb = X_1^\star X_1 + X_2^\star X_2 + X_3^\star X_3.
\]
We suppose that Assumptions \ref{CHOr} with $r=2$ and Assumption \ref{equireg} are satisfied, but we assume that at each point $x\in M$ the space generated by $X_1(x), X_2(x), X_3(x)$ is of dimension $2$.

Let us derive Lemma \ref{locindep} in this particular case, that is, show that locally, we~can represent $-\Delta_\Xb$ in the form
\[
- \Delta_\Xb = \tilde X_1^\star \tilde X_1 + \tilde X_2^\star \tilde X_2,
\]
for suitable vector fields $\tilde X_1,\tilde X_2$, such that for all $x$ in a neighborhood of a given point the vectors $\tilde X_1(x)$ and $\tilde X_2(x)$ are linearly independent and belong to $\mathcal D_1(x)$.

We proceed as for the proof of the Morse lemma. Let us assume that $X_1$ and $X_2$ are linearly independent in some open set $W$ in $M$. Hence we can write for $x\in W$
\[
X_3(x) = a_1(x) X_1(x) + a_2(x) X_2(x)
\]
with two smooth functions $a_1$ and $a_2$ on $W$. The two desired vector fields are given by
\begin{equation}
\begin{array}{ll}
\tilde X_1 &= \sqrt{1+a_1^2} \,X_1 + \frac{a_1 a_2}{\sqrt{1+ a_1^2}} \,X_2, \\
\tilde X_2 &= \sqrt{\frac{1+ a_1^2 + a_2^2}{1+a_1^2}} \,X_2.
\end{array}
\end{equation}
Clearly $\tilde X_1$ and $\tilde X_2$ are linearly independent. Note also that
\[
[\tilde X_1,\tilde X_2] = \sqrt{1+a_1^2 +a_2^2}\,[X_1,X_2] + b_1 X_1 + b_2 X_2
\]
with two smooth functions $b_1$ and $b_2$ on $W$. This proves the assertions of Lemma \ref{locindep} in this example.

In passing we mention that in this case we can take as an adapted flag (see Section~\ref{s1}) the vector fields $\tilde X_1, \tilde X_2, [\tilde X_1,\tilde X_2] $, so that the group $G_x$ is the Heisenberg group.

\subsection{General argument}
The argument in the general case is as follows: we consider a sub-Laplacian $\Delta_\Xb^{M,\mu} =-\sum_{i=1}^p X_i^\star X_i$ and assume that $p>n_1$ (for otherwise there is nothing to prove). For $f\in C^\infty(M)$ we consider
\[
Q[f](x) := \sum_{j=1}^p (X_j f (x))^2.
\]

We choose $n_1$ vector fields, which we can assume to be $X_1,\ldots,X_{n_1}$ (up to relabeling), that are linearly independent in a neighborhood of some point $x_0\in M$. In~the sequel $x$ denotes the variable in this neighborhood. For $j=n_1+1,\ldots,p$, we~find coefficients $a_{ij}$ such that
\[
X_j=\sum_{\ell=1}^{n_1} a_{j\ell}(x)X_\ell.
\]
Expanding, we~obtain for $f$ with compact support in this neighborhood
\[
Q[f](x)=\sum_{\ell, \ell'=1}^{n_1} b_{\ell \ell'}(x)X_\ell f(x) X_{\ell'}f(x)
\]
with some coefficients $b_{\ell \ell'}$ computed in terms of the $a_{j\ell}$. The matrix
\[
B_x :=(b_{\ell \ell'}(x))_{1\leq \ell,\ell'\leq n_1}
\]
is symmetric and positive definite.

In this case, we~can have a normal form (see the proof of the Morse lemma):
\[
B_x = T_x^t\,T_x,
\]
where $T_x = (t_{i\ell}(x))$ is triangular, invertible, depending smoothly on $x$ in the construction. We~set for $i=1,\ldots,n_1$,
\[
\tilde{X}_i(x) := \sum_{\ell=i}^{n_1} t_{i\ell}(x)X_\ell(x)
\]
and get
\[
Q[f](x) =\sum_{i=1}^{n_1} (\tilde{X}_i f(x))^2.
\]
Integrating this identity with respect to $\mu$ we conclude that, locally,
\[
\Delta = - \sum_{i=1}^{n_1} \tilde{X}_i^\star \tilde{X}_i,
\]
as claimed. Clearly, at each point $x$ in the relevant neighborhood the span of the vectors $\tilde X_1(x),\ldots,\tilde X_{n_1}(x)$ is equal to $\mathcal D_1(x)$. Since the $\mathcal D_j$ for $j\geq 1$ depend on the vector fields $X_1,\ldots,X_p$ only through their span $\mathcal D_1$, we~see that the $\tilde X_1,\ldots,\tilde X_{n_1}$ satisfy Assumptions \ref{CHOr} (with the same $r$) and \ref{equireg} (with the dimensions $n_j$). This concludes the proof of Lemma \ref{locindep}.

\subsection{Sub-Riemannian geometric construction}\label{sec:subriemannianortho}
We provide\footnote{Discussions with Y.\,Colin de Verdière, L.\,Hillairet and C.\,Letrouit.} an alternative approach to the argument of the previous subsection. We~refer for example to the books \cite{AgBaBo,Je} or to the appendices of \cite{CHT,CHT1} for further background.

We introduce the following metric on the distribution: for $x\in M$ and $v\in\mathcal D_1(x)$ we set
\[
g_x(v) := \inf\Bigl\{ \textstyle\sum_{j=1}^p u_j^2 :\ \sum_{j=1}^p u_j X_j(x)=v\Bigr\}.
\]
One can show that this defines a positive definite quadratic form, and thus a metric, on $\mathcal D_1(x)$, which depends smoothly on $x$. We~omit a proof of these assertions, which can be verified using for instance some ideas used in the proof of \eqref{e:equality} below.

Let $(\tilde{X}_1(x),\ldots,\tilde{X}_{n_1}(x))$ be a smooth local orthonormal frame for this metric. We~claim that
\begin{equation}\label{e:equality}
\sum_{j=1}^p (X_j f)^2=\sum_{i=1}^{n_1} (\tilde{X}_if)^2.
\end{equation}
Once we have verified this, we~obtain another proof of Lemma \ref{locindep}.

Identity \eqref{e:equality} is well-known and can be proved in several ways, but for the sake of completeness we outline one possible proof. Introducing
\[
E_x(v)= \Bigl\{\textstyle u\in \mathbb R^p\,:\ \sum_{j=1}^{p} u_j X_j(x) =v\Bigr\}
\quad\text{for all}\ v\in \mathcal D_1(x),
\]
we notice that
\[
g_x(v) = \inf\Bigl\{ \textstyle\sum_{j=1}^p u_j^2 : u\in E_x(v)\Bigr\}.
\]
For every $v\in D_1(x)$ this infimum is attained at a unique $u_x(v)\in \mathbb R^p$ and this $u_x(v)$ is the orthogonal projection of $0$ onto $E_x(v)$. In~particular, $u_x(v)$ is orthogonal to $E_x(0)$. Let
\[
u_i := (u_{ij})_{j=1}^p := u_x(\tilde X_i(x))
\quad\text{for}\ i=1,\ldots, n.
\]
(The $u_{ij}$ depend on $x$, but since the proof of \eqref{e:equality} is pointwise we sometimes drop $x$ from the notation.) Thus, we~have
\[
\tilde{X}_i=\sum_{j=1}^p u_{ij} X_j
\quad\text{and}\quad\sum_{j=1}^p u_{ij}^2=1.
\]
As observed above, we~have $u_i$ is orthogonal to $E_x(0)$, so if $i\neq i'$, then $u_i+ u_{i'}\in E_x(\tilde{X}_i+\tilde X_{i'})$ is orthogonal to $E_x(0)$ and $u_i + u_{i'} = u_x(\tilde{X}_i+\tilde X_{i'})$. Combining this with the orthogonality of the $\tilde{X}_i$'s, we~get for $i\neq i'$
\[
2= g(\tilde X_i+\tilde X_{i'})= \sum_j (u_{ij} +u_{i'j})^2 = 2 +2 \sum_{j=1}^p u_{ij} u_{i' j},
\]
so
\[
\sum_{j=1}^p u_{ij} u_{i' j}=0 \quad\text{for } i\neq i'.
\]
To summarize, we~have shown that the $(n_1\times p)$-matrix $U:=(u_{ij})$ satisfies
\begin{equation}
\label{eq:subriemproof1}
U U^\mathrm{T} = {\Id_{\R^{n_1}}}.
\end{equation}

We observe that
\begin{equation}
\label{eq:subriemproof3}
\ker U = E_x(0).
\end{equation}
Indeed, by~the orthogonality of the $u_i$'s to $E_x(0)$ we have $E_x(0)\subset \ker U$. This inclusion is an equality since, as a consequence of \eqref{eq:subriemproof1}, $U$ is surjective and therefore $\dim \ker U = p - n_1 = \dim E_x(0)$.

Since the $\tilde X_i$ are a basis, there are uniquely defined coefficients $c_{ij}$ such that
\[
X_j =\sum_{i=1}^{n_1} c_{ij} \tilde X_i.
\]
We note that the $(n_1\times p)$-matrix $C :=(c_{ij})$ satisfies
\begin{equation}
\label{eq:subriemproof4}
\ker C = E_x(0).
\end{equation}
Indeed, $b\in E_x(0)$ if and only if $\sum_i \sum_j b_j c_{ij} \tilde X_i = 0$ if and only $C b =0$. Similarly, it~follows directly from the definitions of $U$ and $C$ that
\begin{equation}
\label{eq:subriemproof2}
U C^\mathrm{T} = {\Id_{\R^{n_1}}}.
\end{equation}

It follows from \eqref{eq:subriemproof1} and \eqref{eq:subriemproof2} that $U(U^\mathrm{T} - C^\mathrm{T}) = 0$, and it follows from \eqref{eq:subriemproof3} and \eqref{eq:subriemproof4} that $\ran(U^\mathrm{T} - C^\mathrm{T}) = (\ker U)^\mathrm{T}$. Combining these two facts, we~conclude that $U^\mathrm{T} - C^\mathrm{T}=0$, that is,
\[
C = U.
\]

Thus, we~have shown that
\[
X_j = \sum_{i=1}^{n_1} u_{ij} \tilde X_i
\]
and, inserting this decomposition into the left side of \eqref{e:equality} and using the properties of the coefficients $u_{ij}$, we~arrive at the identity \eqref{e:equality}. This completes the alternative proof of Lemma \ref{locindep}.

\subsection{Topological constraints}\label{ssA.4}
The existence of $n_1$ global vector fields in $\mathcal D_1$ giving for each
$x\in M$ a basis of $\mathcal D_1(x)$ is only possible under strong topological conditions on $M$ that involve its orientability, its Euler characteristic, the Euler class of $\mathcal D_1$ and other invariants.\footnote{Discussion with V.\,Colin.}

Note first that the existence of a global non-zero vector field implies that the Euler characteristic of $M$ is zero by the Poincaré--Hopf theorem.
When $M$ has dimension~3 and is orientable, this is not an obstruction since this Euler characteristic is zero. A second obstruction related to the Euler class of the subbundle $\mathcal D_1$ of $TM$ also disappears in the case $M=\mathbb S^3$. If $n_1=2$ (and, as before, $M$ has dimension 3 and is orientable), one can instead of a basis of $\mathcal D_1(x)$ consider the unit normal to the plane $\mathcal D_1(x)$ in $T_x M$. By~orientability of $M$, this normal can be globally defined and we get a map $M \to \Sph^2$. In~particular, when $M=\mathbb S^3$ there is an associated Hopf invariant, which belongs to $\mathbb Z$, and this invariant should be zero. This is an obstruction that is not verified in general. In~particular, in the case of the Hopf fibration, this invariant is one and if $\mathcal D_1$ is perpendicular to the fibers, one cannot find a continuous basis over~$\mathbb S^2$.

\section{Weyl law and heat kernel}\label{AppC}
This section is based on notes kindly transmitted to us by C.\,Letrouit. We~make a connection with the framework of \cite{CHT1}.

\subsection{Nilpotentization of measure}
We follow \cite[\S A.5.6]{CHT1}. Let $\mu$ be a smooth measure on $M$ and let $x \in M$. Recall that privileged coordinates are defined by $\theta_x$, which maps a neighborhood of $x$ in $M$ to a neighborhood of $0$ in $\R^n$. We~also recall the definition of the family of dilations $\delta_t$ on $\R^n$ from \eqref{eq:dilation}. Setting $\delta_x^{(\epsilon)} := \theta_x^{-1} \circ \delta_\varepsilon$, for $\varepsilon>0$ we can define the measures
\[
\mu_x^{(\varepsilon)} := \varepsilon^{-Q}(\delta_x^{(\varepsilon)})^*\mu
\]
on a neighborhood of the origin in $\R^n$. The nilpotentization $\widehat\mu_x$ of $\mu$ at $x$, which we defined in \eqref{eq:nilpmeas} ``by hand'', satisfies
\[
\widehat{\mu}_x=\lim_{\varepsilon\to 0_+}\mu_x^{(\epsilon)}
\]
with convergence in the vague topology.

\subsection{Nilpotentization of spectral function and heat kernel}
We follow \cite[\S A.8.3]{CHT1}. We~recall that the nilpotentized sub-Laplacian at $x$ is defined as an operator on functions on $G_x\simeq \R^n$ by
\[
\widehat{\Delta}_x= \sum_{i=1}^{n_1} (\widehat{X}_{i,x})^2.
\]
It is self-adjoint in $L^2(G_x,\widehat\mu_x)$ with the usual domain and nonnegative.

In this appendix we denote by
\[
\widehat e_x(\lambda,u,v) := \1(-\widehat\Delta_x <\lambda)(u,v)
\]
the spectral function of $-\widehat\Delta_x$, that is, the integral kernel of the spectral projection $\1(-\widehat\Delta_x<\lambda)$ for $\lambda>0$.

The operator $\partial_t- \widehat{\Delta}_x$ is hypoelliptic and therefore $\exp (t \widehat \Delta_x)$ is an integral operator. Its integral kernel is denoted by
\[
\widehat k_x(t,u,v) := (\exp(t\widehat\Delta_x))(u,v),
\]
that is,
\[
(\exp(t \widehat \Delta_x) f)(u)= \int_{G_x} \widehat k_x(t,u,v) f(v)\,d \widehat{\mu}_x(v)
\quad \forall u\in G_x.
\]
According to the functional calculus we have
\begin{equation}
\label{eq:heatkernelspecfcn}
\widehat k_x(t,u,v) = \int_0^\infty e^{-t\lambda} \widehat e_x(\lambda,u,v)\,d\lambda
\quad\text{for all}\ t>0.
\end{equation}

\subsection{Effect of changes of coordinates}
We follow \cite[2nd part of App.\,A.8.2]{CHT1}. The above nilpotentizations of the spectral function and of the heat kernel depend on the choice of privileged coordinates (we have made a specific choice in the main text) and on the measure $\widehat\mu_x$. We~claim that diagonal values
\[
\widehat e_x(\lambda,u,u)
\quad\text{and}\quad
\widehat k_x(t,u,u)
\]
do not depend on this choice.

Indeed, this follows from \cite[third relation in (85)]{CHT}, which says that changing variables both in the operator and in the measure has no effect on the heat kernel. More precisely, if~we denote the integral kernel of $-\Delta^{M,\mu}_\bX$ by $k_{\Delta,\mu}$ and if $\varphi$ is a diffeomorphism of $M$, then
\begin{equation}
\label{eq:diffeoheat}
k_{\varphi^* \Delta \varphi_*, \varphi^*\mu}(t,u,v)=k_{\Delta,\mu}(t,\varphi(u),\varphi(v)).
\end{equation}
We apply this formula with $M$ replaced by $G_x$, with $-\Delta^{M,\mu}_\bX$ replaced by $-\widehat\Delta_x$ and with $\mu$ replaced by $\widehat\mu_x$. A similar argument applies to the spectral function.

Choosing $\varphi$ in \eqref{eq:diffeoheat} as the diffeomorphism on $G_x$ given either by dilation or by translation by a group element, we~see that
\[
\widehat k_x(t,u,u) = c^\heat_x \,t^{-\sfrac Q2}
\quad\text{for all}\ u\in G_x,\\ t>0,
\]
where
\[
c^\heat_x = \widehat k_x(1,0,0).
\]
These are the analogues of \eqref{eq:specfcndiag1} and \eqref{eq:specfcndiag2}. It~follows from \eqref{eq:heatkernelspecfcn} that
\begin{equation}
\label{eq:heatkernelspecfcnconst}
c_x^\heat = \Gamma(\tfrac Q2+1) \,c^\Weyl_x.
\end{equation}

\subsection{Heat kernel asymptotics and Weyl law}
The following appears in \cite[Th.\,I \& 4.1]{CHT1}:

\begin{theorem}\label{weylapp}
In the equiregular case, for every $f\in C^\infty(M)$, we~have
\[
\Tr f e^{t\Delta^{M,\mu}_\bX} =
\int_M f(x)k_{\Delta,\mu}(t,x,x)\,d\mu(x) = t^{-\sfrac Q2}\,F(t)
\]
for some $F\in C^\infty(\R)$ with
\[
F(0)=\int_M f(x) \,c^\heat_x \,d\mu(x).
\]
Moreover, the eigenvalue counting function satisfies
\begin{equation}\label{e:specasymp}
N(\lambda,-\Delta^{M,\mu}_\bX)\sim \int_M c^\Weyl_x \,d\mu(x) \ \lambda^{\sfrac Q2}
\quad\text{as}\ \lambda\to\infty,
\end{equation}
where $c^\Weyl_x$ and $c^\heat_x$ are related by \eqref{eq:heatkernelspecfcnconst}.
\end{theorem}

Behind the proof of this theorem is the fact (see \cite[Th.\,C.1 and Eq.\,(96)]{CHT1}) that $t^{\sfrac Q2} k_{\Delta,\mu}(t,x,x)$ converges to $\widehat{k}_x(1,0,0)=c^\heat_x$ as $t\to 0^+$, uniformly with respect to $x$. This uses the equiregularity assumption. The spectral asymptotics \eqref{e:specasymp} follows from the heat kernel asymptotics by a Tauberian theorem. This gives \eqref{e:specasymp} with $\Gamma(\frac Q2 +1)^{-1} c^\heat_x$ instead of $c_x^\Weyl$. These coefficients coincide in view of \eqref{eq:heatkernelspecfcnconst}.
